\documentclass[11pt,a4paper]{article}
\usepackage[hmargin=27mm,vmargin=25mm]{geometry}
\usepackage[T1]{fontenc}
\usepackage{amsmath,amssymb,amsthm,mathtools,microtype}
\usepackage{mathrsfs}
\usepackage{array}
\usepackage[dvipsnames]{xcolor}
\usepackage[colorlinks=true,linkcolor=MidnightBlue,citecolor=MidnightBlue,urlcolor=MidnightBlue,breaklinks=true]{hyperref}
\usepackage{bookmark}
\hypersetup{pdftitle={Pointwise Correlations of Quadratic Form Values and Splitting Semigroups},pdfauthor={Roy van Rijn}}
\newtheorem{theorem}{Theorem}[section]
\newtheorem{proposition}[theorem]{Proposition}
\newtheorem{lemma}[theorem]{Lemma}
\newtheorem{corollary}[theorem]{Corollary}
\theoremstyle{definition}
\newtheorem{definition}[theorem]{Definition}
\theoremstyle{remark}
\newtheorem{remark}[theorem]{Remark}
\newcommand{\Z}{\mathbb Z}
\newcommand{\Q}{\mathbb Q}
\newcommand{\C}{\mathbb C}
\newcommand{\N}{\mathbb N}
\newcommand{\E}{\mathbb E}
\newcommand{\Prob}{\mathbb P}
\newcommand{\1}{\mathbf 1}
\DeclareMathOperator{\supp}{supp}
\DeclareMathOperator{\lcm}{lcm}
\DeclareMathOperator{\rad}{rad}
\DeclareMathOperator{\Res}{Res}
\DeclareMathOperator{\Cl}{Cl}
\DeclareMathOperator{\Gal}{Gal}
\DeclareMathOperator{\Frob}{Frob}
\DeclareMathOperator{\TV}{TV}
\newcommand{\ee}[1]{e^{2\pi i #1}}
\allowdisplaybreaks[2]
\setlength{\emergencystretch}{1.5em}
\title{Pointwise Correlations of Quadratic Form Values\\
and Splitting Semigroups}
\author{Roy van Rijn}
\date{}
\begin{document}
\maketitle
\begin{abstract}
Under the stated analytic inputs, we establish the fixed-shift pair
asymptotics predicted by Smilansky for the ordinary values of
\(x^2+dy^2\), \(d=1,2,3,4,7\).
In particular, the number of \(n\le X\) for which both \(n\) and
\(n+1\) are sums of two squares is asymptotic to \(X/(2\log X)\).
For primitive positive definite binary forms of fixed discriminants
\(\Delta_i=f_i^2D\) sharing one negative fundamental discriminant \(D\),
we count values admitting coprime coordinates. Replacing either form
within its respective genus preserves the leading binary count, with
a relative logarithmic saving through an independent classical route.
Under the recent analytic inputs we give an explicit mixed leading term
with relative \(o(1)\), including different fixed conductors in that
one field, and no class-number divisor.
The central mechanism is a shared comparison between sparse indicators
and arithmetic models retaining every surviving congruence character.
Following Walsh's proposed indicator extension, we supply the models,
local factors and a uniform reflected-target argument, using the recent
zero-free and weighted-dilation theorems of OpenAI.
Applications to fixed subgroup restrictions and prime-coset self-sums
give target-uniform asymptotics with logarithmic savings and exact
eventual congruence criteria. They include prime support \(1\pmod{12}\)
on targets \(2\pmod{12}\) and prime support \(3\pmod4\) on even targets.
All counts concern numerical values without coordinate or ideal multiplicities.
\end{abstract}
\section{Introduction}
How many positive integers \(n\le X\) have the property that both
\(n\) and \(n+1\) are sums of two squares? Write \(b_d(n)\) for
the indicator of the positive values of \(x^2+dy^2\), with
\(x,y\in\Z\), and put \(b_d(n)=0\) for \(n\le0\).
Coordinates are unrestricted and may vanish; each numerical value
is counted once. Under the analytic inputs stated below, we obtain
\begin{equation}\label{eq:ck-headline}
 \#\{1\le n\le X:b_1(n)=b_1(n+1)=1\}
                         \sim\frac{X}{2\log X}.
\end{equation}
This is the shift-one Connors--Keating prediction. It belongs to
Smilansky's five-form family of pair-correlation predictions.

\begin{corollary}\label{cor:smilansky-five-pairs}
Under the inputs of Section~\ref{sec:inputs}, for
\(d\in\{1,2,3,4,7\}\) and fixed \(h\ne0\),
\[
 \sum_{1\le n\le X}b_d(n)b_d(n+h)
       =\mathscr T_{d,h}\frac X{\log X}+o_{d,h}(X/\log X),
\]
with precisely the coefficients predicted in
\cite[Section 6.3, equations (6.2)--(6.4)]{smilansky}.
In particular,
\begin{equation}\label{eq:minus-eight-shift-one}
 \sum_{n\le X}b_2(n)b_2(n+1)
                       \sim\frac1{\sqrt2}\frac X{\log X}.
\end{equation}
\end{corollary}

Section~\ref{sec:quadratic-pairs} evaluates the coefficients and exceptional-prime conditions.

\subsection{A second headline: genus governs primitive value counts}

Fix one negative fundamental discriminant \(D\) and two fixed
discriminants \(\Delta_i=f_i^2D\). Let \(Q_i,Q'_i\) be primitive
positive definite integral binary forms of discriminant \(\Delta_i\),
in the same genus for each \(i=1,2\), in the one field \(\Q(\sqrt D)\).
If \(P_Q\) is the set of positive values of
\(Q\) at coprime coordinates and
\(R_{Q_1,Q_2}(N)=\#\{(a,b)\in P_{Q_1}\times P_{Q_2}:a+b=N\}\),
then the classical route gives
\begin{equation}\label{eq:genus-headline}
 R_{Q_1,Q_2}(N)=R_{Q'_1,Q'_2}(N)
                  \left(1+O((\log N)^{-1/5})\right)
\end{equation}
on sufficiently large common locally admissible targets.
The slots may have different genera and conductors. Under the recent
analytic inputs the same count has an explicit leading term of order
\(NJ_S(N)/\log N\), with relative error \(o(1)\). No class-number divisor
occurs: a value can belong to every class in its compatible genus.
The common-field theorem and its principal specialization are in
Section~\ref{sec:forms}.

The counting conventions and strengths are different:
\begin{center}\small
\begin{tabular}{@{}>{\raggedright\arraybackslash}p{.19\linewidth}>{\raggedright\arraybackslash}p{.34\linewidth}>{\raggedright\arraybackslash}p{.41\linewidth}@{}}
\hline
Result family&What is counted&Conclusion\\\hline
CK/Smilansky&Ordinary values; unrestricted coordinates&Fixed-shift asymptotic with the stated little-\(o\) error\\[1mm]
Subgroup/coset sums&Integers with prescribed prime support&Target-uniform asymptotic with a logarithmic saving\\[1mm]
Primitive-form sums&Values admitting coprime coordinates in two fixed orders of one field&Explicit leading term with relative \(o(1)\)\\[1mm]
Genus replacement&Comparison of two primitive-form value counts&Relative logarithmic saving through the classical route\\\hline
\end{tabular}
\end{center}
The first three asymptotic families use the imports in
Section~\ref{sec:inputs}; genus replacement does not.
A \emph{primitive form} \([a,b,c]\) has \(\gcd(a,b,c)=1\).
A \emph{primitive representation} \(Q(x,y)\) has
\(\gcd(x,y)=1\). \emph{Coprime summands} in \(a+b=N\)
means \(\gcd(a,b)=1\), a separate condition imposed only when stated.
For example, \(4=2^2+0^2\) is an ordinary sum of two squares,
but any representation of four has both coordinates even and is
not primitive. Consequently the primitive-form theorem does not
extend the ordinary CK statement to arbitrary forms. Its class-error
saving also does not supply a logarithmic remainder for the full
conductor-prefix leading term.

\subsection{The common mechanism, in the CK example}

The central result is the indicator-to-model comparison of
Theorem~\ref{thm:model-comparison}. Multiplication by suitable primes
preserves an indicator and its arithmetic model. One-variable means
center their difference, and a positive sparse sieve controls exceptional
prime-factor records. A shared graph reduction then leaves two frequency
estimates: one for fixed shifts and one for reflected targets. Their
assembly replaces the indicators in the binary count by the models.
The model correlation can then be evaluated from local arithmetic.

Here is the model for \(b_1\). At outer scale \(X\), put
\(L=\log X\), \(y=\exp(L^{2/5})\),
\(\mathcal P_y=\{p:p\le y\}\), and
\(V_y=\prod_{p\le y}(1-1/p)\). Write every \(n\) uniquely as
\(n=rv\), where \(r=n_{\mathcal P_y}\) contains its full small-prime
powers and \(v\) has no prime factor at most \(y\).
The rough ordinary-value series is
\[
 F_y(s)=\prod_{\substack{p>y\\p\equiv1\ (4)}}(1-p^{-s})^{-1}
         \prod_{\substack{p>y\\p\equiv3\ (4)}}(1-p^{-2s})^{-1}
       =(s-1)^{-1/2}H_y(s).
\]
For a sufficiently deep fixed \(K\), its smooth density expansion is
\[
 M(v)=\sum_{j=0}^K\frac{H_y^{(j)}(1)}{j!}
            \frac{(\log v)^{-1/2-j}}{\Gamma(1/2-j)},\qquad
 \mathcal D(rv)=b_1(r)\bigl(1+\chi_{-4}(v)\bigr)\frac{M(v)}{V_y}.
\]
These formulas apply on the polynomial rough scales used in the proof;
the smooth bounded extension is specified in
Section~\ref{sec:model-preparation} and the quadratic application.
The leading term of \(M\) is \(H_y(1)/\sqrt{\pi\log v}\).
Dividing by \(V_y\), the density of rough integers, gives the
conditional density on the rough variable. Every allowed rough value
has \(\chi_{-4}(v)=1\). The correction \(1+\chi_{-4}(v)\)
retains this bias; average density alone would miss a leading contribution.
For a split prime \(p\le y\), multiplication by \(p\) changes only
\(r\), preserves \(b_1(r)\), and leaves both \(v\) and
\(\mathcal D\) unchanged. Replacing
\(b_1(n)b_1(n+h)\) by \(\mathcal D(n)\mathcal D(n+h)\)
with an error \(o(X/\log X)\) is the analytic task. The remaining
model product gives the coupled local factors and the coefficient
\(1/2\) in \eqref{eq:ck-headline}.

\subsection{Semigroups and the reflected-target obstacle}

The same mechanism reaches a different geometry. Let
\(C_{12}=\{n\ge1:p\mid n\Rightarrow p\equiv1\pmod{12}\}\),
including one. Every sufficiently large \(N\equiv2\pmod{12}\)
is a sum of two elements of \(C_{12}\), and
the ordered value-pair count
\(R_{12}(N)=\#\{(a,b)\in C_{12}^2:a+b=N\}\) satisfies
\begin{equation}\label{eq:opening}
 R_{12}(N)=A_{12}J_{12}(N)\frac{N}{(\log N)^{3/2}}
                  \left(1+O((\log N)^{-\delta})\right),
 \qquad\delta>0.
\end{equation}
The positive factors are explicit in Theorem~\ref{thm:c12}.
For \(C^-_4=\{n\ge1:p\mid n\Rightarrow p\equiv3\pmod4\}\),
every sufficiently large even target likewise has an explicit
asymptotic of order \(N/\log N\), with a positive logarithmic saving.
The subgroup and coset theorems allow fixed arithmetic data and retain
all surviving characters, including the complex branches forced by
the prime class \(3\pmod{10}\). Numerical saving exponents are
collected in Section~\ref{sec:assembly}.

One-variable density does not control the particular line \(a+b=N\),
and averages can conceal exceptional targets. The nonconstant linear
parts of \(n\) and \(N-n\) are dependent, so the independent-linear-form
correlation theorems of Loughran--Matthiesen do not directly apply
\cite[Theorems 1.9 and 3.10]{loughranmatthiesen}.
The reflected frequency argument must retain cancellation and
normalization when the target enters the rational denominators.

Iwaniec--Pomyka\l a's quartic-norm work is a direct antecedent
\cite{iwaniecpomykala}. For \(\Q(\zeta_{12})\), its splitting conditions
select the primes \(1\pmod{12}\). Their Theorem~2 gives conditional
pointwise order bounds for coprime sums; hypothesis (P4) is uniform
in growing shifts. Their Theorem~1 treats fixed differences
unconditionally. They count numerical values once.
Our unrestricted leading term gives a coprime asymptotic by M\"obius
inversion, with the extra target factor in Corollary~\ref{cor:coprime-c12}.

\subsection{Attribution and dependencies}

Walsh's multiplicative-invariance framework treats positive integer divisor
orders and fixed nonzero shifts \cite{walsh}. His Section~1.2.4, pp.~5--6,
proposes sums-of-two-squares indicator correlations. Here we supply their
complete character models, arithmetic hypothesis checks and local factors,
and the uniform reflected-target argument.

The recent inputs are the weighted-dilation root-moment theorem
\cite[Section 3]{weightedgraphs} and the all-conductor zero-free half-plane
\(\Re s>11/12\) \cite[Part I, Theorem 3.1]{zerofree}; Part II is not used.
Walsh notes a possible classical recovery of his fixed-shift theorem
\cite[p.~2]{walsh}, but his printed Lemma~5.1 and Section~6.4 retain that
zero-free input. Our polynomial-logarithmic fixed-shift moduli still occur
at Mellin heights up to \(x(\log x)^C\).
Section~\ref{sec:fixed-classical-boundary} proves the classical
intermediate-height estimate and states the remaining high-frequency
obligation. Both recent inputs remain in our fixed-shift result; such a
replacement would also leave the reflected growing-conductor problem.

Appendix~\ref{sec:older-route} proves classical quadratic order bounds
by Kloosterman completion and positive restoration. With the class-error
sieve these give genus comparison independently of the recent inputs.
All data are fixed; thresholds are not numerical. The companion paper
\cite{exactcompanion} gives exact classifications and finite certificates
for four primitive quadratic norm sumsets. The deposited window paper
concerns the placement of summands.

\subsection{Notation}

\begin{center}\small
\begin{tabular}{@{}>{\raggedright\arraybackslash}p{.28\linewidth}>{\raggedright\arraybackslash}p{.66\linewidth}@{}}
\hline
Symbol&Role\\\hline
\(X\), \(L=\log X\), \(h\)&Outer scale, its logarithm, and the fixed shift\\
\(\mathcal P_y\), \(n_{\mathcal P_y}\), \(V_y\)&Small primes, complete small-prime part, and rough-integer density\\
\(f_i,\mathcal D_i,\xi_i\)&Indicators, arithmetic models, and centered errors\\
\(D,D_i\), \(U_i\)&Auxiliary prime products and their dyadic lengths\\
\(H_i\), \(\mathscr A\), \(\mathcal E\)&Residue subgroups, a prime coset, and forbidden prime sets\\
\(B_0\), \(S\)&Finite exceptional-prime sets in the ordinary-indicator and primitive-form applications\\
\(\Delta_i=f_i^2D\)&Order discriminants sharing one fixed fundamental discriminant \(D\); here \(f_i\) denotes a conductor\\
\(P_Q,\Lambda_Q\), \(G\)&Primitive numerical values, their local shadow, and the applicable class group\\\hline
\end{tabular}
\end{center}
The modulus and the arithmetic data stay fixed. Definitions in each
application specify its groups and exceptional primes; \(h\) is
reserved for the additive shift throughout.

\setcounter{tocdepth}{1}
\begingroup\small
\makeatletter
\renewcommand*\l@section{\@dottedtocline{1}{0em}{2em}}
\makeatother
\tableofcontents
\endgroup
\section{Replacing indicators by arithmetic models}\label{sec:comparison}

We isolate the prime invariance, one-variable means and sparse sieve
bounds needed to replace an indicator by its arithmetic model.

Fix an even integer \(L_0\ge2\) and a reduced class
\(c\in(\Z/L_0\Z)^\times\). At the outer scale \(X\), write
\[
 L=\log X,\qquad y=\exp(L^{2/5}),\qquad
 \mathcal P_y=\{p:p\le y\},\qquad E=\exp(\sqrt L).
\]
For a set of primes \(\mathcal P\), let \(n_{\mathcal P}\) be the
part of \(n\) supported on \(\mathcal P\), including its full
valuations. A rough integer has no prime factor in \(\mathcal P_y\).
There are two families of lines in the theorem:
\begin{align*}
 \mathcal L_X^{\rm sum}&=\{(u,v)\in\Z_{>0}^2:Au+Bv=X\},
       &t(u,v)&=Au, &J_X&=[0,X],\\
 \mathcal L_X^{\rm shift}&=\{(u,u+h):u\in\Z,\ X\le u\le2X\},
       &t(u,v)&=u, &J_X&=[X,2X].
\end{align*}
In the first family \(X\) is an integer target and \(A,B\) are
fixed coprime positive integers supported on primes dividing \(L_0\).
In the second, \(h\ne0\) is a fixed integer and \(X\) tends to
infinity; both coordinates are then positive. Set \(k=X\) in the
sum family and \(k=h\) in the shift family. Dependence on the
fixed data, including \(h\) in the second family, is permitted.

Let \(0\le f_i\le1\), \(i=1,2\), be fixed functions on the
positive integers. For every sufficiently large fixed integer \(K\),
let \(\mathcal D_{i,X,K}\) be an arithmetic model, and write
\(\xi_{i,X,K}=f_i-\mathcal D_{i,X,K}\). We often suppress \(X,K\).
The entire model, its cutoff \(y\), and its small-prime records stay
tied to this outer scale during all subsequent dilations.

Choose a positive comparison scale \(\mathcal S_X\) satisfying
\begin{equation}\label{eq:comparison-scale}
        XL^{-\kappa_0}\le\mathcal S_X\le XL^{\kappa_0}
\end{equation}
for a fixed \(\kappa_0\ge0\). Arithmetic factors may be included
in this scale. The hypotheses are as follows.

\begin{description}
\item[\textup{(C1)} Bounded models.]
For every fixed \(K\), sufficiently large \(X\) gives
\(0\le \mathcal D_{i,X,K}(n)\le1\) on \(1\le n\le X^3\).
The bound one is independent of the model depth; its threshold
may depend on \(K\).

\item[\textup{(C2)} A common prime class.]
For every prime \(p\le y\) with \(p\equiv c\pmod{L_0}\),
\[
 f_i(pn)=f_i(n),\qquad \mathcal D_{i,X,K}(pn)=\mathcal D_{i,X,K}(n)
\]
whenever the arguments are in the range under consideration.
This includes repeated multiplication by \(p\). A residue filter
is part of \(f_i\) and must satisfy the same equality.

\item[\textup{(C3)} Centered one-variable rough means.]
For every fixed \(B_0,\mathcal D_0>0\), there is a depth \(K_0\) such that
every fixed \(K\ge K_0\) has the following property. Uniformly for
\(\mathcal P_y\)-supported integers \(r\le E\),
\(X^{1/2}\le Y\le X^2\), bounded periodic functions \(a(v)\)
of period at most \(L^{B_0}\), and smooth functions \(W\) supported
in a fixed compact subinterval of \((0,\infty)\),
\begin{equation}\label{eq:comparison-rough-means}
 \left|\sum_{v\ {\mathrm{rough}}}
       \xi_{i,X,K}(rv)a(v)W(rv/Y)\right|
                       \ll_{B_0,\mathcal D_0,K}\frac Yr L^{-\mathcal D_0}.
\end{equation}
Here \(|a|\le1\), and every fixed logarithmic derivative of \(W\)
used in the estimate is at most a fixed power of \(L^{B_0}\).
The hypothesis permits Mellin factors with polynomial-logarithmic
frequency, by including them in \(W\). Its constants may depend on
the chosen derivative orders. Every surviving character branch,
including a negative or complex pole parameter, must be retained
when verifying this estimate.

\item[\textup{(C4)} A positive sparse record envelope.]
Take disjoint prime banks in the two ranges
\[
 [e^{L^{1/8}},e^{L^{1/4}}],\qquad
 [e^{L^{1/3}},e^{L^{2/5}}],
\]
each bank consisting of the primes in an interval which are
\(c\pmod{L_0}\), after deleting primes dividing \(k\). Assume each
bank has reciprocal mass between \(1/2\) and \(2\), and that the
number in either range is at most \((\log L)/(16\varphi(L_0))\).
Write \(V_g=\sum_{p\in\mathcal P_g}p^{-1}\) for its reciprocal mass.
Let \(\mathcal B\) be their union. Define the independent joint
record law \(\mu_{\mathcal B}=\prod_{p\in\mathcal B}\mu_p\) by
\begin{equation}\label{eq:bank-law}
 \mu_p(0,0)=1-2/p,\qquad
 \mu_p(e,0)=\mu_p(0,e)=(1-1/p)p^{-e}\quad(e\ge1),
\end{equation}
and zero mass on two positive valuations.
For every fixed \(B_0>0\), nonnegative record function
\(F(r,s)\le L^{B_0}\), interval \(J\subset J_X\), and
\(w_i\in\{f_i,\mathcal D_{i,X,K}\}\), require
\begin{align}
 &\sum_{\substack{(u,v)\in\mathcal L_X\\t(u,v)\in J}}
     F(u_{\mathcal B},v_{\mathcal B})w_1(u)w_2(v)
       \ll_{B_0,K}\frac{|J|}{X}\mathcal S_X
                              \E_{\mu_{\mathcal B}}F\notag\\
 &\hspace{34mm}
       +X^{1/2}L^{C_{B_0}}
       +XL^{C_{B_0}}\exp(-c_0L^{1/10}).
                                                        \label{eq:comparison-envelope}
\end{align}
The positive constant \(c_0\) and exponent \(C_{B_0}\) are
independent of the final fixed model depth. Constants and thresholds
may depend on that depth. Uniformity includes the bank intervals,
the exact valuations in the record, and the interval endpoints.
\end{description}

For a bank \(g\), let \(\omega_g(n)\) count its distinct prime
divisors, and define
\[
 I_g(u,v)=\1_{\omega_g(u)=\omega_g(v)=1},\qquad
 \Omega_T(n)=\prod_{g\in T}V_g^{-1}\1_{\omega_g(n)=1}.
\]
Here \(T\) is a set of bank indices. We also write
\(\Omega_{\mathcal P}\) when \(\mathcal P\) denotes the union of
the active banks. The joint valuation law above
is the ordinary CRT law on either line: every bank prime avoids
\(k\) and is a unit at the fixed coefficients.

Condition (C3) is a statement about one variable on a progression;
it allows arbitrary retained small-prime records, but no arbitrary
weight in the rough variable. Condition (C4) is only an upper bound
with nonnegative weights. In the applications it follows from CRT
root counts and an upper sieve, with a power-smaller endpoint error.
Neither condition assumes centered binary cancellation or a major-arc
energy estimate. In particular, choosing a model with the right
average density alone does not verify (C3).

\begin{theorem}[Arithmetic model comparison]\label{thm:model-comparison}
Assume the zero-free and weighted-dilation inputs of
Section~\ref{sec:inputs}, and conditions \textup{(C1)--(C4)}.
Put \(\delta=1/(8192\varphi(L_0))\). For every sufficiently deep
fixed \(K\), uniformly for intervals \(J\subset J_X\),
\begin{equation}\label{eq:model-comparison-conclusion}
 \sum_{\substack{(u,v)\in\mathcal L_X\\t(u,v)\in J}}
     \bigl(f_1(u)f_2(v)-\mathcal D_{1,X,K}(u)\mathcal D_{2,X,K}(v)\bigr)
                       \ll_K\mathcal S_XL^{-\delta}.
\end{equation}
For the sum family the estimate is uniform in the integer target
and its prime valuations. For the shift family the integer \(h\)
is fixed, and the constant may depend on it. The two families have
different major-arc geometries; no uniformity in the intervening
range of growing sublinear shifts is asserted.
\end{theorem}

The proof consists of the shared graph reduction
(Section~\ref{sec:graph}), the fixed-shift and reflected estimates
(Sections~\ref{sec:fixed-shift-comparison}--\ref{sec:reflected-comparison}),
and reconstruction (Section~\ref{sec:assembly}).
Subgroup pairs satisfy (C2) in the identity class; a self-coset uses a
class in that coset. Arbitrary prescriptions may have no common class.
The final asymptotic saving must also pay for evaluating the model.

\section{Analytic inputs and the comparison framework}\label{sec:inputs}

We use Walsh's comparison architecture \cite[Sections 2--7]{walsh},
the zero-free input \cite[Theorem 3.1]{zerofree}, and the
weighted-dilation transfer \cite[Section 3]{weightedgraphs}.
The operator normalizations and parameter order are stated explicitly.

\subsection{The zero-free input}

\begin{proposition}[Imported zero-free theorem]\label{prop:zf}
Every Dirichlet \(L\)-function, including \(\zeta\), is nonzero in
\[
                         \Re s>11/12.
\]
The pole at \(s=1\) of a principal character is allowed. Both conductor
and height are unrestricted.
\end{proposition}

This is \cite[Theorem 3.1, printed p.~10]{zerofree}, restricted to
Dirichlet characters and the open half-plane. We also use the usual
functional equation, Euler product and holomorphic logarithm on
zero-free rectangles.

\subsection{Prime groups and the two operators}

For this subsection only, write
\[
                \Lambda=\log x,\qquad T=\log\Lambda.
\]
Fix constants
\[
 0<\alpha<\beta<\gamma<\delta<0.47,\quad
 c_0,C_0,v_-,v_+>0,\quad 0<\vartheta<1,\quad \ell\ge1.
\]
There are disjoint collections \(\mathcal G_s,\mathcal G_b\) of
pairwise disjoint finite prime sets \(P_g\), with
\[
 c_0T\le|\mathcal G_s|,|\mathcal G_b|\le C_0T.
\]
They satisfy
\begin{equation}
\begin{split}
 e^{\Lambda^\alpha}\le p\le e^{\Lambda^\beta}
       &\quad(g\in\mathcal G_s,\ p\in P_g),\\
 e^{\Lambda^\gamma}\le p\le e^{\Lambda^\delta}
       &\quad(g\in\mathcal G_b,\ p\in P_g),\\
 v_-\le V_g:=\sum_{p\in P_g}\frac1p\le v_+,
       &\qquad \mu_g(p)=\frac1{V_gp}.
\end{split}                                                    \label{eq:inputbanks}
\end{equation}
The sets need not be intervals. This freedom is needed when taking
identity residue classes and deleting primes dividing the target.

Fix a sufficiently large integer \(J\), and put \(M=J+\ell\).
A pattern \(\nu\) specifies subsets
\(C_g\subseteq\{1,\ldots,J\}\), with
\[
 C_g=\varnothing\quad(g\in\mathcal G_s),\qquad
 |C_g|\le m\quad(g\in\mathcal G_b),\qquad t_g=\ell+|C_g|.
\]
The slots in \(\{1,\ldots,J\}\setminus C_g\) are shared between the
two endpoint lists. The remaining \(t_g\) slots at each endpoint
are unshared. Form \(D\) from the shared labels and, in the
\(C_g\)-slots, fresh independent draws with laws \(\mu_g\).
Thus \(D\) has \(J\) prime factors from each group, counted with
multiplicity. Write \(D=D_sD_b\) according to group size.

Each pattern has a complex coefficient \(K_\nu\). It may depend on
the ordered big-group source labels, target labels, and auxiliary
labels used in \(D_b\), and on no small-group labels. Assume
\begin{equation}
                 \sum_\nu\sup |K_\nu|\le\Lambda^{C_1},           \label{eq:inputmass}
\end{equation}
where \(m,C_1\) are fixed independently of \(J\). All operators below
average source and target slot permutations independently in each
applicable group. These averages are orthogonal projections.

The \emph{ideal operator} acts on
\[
 \mathcal H_{\mathrm{id}}=
 L^2\!\left(\prod_{g\in\mathcal G_b}P_g^M,
                    \bigotimes_{g\in\mathcal G_b}\mu_g^{\otimes M}\right).
\]
Numerical repetitions are permitted. Its row operation retains shared
labels, samples the unshared target labels and auxiliary labels
independently, and multiplies by
\[
                         K_\nu D_b^{\,i\zeta}e(\Theta D_b),
             \qquad e(t)=\exp(2\pi it).
\]
Sum over patterns and compose with the permutation projections on
both sides; denote the result by \(\mathcal T(\Theta)\).

A \emph{physical state} consists of an integer \(n\) and an ordered
\(M\)-tuple in every group, with distinct labels within each group,
all dividing \(n\). Give every such state at \(n\) mass
\(\prod_gV_g^{-M}\), and use counting measure on the integer position.
This defines \(\mathcal H_{\mathrm{ph}}\). Put
\[
 \omega_g(n)=\#\{p\in P_g:p\mid n\},\quad
 \omega(n)=\sum_g\omega_g(n),\quad
 s=|\mathcal G_s|+|\mathcal G_b|.
\]
For a fixed integer \(k\) with \(0<|k|\le\Lambda^{C_2}\), the
\emph{physical operator} \(A\) moves from \(n\) to \(n'=n+kD\).
It sums over the physical target lists with the prescribed shared
labels, using row normalization \(\prod_gV_g^{-t_g}\); auxiliary
labels retain their independent probability laws. Every auxiliary
label is forbidden from both endpoint lists, although auxiliary
labels may equal each other. Multiply the row by
\begin{equation}
 K_\nu D^{\,i\zeta}
       \vartheta^{(\omega(n)-Ms)/2}
       \vartheta^{(\omega(n')-Ms)/2},                            \label{eq:inputphysical}
\end{equation}
sum over patterns and apply all endpoint permutation projections.
The adjoint reverses the displacement and conjugates the coefficient
with the endpoint roles interchanged. The joint endpoint
normalization in each group is \(V_g^{-M-t_g}\) in either direction.

\begin{proposition}[Imported absolute bounds]\label{prop:graphabsolute}
Replace every elementary transition coefficient of \(A\) by its
absolute value before summing patterns and averaging permutations.
The resulting operator and its adjoint have row sums at most
\(\Lambda^{C_{\mathrm{abs}}}\). The exponent may depend on \(m,C_1\)
and the fixed group data, but is independent of \(J\).
The same assertion holds with integer position replaced by independent
uniform residues at all group primes.
\end{proposition}

This is \cite[Lemma 3.4]{weightedgraphs}. The independence from \(J\)
is part of the assertion. To see its source, once the shared labels
are fixed, the new target slots in a group cost at most
\[
       \sup_{z\ge0}V_g^{-t_g}(z+t_g)^{t_g}\vartheta^{z/2}.
\]
Here \(t_g\le\ell+m\). This is a fixed constant per group and its
product over \(O(T)\) groups is a fixed power of \(\Lambda\).
Reversing the joint endpoint measure gives the column bound.

\subsection{Root moments, residue replacement, and endpoints}

\begin{proposition}[Imported local transference]\label{prop:rootmoment}
For every fixed \(E>0\), there is \(E_{\mathrm{id}}>0\), depending
only on \(E\) and the fixed group data, such that the following holds.
Suppose
\begin{equation}
                   \sup_{\Theta\in\mathbb R}
                        \|\mathcal T(\Theta)\|_{2\to2}
                              \le\Lambda^{-E_{\mathrm{id}}}.
                                                                    \label{eq:inputideal}
\end{equation}
After \(m,C_1\) have been fixed, take a sufficiently large fixed \(J\).
Put \(R=\lceil\Lambda^{0.52}\rceil\). For every fixed \(\eta>0\)
and \(I=[X,2X)\), with \(x^\eta\le X\le x^{1/\eta}\),
\begin{equation}
 \frac1{|I\cap\Z|}\sum_{n\in I\cap\Z}
       \langle u_n,(AA^*)^Ru_n\rangle
                         \le\Lambda^{-2RE}.                    \label{eq:inputroot}
\end{equation}
Here \(u_n\) is one on every physical mark list at position \(n\)
and zero elsewhere. The threshold may depend on all the fixed
parameters, including \(J,C_2,\eta\). No bound on \(\zeta\) is
required beyond \eqref{eq:inputideal}.
\end{proposition}

This is \cite[Theorem 3.5]{weightedgraphs}. It imposes return of the
integer position. The final mark list need not equal the initial
list, and \(u_n\) is not normalized. Substituting a trace with
identical initial and final mark lists would change the statement.

The arithmetic replacement inside that theorem has the following
precise form, which also explains its scale limitation.
For a fixed compatible path with \(r=2R\) edges, let \(h_i\) be
the integer offsets and \(\mathcal A_i\) its active prime set at
visit \(i\). Set \(\vartheta_i=\vartheta\) at internal visits and
\(\vartheta_0=\vartheta_r=\sqrt{\vartheta}\).
The position-dependent factor is
\begin{equation}
 W(n)=\prod_{i=0}^{r}
       \left(\prod_{p\in\mathcal A_i}1_{p\mid n+h_i}\right)
       \vartheta_i^{\,\#\{p\notin\mathcal A_i:
                     \ p\text{ is a group prime},\ p\mid n+h_i\}}.
                                                                    \label{eq:inputpath}
\end{equation}

\begin{proposition}[Imported residue replacement]\label{prop:graphcrt}
For every fixed \(B>0\), the average of \eqref{eq:inputpath} on
\(I\cap\Z\) equals its expectation with independent uniform residues
at all group primes, multiplied by \(1+O(\exp(-BrT))\).
An incompatible collection of active congruences has zero weight
in both models.
\end{proposition}

This is \cite[Lemma 3.6]{weightedgraphs}. For clarity, its relative
form retains all deterministic damping belonging to active primes.
The product \(Q_{\mathrm{act}}\) of active primes satisfies
\[
              \log Q_{\mathrm{act}}
                   =O_J(rT\Lambda^\delta)=o(\Lambda).
\]
After those factors are extracted, a Bonferroni expansion of degree
\(O_B(rT)\) controls the remaining damping. Its independent
expectation is at least \(\exp(-O(rT))\); its total interval
counting error is \(\exp(-\eta\Lambda+o(\Lambda))\).
The inequality \(0.52+\delta<1\) makes that error negligible
even after division by the active congruence density. Together
with Proposition~\ref{prop:graphabsolute}, this permits summation
over paths with absolute coefficients. These observations describe
the imported replacement and do not assert replacement on intervals
shorter than its stated range.

\begin{corollary}[Imported endpoint consequence]\label{cor:graphendpoint}
Under Proposition~\ref{prop:rootmoment}, let \(f\in\mathcal H_{\mathrm{ph}}\)
be supported on positions \(I=[X,2X)\), constant on the chosen mark
list at each position, with value \(f_n\). If
\[
 \sup_{n\in I}|f_n|\le\exp(O(\sqrt{\Lambda})),\qquad
          \|f\|,\|g\|\le X^{1/2}\Lambda^{C_3},
\]
then
\begin{equation}
                    |\langle f,Ag\rangle|
                        \ll X\Lambda^{-E+2C_3}.                 \label{eq:inputendpoint}
\end{equation}
Only \(f\) must be constant on mark lists.
\end{corollary}

This is \cite[Corollary 3.11]{weightedgraphs}. Its proof partitions
\(I\) into intervals on the displacement scale
\(\exp(O_J(T\Lambda^\delta))\), uses spectral H\"older for
\(AA^*\), and bounds the positive matrix
\((\langle u_n,(AA^*)^Ru_m\rangle)\) by its trace.
It makes no equal-fiber-mass assumption.

\subsection{The application interfaces}

Three steps remain: construct the fractional-density models retaining
every surviving character; prove the progression-filter ideal comparison
with its big-label coefficient bound \eqref{eq:inputmass} fixed before
\(J\); and transfer to the growing reflected target with uniform endpoint
and localization bounds. The root-moment theorem supplies none of these
steps on its own. Sections~\ref{sec:graph}--\ref{sec:assembly} prove the
comparison; the applications then check its arithmetic hypotheses.

\section{The shared graph comparison}\label{sec:graph}

This section proves the graph reduction common to the reflected and
fixed-shift families of Theorem~\ref{thm:model-comparison}. Its output is
\eqref{eq:gr-comparison}--\eqref{eq:gr-comparison-term}: every error
correlation is replaced by comparisons containing two independent
large prime products. Sections~\ref{sec:fixed-shift-comparison} and
\ref{sec:reflected-comparison} estimate those comparisons in their
different frequency windows. No arithmetic form of the indicators
is used here beyond \textup{(C1)--(C3)}.

The labels lie in one fixed reduced class. The character frame must
therefore identify restrictions differing by a known unit phase.

Write \(L=\log x\).  There are \(\asymp\log L\) small banks and big
banks, with primes respectively in
\[
 [\exp(L^{1/8}),\exp(L^{1/4})],\qquad
 [\exp(L^{1/3}),\exp(L^{2/5})].
\]
Every label is \(c\bmod L_0\); primes dividing the nonzero displacement
\(k\) are removed. We use \(k=T=x\) for sums and fixed \(k=h\)
for shifts. All banks are disjoint. Their laws are
\(\mu_g(p)=(V_gp)^{-1}\), where \(V_g=\sum_{p\in\mathcal P_g}p^{-1}\)
is bounded above and below by positive fixed constants.
Banks in this section may be any active subfamilies of the banks
constructed from the class in \textup{(C2)}, provided that each size has
at least a fixed positive multiple of \(\log L\) banks.
All constants may depend on these fixed bounds and on \(L_0\).

\subsection{The character frame for progression banks}

The ideal-kernel construction in \cite[Section~4]{weightedgraphs}
uses primitive characters of conductor at most \(Q=L^{q_*}\).
On labels \(c\bmod L_0\), this list is redundant up to phases.
For example, the principal character and \(\chi_{-4}\) agree on
primes \(1\bmod4\) and are opposite on primes \(3\bmod4\).
The quotient below removes this redundancy while retaining its phase.

Let \(\mathcal X_Q\) be the primitive characters of conductor at most
\(Q\), and define
\begin{equation}
 \eta\sim\eta'
 \quad\Longleftrightarrow\quad
 \operatorname{cond}(\eta\overline{\eta'})\mid L_0.
 \label{eq:gr-equivalence}
\end{equation}
The conductor here belongs to the primitive character inducing the
quotient on a common modulus.  Choose one representative from each
class and denote the resulting set by \(\mathcal R_Q\).
Thus \(|\mathcal R_Q|\le Q^2\).

To verify the equivalence and its meaning, put
\(M_0=\lcm(L_0,\operatorname{cond}\eta,\operatorname{cond}\eta')\).
Reduction from \((\Z/M_0\Z)^\times\) onto
\((\Z/L_0\Z)^\times\) is surjective.  Its kernel consists of the
units \(1\bmod L_0\).  A character is trivial on that kernel exactly
when it factors through the latter group.  This is equivalent to
its primitive conductor dividing \(L_0\). Consequently, if
\(\eta\sim\eta'\) and their quotient is induced by
\(\vartheta\bmod L_0\), then on a product of \(r\) labels,
\[
                 \eta(D)=\vartheta(c)^r\eta'(D).
\]
Both comparison products have the same slot type and the same phase.
Inequivalent characters have zero average on the kernel, and also
on any fixed coset of it: translation multiplies their sum by a unit phase.

Fix a nonempty set \(\Lambda\) of at most \(m\) slots per big bank,
where \(m\) is fixed before the eventual graph mark length.
Let \(D_\Lambda\) be the product of independent labels in these slots.
Partition its logarithm into cells of width \(\tau>0\)
\(I_j=[j\tau,(j+1)\tau)\), and write
\(\nu_j=\Prob(\log D_\Lambda\in I_j)\).
For inequivalent \(\eta,\eta'\in\mathcal R_Q\), and any prescribed
\(A>0\), Siegel--Walfisz gives
\begin{equation}
 \left|\E\big[
  \1_{\log D_\Lambda\in I_j}
        \eta(D_\Lambda)\overline{\eta'(D_\Lambda)}\big]\right|
       \ll_A L^{-A}.                                  \label{eq:gr-cell}
\end{equation}
Here and below the accuracy in Siegel--Walfisz is chosen after all
fixed logarithmic exponents.

Indeed, freeze all but one label.  Its allowable values form an
interval inside its original bank interval, still restricted to
\(c\bmod L_0\). On a common modulus at most \(L_0Q^2\), the sum of
\(\eta\overline{\eta'}\) over the appropriate reduced residue classes
is zero.  Siegel--Walfisz, summed over these classes, therefore gives
\[
 \sum_{\substack{p\le z\\p\equiv c\ (L_0)}}
       \eta(p)\overline{\eta'(p)}
       \ll_K L_0Q^2z(\log z)^{-K}.
\]
Partial summation bounds the reciprocal-prime sum on any subinterval
of \([\exp(L^{1/3}),\exp(L^{2/5})]\) by
\(O_K(L_0Q^2(L^{1/3})^{1-K})\).
This is an absolute bound; no lower bound for the subinterval's
length or harmonic mass is needed.  Divide by \(V_g\), then average
the frozen labels.  Deleting target divisors changes the answer by at most
\begin{equation}
 \sum_{\substack{p\mid k\\p\ge\exp(L^{1/3})}}\frac1p
 \le \frac{\log(2|k|)}{L^{1/3}}\exp(-L^{1/3}),               \label{eq:gr-deletion}
\end{equation}
which is smaller than every fixed inverse power of \(L\).
This proves \eqref{eq:gr-cell} uniformly for the actual banks.

On cells with \(\nu_j\ge\xi\), define
\begin{equation}
 K_\Lambda(D,X)=
 \sum_{\nu_j\ge\xi}
 \frac{\1_{\log D\in I_j}\1_{\log X\in I_j}}{\nu_j}
 \sum_{\eta\in\mathcal R_Q}\overline{\eta(D)}\eta(X).
 \label{eq:gr-kernel}
\end{equation}
Both product variables have the same law.  The elementary bounds are
\begin{equation}
 |K_\Lambda|\le Q^2\xi^{-1},\qquad
 \sup_X\E_D|K_\Lambda(D,X)|\le Q^2,\qquad
 \#\{j:\nu_j>0\}\ll_m1+(\log L)L^{2/5}/\tau.              \label{eq:gr-kernel-bounds}
\end{equation}
The last inequality uses only the range of the logarithm.

On a retained cell, the kernel is the positive frame operator
\[
 f\longmapsto
 \sum_{\eta\in\mathcal R_Q}\eta\,\langle\eta,f\rangle
\]
for the conditional product probability measure.
Its Gram matrix has diagonal entries one, since all labels exceed
\(L_0Q^2\), and off-diagonal entries
\(O_A(\xi^{-1}L^{-A})\) by \eqref{eq:gr-cell}.
Thus its norm is at most
\(1+O_A(Q^2\xi^{-1}L^{-A})\le2\).
Distinct cells are orthogonal blocks, so
\begin{equation}
             \|K_\Lambda\|_{2\to2}\le2.                \label{eq:gr-frame-norm}
\end{equation}
This also holds for tuple tests: conditional expectation from tuples
to their product is a contraction.
If \(\psi\) is a character of modulus \(r\le Q\), its primitive
inducing character is equivalent to one member of \(\mathcal R_Q\).
For \(r_\Lambda\) selected slots and \(\psi=\vartheta\eta\) on
these units, that summand gives
\(\vartheta(c)^{r_\Lambda}\eta(X)=\psi(X)\). The other summands give
\begin{equation}
 \E_D K_\Lambda(D,X)\psi(D)
       =\psi(X)+O_A(Q^2\xi^{-1}L^{-A})                 \label{eq:gr-reproduce}
\end{equation}
at every retained \(X\).  This is where quotienting the frame is
essential: the unquotiented frame can reproduce the same constant
several times.

\begin{lemma}[Progression ideal comparison]\label{lem:gr-ideal}
For every prescribed ideal accuracy \(E_{\rm id}\) and fixed phase
range \(|\zeta|\le L^{C_4}\), the progression banks admit the signed
comparison family of \cite[Theorem~4.1]{weightedgraphs}.
Its ideal norm is at most \(L^{-E_{\rm id}}\), uniformly in every
real additive frequency.  Its coefficients ignore small-bank labels
and shared labels, free at most \(m\) slots per big bank, and have
total supremum at most \(L^{C_1}\).
The parameters \(m,C_1\) and all kernel powers are fixed before the
eventual mark length.
\end{lemma}

\begin{proof}
We give the comparison estimates used in the construction, including
the only step where restricting the prime labels changes the proof.

For a product formed from \(k_g\le m\) draws in each of at most
\(C_0\log L\) disjoint banks, unique factorization gives its exact law
\begin{equation}
 n\Prob(D_\Lambda=n)
   =\prod_g\frac{k_g!}{V_g^{k_g}\prod_{p\in\mathcal P_g}a_p!},
 \qquad
 \Prob(D_\Lambda=n)\le L^{K_0}/n,                     \label{eq:gr-productlaw}
\end{equation}
where
\(K_0=C_0(\log(m!)+m\log\max(1,v_-^{-1}))\) and
\(V_g\ge v_->0\).
The factor \(C_0\) multiplies both terms of this exponent.
For a tuple test, Cauchy--Schwarz on each product fibre gives an
integer coefficient sequence of squared norm at most
\(L^{K_0}/U\) times the tuple squared norm on a product dyad
\([U,2U)\).

The elementary integer bilinear estimate
\cite[Lemma~4.2]{weightedgraphs}, or \cite[Lemma~6.2]{walsh}, gives
\begin{equation}
 \left|\sum_{u,v}a_ub_v\ee{\theta uv}\right|^2
 \ll \|a\|_2^2\|b\|_2^2
             (1+U/r)(V+r\log(2r))                  \label{eq:gr-integer-bilinear}
\end{equation}
when the two supports lie in containing intervals of lengths
\(O(U),O(V)\), \((a,r)=1\), and
\(|\theta-a/r|\le r^{-2}\).
For completeness, Cauchy--Schwarz and expansion leave
\(V+\sum_{1\le j\ll U}\min(V,\|\theta j\|^{-1})\).
Divide the \(j\)'s into blocks of diameter less than \(r/2\).
Within a block the rational parts are \(1/r\) separated and their
perturbations differ by at most \(1/(2r)\); hence the sum is
\(O(V+r\log(2r))\) per block.  This proves
\eqref{eq:gr-integer-bilinear} and permits arbitrary subsets of the
containing intervals.

Apply this estimate to two nonempty selected product sets \(A,B\).
On product dyads \(U,V\), use Dirichlet approximation with denominator
ceiling \(\lfloor UV/S\rfloor\).  When the resulting denominator
exceeds \(Q\), \eqref{eq:gr-productlaw} gives the normalized bound
\begin{equation}
 \ll L^{K_0}\big(o(L^{-A})+Q^{-1}+L^{7/5}/S\big)^{1/2}.
 \label{eq:gr-highden}
\end{equation}
There are at most \(L^{7/5}\) product dyads per variable.
The same bound applies to the kernel terms by moving each kernel
onto its endpoint test using \eqref{eq:gr-frame-norm}.
The factor \((uv)^{i\zeta}\) separates and preserves test norms.

For denominators \(r\le Q\), all products are units modulo \(r\).
The multiplicative-character expansion of the rational phase has
total absolute coefficient at most \(\sqrt{\varphi(r)}\).
Replace each primitive character by its representative in
\(\mathcal R_Q\) and apply \eqref{eq:gr-reproduce}.
Grouping characters after inserting their unit phases cannot increase
coefficient mass, so the reproduction error is
\(O_A(Q^{5/2}\xi^{-1}L^{-A})\).
Write the remaining phase as
\(M(v)=\exp(i\zeta v)\ee{\beta\exp v}\).
On matched cells its derivative is
\(O(L^{C_4}+S)\), and the two logarithmic products move by at most
\(2\tau\).  By \eqref{eq:gr-kernel-bounds}, this costs
\(O(Q^4\tau(S+L^{C_4}))\).
The union of discarded cells has probability at most
\((N_A+N_B)\xi\), where \(N_A,N_B\) count the possible cells.
Cauchy--Schwarz bounds its contribution by the square root of this
probability times the two test norms.
These estimates prove the two-product comparison, for arbitrary
\(L^2\) tuple tests, to any prescribed power \(L^{-A_0}\).

There is no circular choice of precisions.  One possible hierarchy is
\[
 \begin{split}
 B_0&>A_0+2(7/5)+10,\\
 q_*,s_*&>2(B_0+K_0+10)+7/5,\\
 \tau_*&>4q_*+\max(s_*,C_4)+B_0+10,\\
 x_*&>2B_0+7/5+\tau_*+10,
 \end{split}
\]
with \(Q=L^{q_*},S=L^{s_*},\tau=L^{-\tau_*},\xi=L^{-x_*}\).
Choose the accuracy in \eqref{eq:gr-cell} afterwards.
All these choices are independent of the eventual mark length.

Here is the signed cancellation which completes the ideal construction.
Split the big banks into two blocks of comparable size and designate
\(m\) probe slots per group.  For nonempty probe subsets
\(A,B\) in the two blocks, free their slots and give the comparison
coefficient
\begin{equation}
              -(-1)^{|A|+|B|}K_A(D_A,X_A)K_B(D_B,X_B).
 \label{eq:gr-signed}
\end{equation}
Here \(D_A,D_B\) are auxiliary-label products and \(X_A,X_B\)
are the unshared endpoint products.  The coefficient has no small
or shared labels.  There are at most \(2^{mC_0\log L}\) comparisons;
by \eqref{eq:gr-kernel-bounds} their total supremum is at most
\begin{equation}
        1+2^{mC_0\log L}Q^4\xi^{-2}\le L^{C_1}.         \label{eq:gr-coeffcost}
\end{equation}
Expanding both kernels also separates each coefficient into a sum
of products of four bounded factors: one on each auxiliary product
and one on each unshared endpoint list.  The number of terms and
total scalar coefficient mass are fixed powers of \(L\).

We now make the exceptional projection and the cancellation explicit.
Write the two big-bank blocks as \(\mathcal B_1,\mathcal B_2\), of
sizes \(s_1,s_2\), and their designated probe coordinates as
\(\mathcal I_1,\mathcal I_2\).  For a coordinate \(i\), let \(E_i\)
average its independent label.  These commuting operators are
orthogonal projections.  Let \(Q_{\rm tail}\) average the last
\(\ell\) input or target coordinates in every group, and let
\(\Pi\) be within-group symmetrization on all \(M=J+\ell\) coordinates.
The ideal multipliers do not use those last coordinates.
For symmetric endpoint tests \(F,H\), put
\(F'=Q_{\rm tail}F,\ H'=Q_{\rm tail}H\).
Both averaging operations are contractions.

Let \(D_{\rm sh}\) be the product of the first \(J\) coordinates,
and write
\[
 \Phi_{\Theta,\zeta}
       =D_{\rm sh}^{i\zeta}\ee{\Theta D_{\rm sh}},\qquad
 R(F',H')=\E\big[F'\overline{H'}\Phi_{\Theta,\zeta}\big].
\]
Thus \(R\) is a bilinear pairing of norm at most one.
Resolve the probes in the first block of \(F'\) and the second
block of \(H'\):
\[
 \begin{split}
 F'&=\sum_{S\subset\mathcal I_1}F_S,\qquad
 F_S=\prod_{i\in S}(I-E_i)
             \prod_{i\in\mathcal I_1\setminus S}E_iF',\\
 H'&=\sum_{T\subset\mathcal I_2}H_T,\qquad
 H_T=\prod_{i\in T}(I-E_i)
             \prod_{i\in\mathcal I_2\setminus T}E_iH'.
 \end{split}
\]
These are orthogonal decompositions.  Define
\(P_1=\prod_{i\in\mathcal I_1}(I-E_i)\),
\(P_2=\prod_{i\in\mathcal I_2}(I-E_i)\), so that
\(F_{\rm full}=P_1F'\), \(H_{\rm full}=P_2H'\).

In one group the original symmetric vector decomposes orthogonally
over subsets of its \(M\) coordinates.  Symmetry makes the squared
norms equal among subsets of cardinality \(j\).
Requiring all \(m\) probes to occur and all last \(\ell\) slots
to be absent therefore retains exactly the fraction
\[
 \frac{\binom{M-m-\ell}{j-m}}{\binom Mj}
     =\frac{(j)_m(M-j)_\ell}{(M)_{m+\ell}}
     \le 2(j/M)^m(1-j/M)^\ell
     \le2\ell^\ell m^{-\ell}=:\theta_m
\]
when \(M\ge(m+\ell)(m+\ell-1)\); impossible cardinalities give zero.
For a vector which is not a tensor product across groups, first
resolve by the tuple of cardinalities in the groups of a block.
Independent within-group permutations make all subset tuples of
those cardinalities have equal squared norms.  The retained fraction
is consequently the product of the displayed fractions.
Summing over the cardinality tuples proves the operator estimates
\begin{equation}\label{eq:gr-blockprojection}
 \|P_1Q_{\rm tail}\Pi\|\le\theta_m^{s_1/2},\qquad
 \|P_2Q_{\rm tail}\Pi\|\le\theta_m^{s_2/2}.
\end{equation}
Averaging tails outside the block is a further contraction.
We use \(\ell=1\), and each \(s_i\ge c_0\log L/3\).

For a component pair \(F_S,H_T\), put
\(U=\mathcal I_1\setminus S\), \(V=\mathcal I_2\setminus T\).
A comparison freeing \(A\subset\mathcal I_1\),
\(B\subset\mathcal I_2\) vanishes unless \(A\subset U,B\subset V\).
Indeed, if a freed input coordinate belongs to \(S\), it is
mean zero, absent from the target, and absent from the comparison
multiplier: integrating it gives zero.  The corresponding statement
holds for a freed target coordinate belonging to \(T\).
This uses the orientation in \eqref{eq:gr-signed}: the \(A\)
endpoint labels are from the target and the \(B\) endpoint labels
are from the input.

If \(A\subset U,B\subset V\), condition on all shared coordinates
outside \(A\cup B\), with product \(C\).
The input test is independent of the freed \(A\) labels and the
target test is independent of the freed \(B\) labels.
The two-product kernel comparison established above therefore
replaces the unsigned comparison by \(R(F_S,H_T)\), with error
\[
                    L^{-A_0}\|F_S\|_2\|H_T\|_2.
\]
Its frequency is \(\Theta C\); uniformity for every real frequency
is essential.  Cauchy--Schwarz over the conditioned coordinates
gives exactly the displayed unconditional norms.

For nonempty \(U,V\), the total unsigned subset sign is
\[
 \sum_{\substack{\varnothing\ne A\subset U\\
                 \varnothing\ne B\subset V}}
       (-1)^{|A|+|B|}
 =\left(\sum_{\varnothing\ne A\subset U}(-1)^{|A|}\right)
  \left(\sum_{\varnothing\ne B\subset V}(-1)^{|B|}\right)=1.
\]
The initial minus sign in \eqref{eq:gr-signed} cancels the raw
pairing.  If either missing set is empty, no comparison survives.
All exceptional raw component pairs together have the exact identity
\[
 R_{\rm exc}(F,H)=
 R(P_1F',H')+R((I-P_1)F',P_2H').
\]
In operator notation, if \(M_\Phi\) multiplies by
\(\Phi_{\Theta,\zeta}\), their operator on the symmetric space is
\[
 T_{\rm exc}=
 \Pi Q_{\rm tail}
   \big[M_\Phi P_1+P_2M_\Phi(I-P_1)\big]
                 Q_{\rm tail}\Pi .
\]
Consequently
\[
 \|T_{\rm exc}\|
       \le\theta_m^{s_1/2}+\theta_m^{s_2/2}.
\]
This identity sums the exceptional components before taking an
absolute value.  It is the reason their bound has no subset-count
or comparison-coefficient factor.

There are at most \(3^{P}\) admissible component/free-subset triples,
where \(P=m(s_1+s_2)\le mC_0\log L\): each probe is indicated,
missing but not freed, or freed.  Thus the full signed ideal
operator satisfies
\[
 \|T_{\rm signed}-T_{\rm exc}\|
           \le3^{P}L^{-A_0}.
\]
Choose \(m\) so that \(\theta_m<1\) and
\(-c_0\log\theta_m/6>E_{\rm id}+3\), and then choose
\(A_0>E_{\rm id}+mC_0\log3+10\).
The kernel precisions above and the exponent \(C_1\) in
\eqref{eq:gr-coeffcost} are chosen after these choices.
Only afterwards require \(J\ge m\) and
\(J+\ell\ge(m+\ell)(m+\ell-1)\); increasing this fixed \(J\)
to meet the physical root-moment theorem changes no kernel
parameter or coefficient exponent.
The resulting ideal norm is at most \(L^{-E_{\rm id}}\).
This proves the asserted ideal norm and coefficient conditions.
\end{proof}

\subsection{A large displacement in the physical graph}

\begin{lemma}[Large-target root moment and pairing]\label{lem:gr-large}
The root moment and endpoint pairing from
\cite[Theorem~3.5 and Corollary~3.11]{weightedgraphs} hold for any
nonzero integer \(k\) coprime to the bank primes, with the same
logarithmic accuracies.  In particular they hold for \(k=T\asymp x\)
after target-prime deletion.  The first endpoint is required to be
independent of its marks, exactly as in the imported pairing.
\end{lemma}

\begin{proof}
Fix the labels, patterns, and permutations of a closed walk of length
\(2R\), where \(R=\lceil L^{0.52}\rceil\).
Its position offsets are \(ks_i\).  Return to the initial integer
is equivalent to \(s_{2R}=0\), independent of the nonzero \(k\).
The root dependence consists of divisibility and damping conditions
at the offsets \(ks_i\).

The independent-residue replacement
\cite[Lemma~3.6]{weightedgraphs} is uniform in the numerical values
of these offsets.  Indeed its active modulus has logarithm
\(O_M(R(\log L)L^{2/5})=o(L)\).
The Bonferroni truncation has the same modulus and term count,
and each interval CRT error is \(O(1)\), independently of the
specified residues.  Once the deterministic damping from active
primes is factored out, the error is relative and can be summed with
absolute path coefficients using the imported absolute row bound.

In the resulting independent residue space, replace the root residue
\(z_p\) by \(k^{-1}z_p\bmod p\) for every bank prime \(p\).
Every fixed-path divisibility and damping expectation is then exactly
the expectation for \(k=1\).
The remaining signed coefficients are unchanged.
Thus the \(k=1\) root moment gives the identical bound for this \(k\).
Negative root intervals follow either from the same CRT count or
from \(n\mapsto-n,\ k\mapsto-k\).

It remains to avoid a large factor in the endpoint pairing.
Let \(D_{\max}\) bound the edge products, so that
\(\log D_{\max}=O_M((\log L)L^{2/5})\).
Partition the first-endpoint interval \(I\) into blocks of length
\(\lceil |k|D_{\max}\rceil\), as in the source pairing, and refine
each block by residue modulo \(|k|\).
The operator preserves these residues.
Each refined block has \(O(D_{\max}+1)\) integer positions.
For \(B=\mathcal A_k\mathcal A_k^*\), the matrix
\((\langle u_n,B^Ru_m\rangle)\) on one such block is positive
semidefinite.  Its largest eigenvalue is at most its trace.
Here \(u_n\) is the unnormalized vector equal to one on all marks at
\(n\).  Mark independence of the first endpoint \(f_n\) gives
\[
 \langle f_j,B^Rf_j\rangle
 \le \Big(\sum_{n\in I_j}|f_n|^2\Big)
             \sum_{n\in I_j}\langle u_n,B^Ru_n\rangle.
\]
Spectral H\"older and Cauchy--Schwarz give the source block estimate;
its only dimension factor is now
\begin{equation}
 \big(O(D_{\max})\sup|f_n|^2\big)^{1/(2R)}
 \le
 \exp\!\left(O_M\!\left(
 \frac{(\log L)L^{2/5}+\sqrt L}{L^{0.52}}\right)\right)=O(1).
 \label{eq:gr-block}
\end{equation}
Second-endpoint neighborhoods have bounded overlap within a residue
class and are disjoint between residue classes.
H\"older over blocks uses only the sum of root traces on the whole
original interval, bounded by the root theorem.
No CRT estimate on an individual residue class modulo \(k\) is used.
This proves the pairing uniformly in \(k\).
\end{proof}

For bounded invariant endpoint functions, the list moments needed
here do not create a mark-dependent logarithmic loss.  For a fixed
ordered distinct list with product \(P\), divisibility gives
\(O(X/P)\) integers in an interval of length \(O(X)\) contained in
\([-CX,CX]\).  For fixed mark length, \(P=x^{o(1)}\).
After summing the list weights,
\begin{equation}
 \sum_{\text{lists}}\frac1P\prod_gV_g^{-M}\le1.
 \label{eq:gr-list-moments}
\end{equation}
Thus the endpoint \(r\)-moments are \(O_r(X)\).
Condition \textup{(C1)} bounds every fixed-depth model without a
depth-dependent logarithmic exponent. Thus this observation is valid
before the final depth is chosen.

\subsection{Shared labels and the two comparison geometries}

The following consequence of the pairing applies to both geometries.
For a vector \(\mathbf t=(t_g)\) and a bounded function \(b\) on its
ordered distinct lists, put
\begin{equation}
 W_{\mathbf t}^b(n)=
 \prod_g\big(\1_{\omega_g(n)=t_g}V_g^{-t_g}\big)
       \sum_{\substack{\mathbf p\text{ a }\mathbf t\text{-list}\\
                        p_{g,j}\mid n}} b(\mathbf p).
 \label{eq:gr-listweight}
\end{equation}
Here \(\omega_g(n)\) counts distinct bank-prime divisors of \(|n|\).
For fixed \(t_g\le m+1\), these weights are bounded by a fixed
power of \(L\).  In particular \(W_{\mathbf1}^1=\Omega_{\mathcal P}\).

\begin{proposition}[Graph-to-arithmetic reduction]
\label{prop:gr-arithmetic-reduction}
Assume the imported graph inputs of Section~\ref{sec:inputs}.
Use the progression banks of this section, with primes dividing
\(k\) deleted, and take \(k=x\) for reflected sums or fixed
\(k=h\ne0\) for shifts. Let \(\alpha,\beta\) be functions of the
integer endpoint, of modulus at most one, invariant under repeated
multiplication by every bank prime whenever both arguments are in
the endpoint ranges. For sums use the negative first half-line and
positive second half-line, extending the first function by absolute
value; for shifts use positive endpoints. The pullback of the first
endpoint to physical states must be constant on its mark lists.

Take \(\ell=1\), \(M=J+1\), and a smooth \(\Psi\) supported
in a fixed compact subset of \((0,\infty)\). For every fixed
accuracy \(A>0\), Lemmas~\ref{lem:gr-ideal} and
\ref{lem:gr-large} give fixed \(m,C\) and a sufficiently large
fixed \(J\) for which
\begin{equation}
 C_0+\sum_\nu c_\nu C_\nu=O_A(L^{-A}),\qquad
 \#\{\nu\}+\sum_\nu|c_\nu|\le L^{C},                  \label{eq:gr-comparison}
\end{equation}
where
\[
 C_0=\sum_{w\ne0}\frac{\Psi(|w|/x)}{|w|}
       \alpha(w)\beta(w+k)
       W_{\mathbf1}^1(w)W_{\mathbf1}^1(w+k)
\]
and each comparison has the form
\begin{equation}
 \begin{split}
 C_\nu=\E_{D_1,D_2}a_\nu(D_1)b_\nu(D_2)
 \sum_{w\ne0}\frac{\Psi(|w|/(xD_1D_2))}{|w|}
  \alpha(w)\beta(w+kD_1D_2)\\
 {}\times W_{\mathbf t}^{u_\nu}(w)
                  W_{\mathbf t}^{v_\nu}(w+kD_1D_2).
 \end{split}                                                     \label{eq:gr-comparison-term}
\end{equation}
The following restrictions and normalizations are part of the conclusion.
\begin{enumerate}
\item The physical pairing is restricted to
\(\omega_g(n)=\omega_g(n')=M\) in every bank. Its state mass is
\(\prod_gV_g^{-M}\), its transition row normalization is
\(\prod_gV_g^{-t_g}\), and hence its joint endpoint normalization
in bank \(g\) is \(V_g^{-M-t_g}\). After shared labels are removed,
the exact counts are \(\omega_g(w)=\omega_g(w+kD_1D_2)=t_g\),
as imposed by \eqref{eq:gr-listweight}, where
\(t_g=1+\#\{\text{freed slots in }g\}\le m+1\).
In every small bank \(t_g=1\).
\item Each comparison frees nonempty probe sets in two disjoint
blocks of big banks. The products \(D_1,D_2\) use these respective
sets, with at most \(m\) draws per bank. All auxiliary draws are
independent with laws \(\mu_g\), including between the two products;
repetitions among auxiliary draws are allowed. No auxiliary-label
exclusion from endpoint lists remains in the displayed comparison.
\item Before separation, a signed pattern coefficient may depend
on the freed big-bank auxiliary labels and the unshared big-bank
endpoint labels specified in \eqref{eq:gr-signed}. It uses no small-bank
labels, shared labels, or distinguished last endpoint slots.
After separation, \(a_\nu\) depends only on \(D_1\),
\(b_\nu\) only on \(D_2\), and \(u_\nu,v_\nu\) only on
the ordered unshared big-bank labels at their respective endpoints.
Each of these four factors has modulus at most one. In particular,
both list coefficients ignore every small-bank label.
\item The parameters \(m,C\), the comparison family and its scalar
coefficient bounds are fixed before the final mark length \(J\)
and the final model depth. Increasing \(J\) supplies shared slots
without changing them. For endpoints satisfying \textup{(C1)},
no logarithmic loss exponent depends on the model depth; the
threshold in \(x\) may depend on that depth and on \(J\).
\end{enumerate}
\end{proposition}

\begin{proof}
The signed coefficients are those of Lemma~\ref{lem:gr-ideal};
their allowed label dependence and four-factor expansion were proved
in \eqref{eq:gr-signed}--\eqref{eq:gr-coeffcost}.
Each auxiliary factor depends only on its numerical product:
\eqref{eq:gr-kernel} uses only that product's log cell and character.
Thus the expectation over ordered draws pushes forward to the product
law \eqref{eq:gr-productlaw}, including its repetition multiplicities.
The first endpoint
is mark independent, as required by the imported pairing, while
all pattern-dependent mark factors remain in the operator.
A pattern has \(M-1\) dilation slots per bank, some shared
and some freed, and one distinguished unshared endpoint slot.
Write its edge product as \(D=D_RD_C\), with shared product \(D_R\)
and auxiliary product \(D_C\).
For a source of either sign in \(|n|\in[Z,2Z)\), an originally nonzero edge
weight satisfies
\(\Psi(|n|/(xD))\ne0\), hence \(D\ll Z/x\).
Since \(|k|\ll x\), this implies \(|n+kD|\ll Z\).
Insert a second-endpoint cutoff on \([-C_\Psi Z,C_\Psi Z]\)
before Mellin inversion.  It equals one on every contributing edge.
After inversion retain this same cutoff.  An edge originally having
zero \(\Psi\) still has zero full Mellin integral, so the insertion
does not change the identity.  No cutoff involving shared-label
products has been introduced.

Mellin inversion of \(\Psi\), with the harmonic weight \(1/|n|\),
now applies the pairing of Lemma~\ref{lem:gr-large} on each of
\(O(L)\) source dyads.  Its tails are bounded by the absolute
row/column estimate, its main part by the small ideal norm, and
its endpoint norms by \eqref{eq:gr-list-moments}.
Compressing the physical operator to exact-count vertices
\(\omega_g(n)=\omega_g(n')=M\) is legitimate:
the corresponding projection depends only on the integer, preserves
mark independence, and sets both damping factors equal to one.

For completeness, put \(n=D_Rw\), \(n'=D_Rw'\), so
\(w'=w+kD_C\), and put \(t_g=1+\#\{\text{freed slots in }g\}\).
Invariance and harmonic scaling give
\[
 \alpha(D_Rw)=\alpha(w),\quad \beta(D_Rw')=\beta(w'),\quad
 \frac{\Psi(|n|/(xD))}{|n|}
       =\frac1{D_R}\frac{\Psi(|w|/(xD_C))}{|w|}.
\]
The graph normalization at a shared prime combines with \(1/D_R\)
to give its probability law \(\mu_g\).
More explicitly, with \((D_R)_g\) its product in bank \(g\),
\begin{equation}
 \frac{V_g^{-M-t_g}}{(D_R)_g}
       =V_g^{-2t_g}
           \prod_{\text{shared slots in }g}\mu_g(p).   \label{eq:gr-sharednormal}
\end{equation}
For fixed reduced endpoints and fixed auxiliary labels,
the shared-label probability is
\begin{equation}
 \prod_g\1_{\omega_g(w)=\omega_g(w')=t_g}
              +O(L^{C_M}e^{-L^{1/8}}).               \label{eq:gr-sharedprob}
\end{equation}
Indeed there are only \(O_M(\log L)\) draws, each atom is
\(O(e^{-L^{1/8}})\), and \(ww'D_C\) has \(O(L)\) prime divisors.
If the reduced counts are \(t_g\), distinct draws avoiding those
divisors reconstruct the exact graph edge.  If a reduced count is
larger, a valid reconstruction must choose a shared prime already
dividing that endpoint, an event of the displayed small probability.
This proves \eqref{eq:gr-sharedprob}, also for negative \(w\).

Summing its absolute error costs only fixed logarithmic powers,
by bounded endpoints and the harmonic weight.
The restrictions excluding auxiliary labels from endpoint lists
can then be removed: for an auxiliary prime \(p\mid D_C\),
\(p\mid w\) if and only if \(p\mid w+kD_C\); these positions have
density \(O(1/p)\) in the containing interval.
Every auxiliary prime is at least \(\exp(L^{1/3})\).
An unshared prime common to the two endpoints divides \(kD_C\);
bank-prime deletion implies it divides \(D_C\), so the same estimate
covers it.  Thus these errors remain smaller than every logarithmic
power.  Finally expand the separated coefficients from
\eqref{eq:gr-signed}.  Formula \eqref{eq:gr-sharednormal} yields
exactly \eqref{eq:gr-listweight}, while \(D_C=D_1D_2\).
The raw pattern has \(D_C=1,\mathbf t=\mathbf1\).
This proves \eqref{eq:gr-comparison}--\eqref{eq:gr-comparison-term}
with their stated independence properties.
The kernel parameters and coefficient exponents were fixed before
\(J\) in Lemma~\ref{lem:gr-ideal}. Endpoint boundedness and
\eqref{eq:gr-list-moments} pay the list norms without a
depth-dependent logarithmic power. Thus the final model depth can
be chosen later to meet the one-variable accuracy, changing only
the threshold for boundedness in \textup{(C1)}.
\end{proof}

\section{Frequency tools for both branches}\label{sec:prime-estimates}

The shared comparison has left two independent products of large
prime labels. We first record its common Fourier normalization, then
the prime estimates used by the two frequency branches.

\subsection{The common Fourier reduction}

Dyadically localize the products in
\eqref{eq:gr-comparison-term} to \(D_i\asymp U_i\). Their ranges give
\[
 U_i\ge\exp(L^{1/3}),\qquad
 U_i\le\exp(O_m((\log L)L^{2/5})),\qquad Y=xU_1U_2.
\]
After inserting endpoint dyadic cutoffs before Mellin inversion, the
ordinary comparison has the form
\begin{equation}\label{eq:gr-fouriercomparison}
 \frac1Y\int_0^1\widehat F(\theta)\widehat G(-\theta)
                           m(k\theta)\,d\theta,\qquad
 m(\theta)=\E b_1(D_1)b_2(D_2)\ee{\theta D_1D_2}.
\end{equation}
The endpoint reflection for sums changes only signs and fixed character
factors. External Mellin transforms have fixed logarithmic costs.
The centered endpoint has the following common form.
Let \(\xi\) be either centered function in \textup{(C3)}, at a
fixed depth which will be chosen after the required precision.
For a comparison term, define
\begin{equation}
 F(n)=\psi(n/Y)n^{i\sigma}\xi(n)W_{\mathbf t}^b(n),
 \quad Y=x^{1+o(1)},\quad |\sigma|\le L^C,             \label{eq:gr-energy-F}
\end{equation}
where the support of \(\psi\) is a fixed compact positive interval.
All its logarithmic derivatives are bounded by fixed powers of \(L\).
The coefficient \(b\) ignores a selected small-bank label, whose count
is exactly one.  This holds for every comparison above.
Reflection changes only fixed character factors and additive signs,
so we state the estimate at a positive endpoint.


The other endpoint satisfies \(\sum_n|G(n)|^2\ll YL^C\), by
bounded invariant endpoints and the pointwise list bound. No
multiplicativity of the list coefficients is required.

By \eqref{eq:gr-productlaw} and
\eqref{eq:gr-integer-bilinear}, a rational approximation
\(|\theta-a/r|\le r^{-2}\) gives
\begin{equation}\label{eq:gr-multiplier}
 |m(\theta)|^2\ll L^{C_{\rm prod}}
 \left(r^{-1}+U_1^{-1}+\frac{\log(2r)}{U_2}
                    +\frac{r\log(2r)}{U_1U_2}\right).
\end{equation}
The two branches apply this bound at \(k=h\) and \(k=x\),
respectively; their resulting physical major windows differ.

\subsection{Sharp twisted prime sums}


The sharp estimate uses the zero-free input at conductors polynomial
in \(x\). The reflected branch needs that range; the fixed-shift
branch uses the smooth version below at polynomial-logarithmic
conductors and polynomial height.

\begin{lemma}[Prime factor estimate]\label{lem:gr-primes}
Assume the fixed all-conductor zero-free half-plane stated in
Section~\ref{sec:inputs}.
Fix \(\sigma_1\in(11/12,1)\) and \(a,B,B_1>0\).
For \(q\le x^B\), every character \(\chi\bmod q\), and
\(\exp(L^a)\le P\le x^{B_1}\), uniformly on subintervals
\(I\subset[P,2P]\),
\begin{equation}
 \sum_{p\in I}\chi(p)p^{it}
 =\delta_\chi\int_I\frac{u^{it}}{\log u}\,du
       +O((1+|t|)P^{\sigma_1}L^C).                  \label{eq:gr-primes}
\end{equation}
Here \(\delta_\chi=1\) for a principal character and zero otherwise,
and the estimate holds for every real \(t\). The exponent \(C\)
and implied constant may depend on \(a,B,B_1,\sigma_1\) and the
fixed bank modulus \(L_0\).
For progression banks the main vanishes unless \(\chi\) is induced
from a character modulo \(L_0\).
Removing modulus divisors or prime divisors of a nonzero target
\(|k|\le x^B\) preserves the error bound.
\end{lemma}

\begin{proof}
A fixed distance to the right of the zero-free boundary, the local
logarithmic-derivative formula and functional equation give
\begin{equation}
 \left|\frac{L'}{L}(s,\chi)+\frac{\delta_\chi}{s-1}\right|
                 \ll \log(q(|\Im s|+3)).
 \label{eq:gr-logderivative}
\end{equation}
Imprimitive Euler factors contribute \(O(\log q)\), since the real
part is fixed and positive.
Apply truncated Perron to
\(-L'/L(s,\chi)\), at
\(c=1+1/\log Z\), height \(Z^2\), and a half-integer endpoint
\(Z\in[P/2,4P]\).
The truncation error is \(O(\log^2 Z)\): near \(Z\), sum the bounds
\(O(Z\log Z/(Z^2|n-Z|))\); away from \(Z\), use absolute convergence
at \(c\).  Shift to \(\sigma_1\).
The principal pole gives \(\delta_\chi Z\), the new vertical integral
is \(O(Z^{\sigma_1}\log^2(qZ))\), and horizontal edges cost
\(O(Z^c\log(qZ)/(Z^2\log Z))\).
Prime powers cost \(O(Z^{1/2}\log^2 Z)\).
Replacing a real endpoint by the adjacent half-integer changes its
prime-weighted count by \(O(\log Z)\).  Therefore
\[
       \sum_{p\le Z}\chi(p)\log p
          =\delta_\chi Z+O(Z^{\sigma_1}L^C).
\]
Here \(Z\asymp P\le x^{B_1}\) and \(q\le x^B\) give
\(\log(qZ)\ll_{B,B_1}L\), justifying the displayed logarithmic cost.
Partial summation against \(u^{it}/\log u\) proves
\eqref{eq:gr-primes} with its linear \(1+|t|\) growth.

Expand the progression condition into the finitely many characters
modulo \(L_0\).  A constituent is principal exactly when the original
character is induced from \(L_0\).
There are at most \(O(\log(q(2+|k|)))=O(L)\) deleted prime divisors;
their absolute unweighted contribution is absorbed by the error.
\end{proof}

Every prime scale used below lies in one of the two bank ranges,
so \(P\le\exp(L^{2/5})\le x\) for large \(x\). Thus \(B_1=1\)
suffices for both comparison geometries; in the reflected application
the conductors are at most \(xL^{O(1)}\), and in the fixed-shift
application they are polynomial logarithmic.

In particular, for a nonexceptional character the prime factor saves
\(P^{\sigma_1-1}\) against its trivial length.  The all-height form
above is convenient: smooth product-cutoff Mellin transforms have
integrable \((1+|t|)\)-weighted absolute value bounded by a fixed
power of \(L\), so there is no boundary-smoothing loss from a sharp bank.


\subsection{Smooth prime sums at polynomial height}


The high-frequency argument uses a smooth version of
Lemma~\ref{lem:gr-primes}. Its height dependence is essential:
the linear \(1+|t|\) loss in the sharp-interval estimate is not
adequate when \(|t|\) can be of order \(x\).
Fix \(\sigma_1\in(11/12,1)\). For \(w\) smooth and supported on
\([1/2,2]\), with each required derivative bounded by a fixed
power of \(L\), the same zero-free input gives
\begin{equation}\label{eq:fixed-smooth-primes}
 \sum_p\chi(p)w(p/P)p^{it}
  =\delta_\chi\int_0^\infty\frac{w(u/P)u^{it}}{\log u}\,du
                       +O(P^{\sigma_1}L^{C_w})
\end{equation}
uniformly for \(q\le L^B\), characters \(\chi\bmod q\),
\(|t|\le x^C\), and \(P\ge\exp(cL^{1/8})\).
Here \(C_w\) depends on the fixed derivative bounds and on \(B,C\).
This is the smooth prime estimate used in
\cite[Lemma 5.1 and Section 6.4]{walsh}.

To obtain it directly from the present input, insert the von
Mangoldt weight and the smooth factor \(w(u/P)/\log u\), apply
Mellin inversion, and move the line from \(\Re s>1\) to
\(\Re s=\sigma_1\). Formula \eqref{eq:gr-logderivative} bounds
the regularized logarithmic derivative on this line by
\(O(\log(q(|\Im s|+3)))\). The Mellin transform is translated
in height by \(t\); its decay depends on the derivatives of
\(w\), rather than on \(|t|\). Its weighted integral is consequently
\(O(L^{C_w})\) even for \(|t|\le x^C\). The principal pole
gives the integral in \eqref{eq:fixed-smooth-primes}; prime powers
cost \(O(P^{1/2}L^{C_w})\), absorbed by the displayed error.
Truncation followed by dominated convergence justifies the contour
move. This uses the same fixed zero-free half-plane, without a
second prime-distribution theorem.

Impose simultaneously the bank class \(c\pmod{L_0}\) and a
prime class modulo \(q\). Incompatible classes contribute zero;
otherwise character orthogonality at \(\lcm(q,L_0)\) applies.
For large \(x\), every prime in the small bank exceeds
\(qL_0|h|\), so nonreduced prime classes and primes dividing
\(h\) do not occur. Integration by parts in the principal
integral, for every fixed \(j\), shows that the normalized
prime sum \(P^{-1}\sum_p\cdots\) is bounded by
\begin{equation}\label{eq:fixed-high-prime}
        L^{C_{w,j}}(1+|t|)^{-j}
                       +L^{C_w}\exp(-cL^{1/8}).
\end{equation}
All character and congruence costs here are polynomial logarithmic.


\section{Fixed shifts and wide major arcs}\label{sec:fixed-shift-comparison}

For fixed \(k=h\ne0\), the shared-label comparison
\eqref{eq:gr-comparison}--\eqref{eq:gr-comparison-term} remains valid.
The frequency estimate has a different scale. After localizing
\(D_i\asymp U_i\), put
\[
 U_*=U_1U_2,\qquad Y=xU_*,\qquad W=U_*L^{-B_0}.
\]
The multiplier is now \(m(h\theta)\). The integer bilinear bound
\eqref{eq:gr-multiplier} confines its significant part to
\(O_h(L^{2B_0})\) arcs about rationals \(a/q\), with
\[
 q\le |h|L^{B_0},\qquad
       |\theta-a/q|\le C_h/W.
\]
In particular, their physical width is \(L^{B_0}/U_*\), rather than
\(L^{B_0}/Y\). The target-grid proposition does not estimate these
arcs: on them \(|\beta|Y\) can be of order \(xL^{B_0}\).
We extract the already established Dirichlet-frequency argument of
\cite[Sections 6.2--6.4]{walsh}, with the arithmetic hypotheses
\textup{(C1)--(C3)} in place of its particular divisor model.

\begin{proposition}[Fixed-shift comparison energy]
\label{prop:fixed-comparison-energy}
Let \(F\) have the form \eqref{eq:gr-energy-F}, with a selected
small bank whose label does not enter \(b\), and with its exact
prime count equal to one. Suppose
\[
 U_i\ge\exp(L^{1/3}),\qquad
 U_i\le\exp(O((\log L)L^{2/5})),\qquad Y=xU_1U_2.
\]
For every fixed \(B_0,D>0\), conditions \textup{(C1)--(C3)} and
the zero-free input imply, at sufficiently deep fixed model depth,
\begin{equation}\label{eq:fixed-wide-energy}
 \int_{|\beta|\le C_h/W}
      |\widehat F(a/q+\beta)|^2\,d\beta
                            \ll_{D,h}YL^{-D}
\end{equation}
uniformly for \(q\le |h|L^{B_0}\), where
\(\widehat F(\theta)=\sum_nF(n)\ee{\theta n}\).
Consequently every fixed-shift comparison in
\eqref{eq:gr-comparison-term} has an arbitrary prescribed
logarithmic saving. No binary cancellation estimate is assumed.
\end{proposition}

\subsection{Conversion to Dirichlet frequencies}

We use the following existing two mean-square inequalities, in their
normalizations from \cite[Lemma 6.3 and (6.14)]{walsh}. If \(a_n\)
is supported on \(n\asymp Y\), then
\begin{equation}\label{eq:fixed-additive-mellin}
 \int_{|\beta|\le C/W}\left|\sum_na_n\ee{\beta n}\right|^2d\beta
 \ll_C\frac1Y\int_{\mathbb R}(1+W|t|/Y)^{-2}
                         \left|\sum_na_nn^{it}\right|^2dt.
\end{equation}
Here is the averaging step in that inequality, following
\cite[Lemmas 2.1--2.2]{soundararajan}. Put
\(\eta=W/(4C)\), \(A(v)=\sum_{n\le v}a_n\),
\(P(t)=\sum_na_nn^{it}\), and
\(J_\eta=\int_0^\infty|A(v+\eta)-A(v)|^2dv\).
The interval kernel has Fourier transform of modulus at least
\(c\eta\) on \(|\beta|\le C/W\), so Plancherel bounds the
left side by \(O(\eta^{-2}J_\eta)\).
For \(2\eta\le\nu\le3\eta\), split the increment at \(v+\nu\),
square, and average in \(\nu\). Changing \(v\) to \(v+\eta\)
in the second increment gives, if the coefficients are supported
on \([aY,bY]\),
\[
 J_\eta\ll\eta^{-1}\int_\eta^{3\eta}
       \int_{aY/2}^{2bY}|A(v+\nu)-A(v)|^2dv\,d\nu.
\]
Substitute \(\nu=\Delta v\). Since \(v\asymp Y\), this is at most
\[
 \frac{CY^2}{\eta}
 \int_{c_1\eta/Y}^{c_2\eta/Y}\int_0^\infty
       |A(v(1+\Delta))-A(v)|^2\frac{dv}{v}\,d\Delta.
\]
Logarithmic Plancherel evaluates the inner integral exactly as
\[
 \frac1{2\pi}\int_{\mathbb R}|P(t)|^2
          \left|\frac{(1+\Delta)^{it}-1}{t}\right|^2dt.
\]
The kernel is \(O(\min\{\eta^2/Y^2,|t|^{-2}\})\), and
the \(\Delta\)-interval has length \(O(\eta/Y)\). Thus
\(J_\eta\ll Y\int\min\{\eta^2/Y^2,|t|^{-2}\}|P(t)|^2dt\).
Multiplying by \(\eta^{-2}\) proves
\eqref{eq:fixed-additive-mellin}; no identification of fixed additive
windows with fixed multiplicative windows is used. The other inequality,
the ordinary Dirichlet-polynomial mean-value theorem, is
\begin{equation}\label{eq:fixed-dirichlet-mean}
 \int_I\left|\sum_{n\le CR}a_nn^{it}\right|^2dt
              \ll_C (|I|+R)\sum_n|a_n|^2,
\end{equation}
uniformly in the location of the interval \(I\). Neither inequality
is a new distributional input.

Apply \eqref{eq:fixed-additive-mellin} to
\(a_n=F(n)\ee{an/q}\) and set
\[
 \mathcal T(t)=Y^{-1}\sum_nF(n)\ee{an/q}n^{it},\qquad
 T_0=(Y/W)L^{B_1}=xL^{B_0+B_1}.
\]
The pointwise list bound and \textup{(C1)} give
\(\sum_n|F(n)|^2\ll YL^C\). Thus on a dyadic interval
\(T\le |t|\le2T\), \eqref{eq:fixed-dirichlet-mean} gives
\[
 \int_{T\le|t|\le2T}|\mathcal T(t)|^2dt
                                  \ll L^C(1+T/Y).
\]
Summing this with the weight in
\eqref{eq:fixed-additive-mellin}, for \(|t|>T_0\), gives
\begin{equation}\label{eq:fixed-frequency-tail}
 \int_{|t|>T_0}(1+W|t|/Y)^{-2}|\mathcal T(t)|^2dt
       \ll L^C\bigl(L^{-2B_1}+W^{-1}L^{-B_1}\bigr).
\end{equation}
The exponent \(B_1\) is chosen after the required saving, so this
tail is arbitrarily small. It remains to prove
\begin{equation}\label{eq:fixed-dirichlet-energy}
                \int_{|t|\le T_0}|\mathcal T(t)|^2dt\ll_D L^{-D}.
\end{equation}

\subsection{Low frequencies and recordwise centering}

Choose a fixed exponent \(B_2\), to be specified after the smooth
prime estimates, and put \(T_c=L^{B_2}\). In the low range write
\(n=rv\), where \(r=n_{\mathcal P_y}\) is the complete smooth record and
\(v\) is rough. Every list weight is a function of \(r\) alone
and is bounded by \(L^C\). The terms with \(r>E\) have population
\(O(YL^C\exp(-cL^{1/10}))\), by Rankin's inequality with parameter
\(c/\log y\), as in \eqref{eq:record-tail}; their contribution to
\(\mathcal T(t)\)
is negligible uniformly in \(t\). For each remaining record,
\(\ee{arv/q}\) is a bounded periodic test of period at most \(q\).
The phase \((rv)^{i(t+\sigma)}\psi(rv/Y)\) has fixed logarithmic
derivatives bounded by fixed powers of \(L\), throughout
\(|t|\le T_c\). Condition \textup{(C3)} therefore gives any
prescribed saving for that record. Summing uses only
\[
          \sum_{r\ \mathcal P_y\text{-smooth}}\frac1r
             =\prod_{p\le y}(1-1/p)^{-1}\ll L^C.
\]
Choose the precision in \textup{(C3)} after \(B_2\) and all list
costs. This gives
\begin{equation}\label{eq:fixed-low}
             \int_{|t|\le T_c}|\mathcal T(t)|^2dt\ll_D L^{-D}.
\end{equation}
The arithmetic progression test is on \(v\), so this argument
does not require \(r\) or the endpoint to be a unit at \(L_0\).

\subsection{High frequencies and a long cofactor}

For this part remove selected-bank prime squares, whose population
is \(O(Y/P_{\min})\). Applied to their bounded coefficients,
\eqref{eq:fixed-dirichlet-mean} bounds their normalized integrated
energy on \([-T_0,T_0]\) by
\(L^C(1+T_0/Y)/P_{\min}\), hence by less than every logarithmic
power. Do not impose the cutoff \(n_{\mathcal P_y}\le E\) here: it would
couple the selected prime to its cofactor.

The remaining terms have a unique selected-bank prime and can be
written \(n=pm\). Exact invariance \textup{(C2)} gives
\(\xi(pm)=\xi(m)\). The distinguished list coefficient ignores
this prime, so after removing its factor \(V_g^{-1}\) the
remaining coefficient \(\gamma(m)\) is independent of \(p\) and
bounded by \(L^C\). It contains the condition that \(m\) has no
selected-bank prime, together with the other list weights.

Partition the selected bank into smooth dyadic pieces
\(p\asymp P\), including its two sharp outer endpoints. Replace
each sharp endpoint by a smooth cutoff across a strip of length
\(\varepsilon P\), where \(\varepsilon=L^{-A_{\rm sm}}\).
The normalized sum of prime labels in such a strip is at most
\(O(\varepsilon+P^{-1})\), simply by counting integers. The
Dirichlet mean-value bound for the cofactor below and Minkowski
show that this replacement changes the normalized integrated
energy by at most
\[
                  L^C(\varepsilon+P_{\min}^{-1})^2.
\]
There are only polynomially many dyadic pieces in \(L\).
Choose \(A_{\rm sm}\) large after the desired saving and these
costs. The resulting smooth prime cutoffs have derivative bounds
which are fixed powers of \(L\).

Put \(R=Y/P\), and insert a smooth cofactor cutoff
\(g(m/R)\) equal to one wherever the product cutoff can contribute.
Separate the additive phase by the residue classes of \(p,m\)
modulo \(q\), and Mellin-invert the product cutoff. Apart from
the errors just bounded, \(\mathcal T(t)\) is a sum of at most
polynomially many terms of the form
\begin{equation}\label{eq:fixed-prime-cofactor}
 \frac{1}{2\pi V_g}\sum_{r,r'\bmod q}\ee{arr'/q}
       \int_{\mathbb R}\widetilde\psi(\nu)Y^{-i\nu}
          P_r(t+\sigma+\nu)R_{r'}(t+\sigma+\nu)\,d\nu,
\end{equation}
where harmless Mellin sign choices are fixed by this formula and
\[
 \begin{split}
 P_r(s)&=P^{-1}\sum_{\substack{p\equiv r\ (q)\\p\equiv c\ (L_0)}}
                    w(p/P)p^{is},\\
 R_{r'}(s)&=R^{-1}\sum_{m\equiv r'\ (q)}
                      g(m/R)\gamma(m)m^{is}.
 \end{split}
\]
There is no unit restriction on \(r'\). For any interval of
length at most \(4T_0\), \eqref{eq:fixed-dirichlet-mean} gives,
uniformly in its location,
\begin{equation}\label{eq:fixed-long-cofactor}
       \int_I|R_{r'}(t)|^2dt\ll L^C(1+T_0/R)\ll L^C.
\end{equation}
The final inequality uses the decisive scale separation
\[
 \frac{T_0}{R}\ll\frac{P L^{B_0+B_1}}{U_1U_2}
   \le\exp\bigl(CL^{1/4}-2L^{1/3}\bigr)L^{B_0+B_1}=o(1).
\]
It comes from the two independent big-prime products already
present in the graph comparison, rather than a new energy bound.

The Mellin transform in \eqref{eq:fixed-prime-cofactor} has
Schwartz decay with polynomial-logarithmic seminorms. Choose a
fixed exponent \(C_\nu\) large enough that its tail
\(|\nu|>L^{C_\nu}\), in weighted \(L^1\), is smaller than the
required logarithmic power. Uniformity in the location of \(I\)
in \eqref{eq:fixed-long-cofactor}, and the trivial bound
\(|P_r|\ll1\), justify removing this tail. Now choose \(B_2\)
larger than \(C_\nu\), the bound for \(|\sigma|\), and the
smooth-cutoff derivative costs. Throughout
\(T_c<|t|\le T_0\), the argument
\(s=t+\sigma+\nu\) then has \(|s|\gg L^{B_2}\).
In \eqref{eq:fixed-high-prime} fix, for example, the
integration-by-parts order \(j=2\) first. After choosing the
requested exponent \(A\), enlarge \(B_2\) beyond the resulting
fixed derivative costs. This gives
\[
           \sup_{T_c<|t|\le T_0,\ |\nu|\le L^{C_\nu}}
                    |P_r(t+\sigma+\nu)|\ll_A L^{-A}
\]
for any requested \(A\). The smooth prime estimate applies
because \(T_0=xL^{B_0+B_1}\le x^2\) for large \(x\).
Choose this \(A\) after the fixed partition and Mellin costs.
Combine the estimate with \eqref{eq:fixed-long-cofactor}, then apply
Minkowski to the Mellin integral and sum the dyadic and residue
pieces. This proves
\[
              \int_{T_c<|t|\le T_0}|\mathcal T(t)|^2dt\ll_D L^{-D}.
\]
Together with \eqref{eq:fixed-low} this is
\eqref{eq:fixed-dirichlet-energy}; together with
\eqref{eq:fixed-frequency-tail} and
\eqref{eq:fixed-additive-mellin} it proves
\eqref{eq:fixed-wide-energy}.

The parameter order is explicit. First fix the graph comparison,
its list bounds, and its external cutoff costs. Then fix the major
arcs and the required energy saving, the tail cutoff \(B_1\),
the bank smoothing precision, and the Mellin truncation. Choose
the high-frequency boundary \(B_2\) after those choices. Finally
choose the model depth through \textup{(C3)} at the required
low-frequency precision. No exponent in the graph or the endpoint
moment bound depends on that final depth.

On the arcs, Cauchy--Schwarz, \eqref{eq:fixed-wide-energy}, and
Parseval for the other endpoint bound the normalized Fourier
comparison by every requested inverse power of \(L\), after
paying the polynomially many arcs. Off the arcs,
\eqref{eq:gr-multiplier} and the two endpoint norms do so.
The normalization is \(1/Y\), exactly as in
\eqref{eq:gr-fouriercomparison}, with \(k=h\).
The harmonic conversion \(\Psi(t)=t\Phi(t)\) therefore gives
the same centered banked replacement as for reflected sums.
This completes the proof of
Proposition~\ref{prop:fixed-comparison-energy} and supplies the
fixed-shift branch used in Section~\ref{sec:assembly}.

\subsection{Classical estimates and the height range}\label{sec:fixed-classical-boundary}

The polynomial-logarithmic modulus here is an advantage over the
reflected target grid, but the height is a separate constraint.
Write \(Z=\exp(L^{1/16})\). The classical Dirichlet zero-free
region, together with Siegel's theorem in its ineffective all-character
form, gives \eqref{eq:fixed-smooth-primes} with error
\[
          O\!\left(P L^{C_w}e^{-cL^{1/16}}\right)
                      \qquad(|t|\le 2Z, q\le L^B).
\]
For exact classical statements see \cite[Proposition 51,
Theorem 62(ii) and Exercise 45(i)--(ii)]{taonotes}.
Indeed, a rectangle of height \(3Z^2\) has a zero-free strip of
width \(2\delta\), with \(\delta=cL^{-1/16}\).
For a possible exceptional character, choose
\(\varepsilon<1/(32(B+1))\) in Siegel's theorem; its zero then
has distance \(\gg L^{-1/32}\) from one, larger than the strip
width for sufficiently large \(x\). On the line
\(\Re s=1-\delta\), the regularized logarithmic derivative
has polynomial-logarithmic size. The same translated smooth Mellin
contour as above, truncated at imaginary part \(Z^2\), therefore
saves \(P^{-\delta}\le e^{-cL^{1/16}}\); fixed sufficiently high
derivatives control the horizontal and truncation terms. Prime powers
are smaller. Constants and the threshold in this classical estimate
may be ineffective.

The factor two absorbs the retained polynomial-logarithmic Mellin
translations when the outer frequency has absolute value at most \(Z\).
This covers the high-prime argument only up to \(Z\). Its required
upper height is \(T_0=xL^{B_0+B_1}\). At that endpoint the
classical strip has width of order \(L^{-1}\), whereas
\(\log P\le CL^{1/4}\); the contour saving
\(\exp(-c\log P/L)\) does not tend to zero. Thus a direct
Siegel--Walfisz substitution does not prove the remaining weighted
prime--cofactor energy on \(Z<|t|\le T_0\).
Walsh's classical-recovery remark \cite[p.~2]{walsh} motivates a
different high-frequency argument, but the proof here retains the
zero-free input on that range. Such a replacement for fixed shifts
would not remove the target-sized conductor issue for reflected sums.

\section{Reflected targets and exact rational grids}\label{sec:reflected-comparison}

This branch treats a displacement equal to the target. Unlike the
fixed-shift windows in Section~\ref{sec:fixed-shift-comparison}, its
rational denominators contain that target. We prove the residue-energy
estimate with those denominators retained.

\subsection{The exact target-grid energy}


Let \(\mathcal P_y=\{p:p\le y\}\), \(y=\exp(L^{2/5})\), and
\(E_0=\exp(\sqrt L)\), as in Section~\ref{sec:comparison}.
The endpoint \(F\) is the centered function in \eqref{eq:gr-energy-F}.

\begin{proposition}[Target-grid residue energy]\label{prop:gr-energy}
Suppose that
\begin{equation}
       q\le xL^B,\qquad
       \frac{Y}{qP_{\max}}\gg_A L^A
                       \quad\text{for every fixed }A,
 \label{eq:gr-energy-ratio}
\end{equation}
where \(P_{\max}\le\exp(CL^{1/4})\) is the selected small-bank scale.
For every prescribed \(D>0\), a fixed sufficiently deep density model
gives
\begin{equation}
 E_q(F):=\sum_{a\bmod q}\left|
           \sum_{n\equiv a\ (q)}F(n)\right|^2
                      \ll_D\frac{Y^2}{q}L^{-D}.       \label{eq:gr-energy}
\end{equation}
The estimate is uniform in all valuations of \(q\).
It also holds with \(F(n)\) multiplied by \(\ee{\beta n}\)
when \(|\beta|Y\) is bounded by a fixed power of \(L\).
\end{proposition}

\begin{proof}
Refine to \(Q=\lcm(q,L_0)\).
Cauchy--Schwarz on the refined classes gives
\(E_q(F)\le(Q/q)E_Q(F)\le L_0E_Q(F)\).
For every gcd profile \(d\mid Q\), put \(r=Q/d\).
Multiplicative orthogonality gives the exact identity
\begin{equation}
 E_Q(F)=\sum_{d\mid Q}\frac1{\varphi(r)}
       \sum_{\chi\bmod r}\left|\sum_m F(dm)\chi(m)\right|^2.
 \label{eq:gr-gcdidentity}
\end{equation}
The character zeros impose \((m,r)=1\), so this identity retains
prime powers, including primes dividing \(L_0\). We do not assume
that the endpoint is prime to \(L_0\). Call a character exceptional
when its primitive inducing character has conductor dividing \(L_0\).
There are only a fixed number of such characters for each \(r\).
In particular \(L_0\mid r\) is not required.

We first identify the negligible sets in the norm actually needed.
For a set \(\mathcal B\) and \(|G|\le L^C\), Cauchy--Schwarz in each
residue class gives
\begin{equation}
 E_Q(G\1_{\mathcal B})
 \le(1+CY/Q)\sum_{n\asymp Y}|G(n)|^2\1_{\mathcal B}(n).
 \label{eq:gr-tail-energy}
\end{equation}
The factor is \(O(Y/Q)\) here.
The following populations are exponentially small relative to \(Y\),
up to fixed logarithmic factors:
\begin{equation}
 \begin{array}{ll}
 n_{\mathcal P_y}>E_0:& O(YL^C e^{-cL^{1/10}}),\\
 p\mid n,\ p\mid Q,\ p>y:& O(YL/y),\\
 p\mid n,\ p\mid Q,\ p\text{ in the selected bank}:&
                  O(YL/P_{\min}),\\
 p^2\mid n,\ p\text{ in the selected bank}:&
                  O(Y/P_{\min}).
 \end{array}                                                       \label{eq:gr-tailsets}
\end{equation}
For the first bound, take \(a=c/\log y\). Counting multiples of
each possible smooth record gives
\[
 \#\{n\le CY:n_{\mathcal P_y}>E_0\}
 \le CY E_0^{-a}\prod_{p\le y}(1-p^{a-1})^{-1}
                 \ll YL^C e^{-cL^{1/10}}.
\]
This estimate uses no property of the indicators. The other bounds
follow by counting multiples,
using \(P_{\min}\ge\exp(cL^{1/8})\).
Thus all these sets are negligible in \eqref{eq:gr-tail-energy}.

For nonexceptional characters remove only selected primes dividing
\(Q\) and selected-prime squares.  Do not yet impose \(n_{\mathcal P_y}\le E_0\),
as that would couple the extracted prime to its cofactor.
The unique selected small-bank prime gives \(n=dpm\).
On a dyadic prime interval \(p\asymp P\), exact dilation invariance
and the independence of \(b\) from this slot make its remaining
coefficient \(\gamma_d(m)\) independent of \(p\), with
\(|\gamma_d(m)|\le L^C\).
It includes the condition that \(m\) has no other selected-bank prime,
and every other bank weight.  Its containing length is
\(M\asymp Y/(dP)\).
Mellin inversion of the product cutoff separates the character sum as
an integral of
\begin{equation}
 A_\chi(t)B_{d,\chi}(t),\qquad
 A_\chi(t)=\sum_{p\asymp P,\ p\in\mathcal P_g}\chi(p)p^{it},
 \quad
 B_{d,\chi}(t)=\sum_{m\asymp M}\gamma_d(m)\chi(m)m^{it}.
 \label{eq:gr-primefactorization}
\end{equation}
The omitted Mellin scaling has modulus one.
The transform has polynomial-logarithmic weighted \(L^1\) norm,
including weight \(1+|t|\).
Condition \eqref{eq:gr-energy-ratio} gives
\[
                       M/r\asymp Y/(QP)\gg1,
\]
independently of \(d\).
Elementary character orthogonality therefore yields
\begin{equation}
 \frac1{\varphi(r)}\sum_{\chi\bmod r}|B_{d,\chi}(t)|^2
 \le(1+CM/r)\sum_{m\asymp M}|\gamma_d(m)|^2
                  \ll \frac{M^2}{r}L^C.             \label{eq:gr-cofactorenergy}
\end{equation}
Apply Lemma~\ref{lem:gr-primes} to \(A_\chi\) and Minkowski to the
integral in \eqref{eq:gr-primefactorization}.
Its nonexceptional energy at profile \(d\) is
\begin{equation}
 \ll \frac{Y^2}{Qd}L^C\exp(-cL^{1/8}).               \label{eq:gr-nonexceptional}
\end{equation}
The prime dyadic decomposition costs only a fixed power of \(L\).
Summing \(d\) is harmless because
\(\sum_{d\mid Q}1/d\le Q/\varphi(Q)\ll\log\log(3Q)\).
There is no number-of-characters or unweighted divisor-count loss.

For exceptional characters use the further restriction
\(n_{\mathcal P_y}=A\le E_0\), and remove \(Q\)-prime divisors above \(y\).
The errors were already controlled in energy.
Now the gcd profile is determined by the smooth record:
\(\gcd(A,Q)=d\), \(n=Av\), with \(v\) \(\mathcal P_y\)-rough.
If \(\chi^*\) is the primitive character inducing an exceptional
\(\chi\bmod r\), then on the retained profile
\[
             \chi(n/d)=\chi^*(A/d)\chi^*(v).
\]
Indeed \(\operatorname{cond}\chi^*\mid r\) and
\((A/d,r)=1\). One must use \(\chi^*\) here: a literal character
modulo \(L_0\) could vanish on a retained bad-prime prefix.
Every other list weight depends only on \(A\), is bounded by \(L^C\),
and the remaining periodic test in \(v\) has fixed period \(L_0\).
The one-variable hypothesis must center this full periodic test;
matching only a scalar principal mean would not suffice.

There is still the condition that \(v\) avoid primes \(p\mid Q\)
above \(y\).  Remove it before applying the fixed-period centering
lemma, extending the same bounded periodic test to those new points.
At each retained record, multiples of such \(p\) number at most
\(CY/(Ap)\) on the containing interval \(0<v\le CY/A\).
Consequently the change in the profile sum is at most
\begin{equation}
 L^CY\sum_{\substack{A\ \mathcal P_y\text{-smooth}\\d\mid A}}\frac1A
                \sum_{\substack{p\mid Q\\p>y}}\frac1p
 \le \frac{Y}{d}L^CV_y^{-1}\frac{\log Q}{y\log y}
 \ll_D (Y/d)L^{-D}.                                  \label{eq:gr-avoidance}
\end{equation}
We used the exact harmonic identity
\begin{equation}
 \sum_{\substack{A\ \mathcal P_y\text{-smooth}\\d\mid A}}\frac1A
                       =\frac1{dV_y}                 \label{eq:gr-record-harmonic}
\end{equation}
when \(d\) is \(\mathcal P_y\)-smooth; otherwise this restricted profile is empty.

After this extension, hypothesis \textup{(C3)} applies record by
record at length \(Y/A\).
It permits arbitrary bounded fixed-period tests, arbitrary choices
of retained records, and the stated smooth phases.
Choose its logarithmic accuracy after the list and cutoff costs.
Using \eqref{eq:gr-record-harmonic} gives, for every prescribed \(D_0\),
\begin{equation}
            \left|\sum_mF_{\rm good}(dm)\chi(m)\right|
                         \ll (Y/d)L^{-D_0}.          \label{eq:gr-exceptional}
\end{equation}
Restoring the avoidance condition costs \eqref{eq:gr-avoidance}.
In \eqref{eq:gr-gcdidentity} the total exceptional contribution is at most
\begin{equation}
 \frac{Y^2}{Q}L^{-2D_0}
           \sum_{d\mid Q}\frac1d\,\frac{Q/d}{\varphi(Q/d)}
                     \ll \frac{Y^2}{Q}L^{-2D_0+C}.   \label{eq:gr-exceptionalenergy}
\end{equation}
Choose \(D_0\) large, add the negligible sets, and use the refinement
inequality to prove \eqref{eq:gr-energy}.
Multiplying the cutoff by \(\ee{\beta Yu}\) changes only its fixed
logarithmic derivative bounds when \(|\beta|Y\le L^B\), proving the
last assertion.
\end{proof}

This proof uses only boundedness, exact prime invariance, and the
recordwise periodic means in \textup{(C1)--(C3)}. In particular no
subgroup support, pole order, positivity of an individual character
branch, or unit-support condition at the fixed modulus enters the
energy estimate.

\subsection{From target-grid energy to comparison cancellation}

Specialize the common Fourier reduction of
Section~\ref{sec:prime-estimates} to the reflected family \(k=T=x\).
Use the product scales \(U_i\) and \(Y=xU_1U_2\) from that reduction.
Choose an outer cutoff \(B_{\rm out}\), after the finite comparison
family has been fixed.
Dirichlet approximation with denominator ceiling
\(U_1U_2/L^{B_{\rm out}-1}\) makes the multiplier smaller than any
requested inverse logarithmic power unless \(T\theta\) lies within
\(L^{B_{\rm out}}/(U_1U_2)\) of a rational of denominator
\(r\le L^{B_{\rm out}}\).
The corresponding physical frequencies are covered by grids
of denominator \(q=Tr\), with widths
\(O(L^{B_{\rm out}}/(TU_1U_2))=O(L^{B_{\rm out}}/Y)\).
Moreover
\[
 \frac{Y}{qP_{\max}}
 \gg \exp(cL^{1/3}-CL^{1/4})L^{-B_{\rm out}},
\]
so \eqref{eq:gr-energy-ratio} holds uniformly in every gcd profile.

Discrete additive Parseval gives the exact normalization
\begin{equation}
 \sum_{j\bmod q}|\widehat F(j/q+\beta)|^2
                         =qE_q(F\ee{\beta n}).        \label{eq:gr-gridparseval}
\end{equation}
Integrating over width \(L^{B_{\rm out}}/Y\) and applying
Proposition~\ref{prop:gr-energy} gives \(YL^{-D}\) for every desired
\(D\), after demanding a larger residue-energy saving.
The polynomial-logarithmic family of \(r\)'s is paid at the same step.
In particular the \(T\) grid points do not create a factor \(T\):
they are already included in \eqref{eq:gr-gridparseval}.
Cauchy--Schwarz and Parseval on the other endpoint bound
\eqref{eq:gr-fouriercomparison} on these grids by an arbitrary inverse
power of \(L\).  Off the grids, use
\eqref{eq:gr-multiplier} and the endpoint norms.
This proves that every comparison in \eqref{eq:gr-comparison} is
negligible to the requested logarithmic accuracy.

Finally \(\Psi(t)=t\Phi(t)\) makes
\[
               \frac{\Psi(|w|/x)}{|w|}
                       =\frac{\Phi(|w|/x)}x.
\]
Thus the harmonic raw correlation is \(1/x\) times the ordinary
banked correlation on the reflected line.
The comparison at scale \(Y\) has exactly the corresponding factor
\(1/Y\).  Paying the finite comparison coefficients therefore proves
the centered banked replacement required in the assembly.

Fixed bad-supported coefficients in \(Aa+Bb=T\) cause no new
analytic condition.  Embed, for example,
\(\widetilde F(n)=\1_{A\mid n}F(n/A)\).
Its exact residue energy is
\[
                  E_q(\widetilde F)=E_{q/(q,A)}(F),
\]
and its core length is divided by the fixed \(A\).
All energy ratios change only by fixed factors.
Bank primes are units at \(A,B\), so dilation invariance and
shared-label removal persist.  This embedding also avoids treating
the physical summands as units modulo \(L_0\) when they are not.

\section{Reconstruction and quantitative comparison}\label{sec:assembly}

The results above require a positive logarithmic saving. We now record
the numerical choices and prove that every approximation fits them.
For the fixed subgroup theorem one may take
\(\delta=1/(8192\varphi(L_0))\). In particular
\(\delta_{12}=1/32768\), the coprime \(C_{12}\) saving is
\(\delta_{12}/2=1/65536\), and the prime-\(3\pmod4\) self-sum
saving is \(\delta_-=1/16384\). These choices concern the
semigroup asymptotics; the ordinary quadratic-indicator statement
has the remainder stated in Section~\ref{sec:quadratic-pairs}.
We assemble the two frequency estimates into
Theorem~\ref{thm:model-comparison} at the sparse scale in (C4).

Write \(\phi=\varphi(L_0)\), \(L=\log X\),
\(\delta_0=1/(8192\phi)\), and use \eqref{eq:comparison-scale}.

\subsection{An explicit discarded-bank exponent}
Primes in the reduced bank class \(c\bmod L_0\) have reciprocal
mass \((1/(8\phi)+o(1))\log L\) in the small range and
\((1/(15\phi)+o(1))\log L\) in the big range. Choose
\[
 b_X=\left\lfloor\frac{\log L}{32\phi}\right\rfloor
\]
disjoint banks of each size, greedily giving each reciprocal mass
\(V_g=1+o(1)\). The overshoot per bank is at most the largest
reciprocal prime. Delete bank primes dividing the displacement.
For a reflected target their total reciprocal mass is
\(O(L/(P_{\min}\log P_{\min}))\); for a fixed shift the deletion
is eventually empty. Thus the total mass error is smaller than every
fixed inverse logarithmic power, uniformly in the reflected target.

The independent law \eqref{eq:bank-law} gives
\begin{equation}\label{eq:success}
 \E I_g=\prod_{p\in\mathcal P_g}(1-2/p)
 \sum_{\substack{p,q\in\mathcal P_g\\p\ne q}}
                  \frac1{(p-2)(q-2)}
 =V_g^2e^{-2V_g}+O\left(\sum_{p\in\mathcal P_g}p^{-2}\right)
 \longrightarrow e^{-2}>1/8.
\end{equation}
Banks are independent in this law. Chernoff's inequality bounds the
probability of fewer than \(b_X/16\) successful banks in either range by
\begin{equation}\label{eq:rare-budget}
 2e^{-b_X/64}\ll L^{-1/(4096\phi)}=L^{-2\delta_0},
\end{equation}
since \(b_X\ge\log L/(64\phi)\) eventually. Apply (C4) to this
bad-bank indicator and use \(|f_i-\mathcal D_i|\le f_i+\mathcal D_i\).
The discarded centered correlation is
\(O(\mathcal S_XL^{-2\delta_0})\), plus its power-smaller
and exponential endpoint errors.

The reconstruction is exact. For \(b\) banks, the Boolean function
asserting at least \(k\) successes has expansion
\(\sum_Tc_T\prod_{g\in T}I_g\), where \(c_T=0\) for \(|T|<k\)
and \(|c_T|=\binom{|T|-1}{k-1}\) otherwise. Binomial inversion in
the number of successful banks proves this identity, and its total
coefficient mass is at most \(3^b\). Multiply the two size-range
identities. For every active set \(T\), exact prime invariance gives
\begin{equation}\label{eq:reconstruction}
 \xi_i(x)w_j(y)\prod_{g\in T}I_g(x,y)
 =\left(\prod_{g\in T}V_g^2\right)
                  \xi_i(x)w_j(y)\Omega_T(x)\Omega_T(y),
 \qquad w_j\in\{f_j,\mathcal D_j\}.
\end{equation}
These weights introduce no representation multiplicity.
The full reconstruction coefficient mass is at most
\[
 3^{2b_X}\prod_g\max(1,V_g^2)\le L
\]
for large \(X\). Indeed the second product is \(1+o(1)\), and
\(2b_X\log3\le(\log3)\log L/(16\phi)\).
Prime squares in active banks have population at most a fixed
logarithmic factor times \(X\sum_{p\in\mathcal B}p^{-2}\),
which is exponentially small. This estimate follows from bounded
endpoints and counting multiples, independently of a pair asymptotic.

\subsection{Moving cutoffs and reflected target localization}
Put \(\epsilon=L^{-\delta_0}\). In the reflected case choose
\(0\le\Phi_\epsilon\le1\), equal to one on
\([2\epsilon,1-2\epsilon]\), supported in
\([\epsilon,1-\epsilon]\), with transitions of width \(\epsilon\).
The same derivative estimates apply when additional interval endpoints
are smoothed at that width. Set
\(\Psi_\epsilon(t)=t\Phi_\epsilon(t)\) and
\(F_\epsilon(u)=\Psi_\epsilon(e^u)\). Differentiation gives
\[
 \|F_\epsilon\|_1\ll1,\qquad
 \|F_\epsilon^{(j)}\|_1\ll_j\epsilon^{1-j}\quad(j\ge1).
\]
At the lower transition logarithmic differentiation cancels its small
scale. Fourier integration by parts yields
\begin{equation}\label{eq:mellin-budget}
 \int|\widehat F_\epsilon(\zeta)|(1+|\zeta|)^r\,d\zeta
 \ll_r\epsilon^{-r-1},\qquad
 \int_{|\zeta|>L^3}|\widehat F_\epsilon(\zeta)|\,d\zeta
 \ll_j L^{-(3-\delta_0)(j-1)}.
\end{equation}
The external phase range stays at \(L^3\); its tail derivative order
is chosen last.

On a contributing reflected edge \(n<0\), \(n'=n+XD>0\),
put \(t=|n|/(XD)\). Then
\begin{equation}\label{eq:two-sided-localization}
 \frac{n'}{|n|}=\frac{1-t}{t}
 \in\left[\frac\epsilon{1-\epsilon},
                \frac{1-\epsilon}{\epsilon}\right].
\end{equation}
For a source dyad \(|n|\asymp Y\), insert a target cutoff on
\([c\epsilon Y,C\epsilon^{-1}Y]\), equal to one on that range,
\emph{before} Mellin inversion, and retain it in each pairing.
Removed edges originally had zero full Mellin integral.
Shared-label division preserves the ratio. There are \(O(\log L)\)
target dyads \(Y'\), each with uniform rescaled logarithmic derivatives.
Their exact-list endpoint moments are \(O(Y)\), \(O(Y')\), by
counting multiples of a fixed list product. This introduces no
exponent depending on the eventual mark length or density depth.
Pairing costs at most \(\sqrt{Y'/Y}\ll\epsilon^{-1/2}\).
The Mellin norm and dyadic count together cost at most
\(L^{5\delta_0/2}\); reserve \(L^{8\delta_0}\) for all localization.

For comparison scale \(Y\asymp XD_1D_2\), use the actual target
dyad \(c\epsilon Y\le Y'\le C\epsilon^{-1}Y\).
The reflected major window has width \(L^{B_{\rm out}}/Y\), so
\(|\beta|Y'\ll L^{B_{\rm out}+\delta_0}\).
Proposition~\ref{prop:gr-energy} applies at this enlarged logarithmic
window exponent. Its integrated rescaling costs at most
\(\epsilon^{-2}\), within the reserve. At the smallest dyad,
\begin{equation}\label{eq:smallest-dyad}
 \frac{Y'}{qP}\gg
 e^{cL^{1/3}-CL^{1/4}}L^{-B_{\rm out}-\delta_0}\gg_A L^A.
\end{equation}
Removing a gcd divisor scales length and modulus equally, so this
bound holds in every target profile.
For a fixed shift the two endpoints are comparable on the original
support; the cutoff and pairing cost is no larger. Its wider major
windows are handled by Proposition~\ref{prop:fixed-comparison-energy},
not by extending \eqref{eq:gr-energy} beyond its proved range.

\subsection{Precision order and completion}
Let \(Q_0=\lceil\kappa_0\rceil+30\).
Request normalized graph-pairing error \(L^{-(Q_0+3)}\) per source
dyad. There are at most \(L^2\) such dyads. Fix root and ideal
accuracies for this request, then the probe complexity and eventual
mark length, keeping the phase range \(L^3\).
Only after this family is fixed let \(C_{\rm cmp}\) dominate its
coefficient and dyadic costs, and \(C_w\) its endpoint pointwise
exponent. These exponents do not depend on the density depth.
If \(m_*\) bounds free draws per group, the dyadic product laws obey
\[
 \nu_i(n)\le L^{K_*}/n,\qquad
 K_*=\frac{\log(m_* !)+m_*\log2}{16\phi}.
\]
Their mass is at most one, so their squared norm at length \(T_i\)
is at most \(L^{K_*}/T_i\). In \eqref{eq:gr-multiplier} take
\(C_{\rm prod}=2K_*+2\), and choose
\begin{align}
 A_{\min}&>C_{\rm cmp}+2C_w+\kappa_0+50+8\delta_0,\notag\\
 B_{\rm out}&>2A_{\min}+C_{\rm prod}+4.\label{eq:outer-budget}
\end{align}
Dirichlet approximation with denominator ceiling
\(T_1T_2/L^{B_{\rm out}-1}\) gives multiplier \(O(L^{-A_{\min}})\)
off the outer grids. Their denominators are \(Xr\) in the reflected
case and \(|h|r\) for a fixed shift, with \(r\le L^{B_{\rm out}}\).
The middle terms in the multiplier bound are superlogarithmically
small because \(T_i\ge\exp(L^{1/3})\).

Request integrated major-energy saving
\(T_E>2(C_{\rm cmp}+C_w+\kappa_0+50+8\delta_0)\).
In the reflected case enlarge the residue-energy accuracy to pay
both arc widths and the number of outer denominators. Exact discrete
Parseval accounts for all target-grid points. In the fixed-shift case
use the low/high Dirichlet-frequency proof at that accuracy, choosing
its smoothing and low-frequency threshold after the fixed comparison
costs. The cofactor-length margin pays its mean-value bound.
Both major contributions are then at most
\(L^{-T_E/2+C_w+C_{\rm cmp}+8\delta_0}\).
The minor bound has the same ample margin.

Choose the model depth last, using (C3) at the largest period,
smoothness and precision exponents now required. Its uniform-in-depth
pointwise and sieve exponents in (C1) and (C4) make this order
noncircular. Choose the external Mellin tail derivative order last
as well. Nothing changes the earlier ideal family or phase range.

\begin{proof}[Completion of Theorem~\ref{thm:model-comparison}]
After localization, source-dyad summation and reconstruction, the graph
and comparison errors are
\[
 O(XL^{-Q_0+8\delta_0})=o(\mathcal S_XL^{-\delta_0}).
\]
The rare-bank term is \(O(\mathcal S_XL^{-2\delta_0})\).
All record, square and CRT tails are exponential or power-smaller
than \(\mathcal S_XL^{-\delta_0}\).
For each ordered endpoint pair use
\[
 f_1(x)f_2(y)-\mathcal D_1(x)\mathcal D_2(y)
 =(f_1(x)-\mathcal D_1(x))f_2(y)+\mathcal D_1(x)(f_2(y)-\mathcal D_2(y)).
\]
The banked comparison bounds each centered term to the required
precision. The argument is symmetric in the endpoints; for a fixed
shift, interchange them and replace \(h\) by \(-h\).

To restrict to an arbitrary interval, smooth its endpoints at width
\(\epsilon X\); in the reflected case first remove the two extreme
strips. At most four strips of total length \(O(\epsilon X)\)
are lost. Condition (C4), applied to each nonnegative pair weight,
bounds their total contribution by
\(O(\epsilon\mathcal S_X+X^{1/2}L^C+XL^Ce^{-cL^{1/10}})\).
For an interval shorter than this width use that upper bound directly.
Otherwise total variation remains \(O(1)\), and a transition at
\(t_0\) has logarithmic derivative norm at most
\(O_j(t_0^j\epsilon^{1-j})\). The Mellin and localization budgets
are therefore uniform in both endpoints. Since
\(\epsilon=L^{-\delta_0}\), every error has the asserted size.
This proves the comparison theorem in both regimes.
\end{proof}

\section{Fractional rough models and one-variable means}\label{sec:model-preparation}
We develop the rough Hankel model and its one-variable centering
estimate before applying the comparison theorem. Fix an even
\(L_0\), put \(G=(\Z/L_0\Z)^\times\), and take a proper subgroup
\(H\le G\). Write
\[
 \mathcal P_H=\{p:p\nmid L_0,\ p\bmod L_0\in H\},\qquad
 C_H=\langle\mathcal P_H\rangle,\qquad f=\1_{C_H},
\]
and set \(\rho=|H|/|G|\), \(k_H=[G:H]=1/\rho\). The outer
scale is \(X\); later reflected applications take \(X=N\).
Put
\[
 L=\log X,\quad y=\exp(L^{2/5}),\quad \mathcal P_y=\{p:p\le y\},\quad
 E=\exp(\sqrt L),\quad V_y=\prod_{p\le y}(1-1/p).
\]
Let \(n_{\mathcal P_y}\) denote the \(\mathcal P_y\)-supported part of \(n\), and
\(n_{>y}=n/n_{\mathcal P_y}\). Define
\[
 E_+(y)=\prod_{\substack{p\le y\\p\in\mathcal P_H}}(1-1/p),\quad
 E_-(y)=\prod_{\substack{p\le y\\p\nmid L_0,\ p\notin\mathcal P_H}}(1-1/p),
 \quad \alpha_0=\varphi(L_0)/L_0.
\]
Then \(V_y=\alpha_0E_+(y)E_-(y)\) for large \(X\).

\subsection{Fractional poles and centered rough means}
Character orthogonality gives, for \(\Re s>1\),
\begin{equation}\label{eq:euler-branches}
 F_H(s)=\prod_{p\in\mathcal P_H}(1-p^{-s})^{-1}
       =G_H(s)\prod_{\eta\in H^\perp}L(s,\eta)^\rho.
\end{equation}
At a forbidden good prime of coset order \(r>1\), the correction
is \((1-p^{-rs})^{1/r}\); thus \(G_H\) converges absolutely for
\(\Re s>1/2\). The literal principal character carries the bad Euler factors.
The positive marginal constant is
\[
 c_H=\frac1{\Gamma(\rho)}
             \lim_{s\downarrow1}(s-1)^\rho F_H(s);
\]
its Euler expression is \eqref{eq:marginal-product}. Every
\(\eta\in H^\perp\) equals one on the support of \(f\); every such
twist therefore has the same pole branch.

Define the rough series and its regular part by
\[
 F_R(s)=F_H(s)\prod_{\substack{p\le y\\p\in\mathcal P_H}}(1-p^{-s})
        =(s-1)^{-\rho}H_R(s).
\]
Here \(H_R(1)=\Gamma(\rho)c_HE_+(y)\). Fixed-order Cauchy estimates
on a disk of radius \(c/\log y\) give
\begin{equation}\label{eq:rough-derivatives}
 |H_R^{(j)}(1)|/j!\ll_j H_R(1)(C\log y)^j.
\end{equation}
For a fixed depth \(K\), define on polynomial scales in \(X\)
\begin{equation}\label{eq:hankel-model}
 M_{R,K}(t)=\sum_{j=0}^K\frac{H_R^{(j)}(1)}{j!}
                 \frac{(\log t)^{\rho-j-1}}{\Gamma(\rho-j)}.
\end{equation}
For \(X^a\le t\le X^C\), with fixed \(0<a<C\), this is
\(c_HE_+(y)(\log t)^{\rho-1}(1+O_K(L^{-3/5}))>0\).
For definiteness fix a smooth real function \(\eta_0\), with
\(0\le\eta_0\le1\), zero on \(({-}\infty,1/4]\) and one on
\([1/3,\infty)\). Set \(M_{R,K}(t)=0\) for \(t\le X^{1/4}\),
and multiply the displayed expression by \(\eta_0(\log t/L)\)
above that point. On the transition range it has the same positive
upper bound. All rough scales used in (C3) are at least
\(X^{1/2}/E>X^{1/3}\), so the cutoff never changes their means.
Only the polynomial bulk enters principal means.

\begin{lemma}\label{lem:rough-means}
For fixed \(B,A_0>0\), sufficiently large fixed \(K\) has the
following property. If \(b\) is bounded and periodic of period at
most \(L^B\), and \(W\) is a smooth dyadic weight at a scale
\(X^{2/5}\le T\le X^2\), with fixed derivatives bounded by
fixed powers of \(L\),
then
\begin{align*}
 &\sum_n W(n)b(n)f(n)\1_{(n,\prod_{p\le y}p)=1}\\
 &\quad-\sum_n W(n)b(n)\1_{(n,\prod_{p\le y}p)=1}
       \left(\sum_{\eta\in H^\perp}\eta(n)\right)
                  \frac{M_{R,K}(n)}{V_y}\ll T L^{-A_0}.
\end{align*}
The same holds with a factor \(n^{iu}\), \(|u|\le L^B\).
\end{lemma}
\begin{proof}
Let \(q\le L^B\) be the period of \(b\). Refine it to \(Q=\lcm(q,L_0)\) and expand on unit residues
by characters. All primes dividing \(Q\) are at most \(y\).
For \(\psi\bmod Q\), the twisted rough series is
\[
 \prod_{\eta\in H^\perp}L(s,\psi\eta)^\rho\,G_\psi(s)
 \prod_{\substack{p\le y\\p\in\mathcal P_H}}(1-\psi(p)p^{-s}).
\]
If \(\psi\) is inflated from \(H^\perp\), this equals \(F_R\);
otherwise all its character factors are nonprincipal.
Use smooth Mellin inversion at height \(\exp(L^{1/10})\), moving
to \(1-c/\log y\). The classical zero-free region is wider there
at these heights and polylogarithmic conductors. Siegel's theorem
places a possible exceptional real zero outside the strip for large
\(X\), after choosing its fixed exponent depending on \(B\).
The finite Euler product costs \(L^{O(1)}\); the displacement gives
\(\exp(-cL^{3/5})\), dominating polynomial \(L\)-function growth at
that height. Smooth tails are smaller than every fixed logarithmic
power. The threshold can be ineffective.

Hankel integration gives \eqref{eq:hankel-model} for each principal
branch. Its relative Taylor error is \(O_K(L^{-3(K+1)/5})\),
by \eqref{eq:rough-derivatives}.
Choose \(K\) after the derivative and character-count costs.
For the unweighted rough population the series
\(\zeta(s)\prod_{p\le y}(1-p^{-s})\) has principal density \(V_y\),
and only the trivial character contributes. The projection
\(\sum_{\eta\in H^\perp}\eta\) therefore reproduces every principal
branch of the first sum with its exact coefficient. The factor
\(n^{iu}\) changes only the derivative budget and choice of \(K\).
\end{proof}

The contour argument uses the fractional pole and the absolute
convergence and fixed-order derivative bounds of its Euler correction.
It therefore also applies when that correction is inverted and remains
absolutely convergent for \(\Re s>1/2\). The quadratic-indicator
application uses precisely this variant, with exponent \(1/2\)
and its two principal character branches.

For \(U\subset H\), set
\begin{equation}\label{eq:corrected-model}
 \mathcal D_U(n)=k_H\1_{n\bmod L_0\in U}f(n_{\mathcal P_y})
                       \frac{M_{R,K}(n_{>y})}{V_y},\qquad
 \xi_U(n)=f(n)\1_{n\bmod L_0\in U}-\mathcal D_U(n).
\end{equation}
Then \(\mathcal D_U\ge0\), \(\mathcal D_U\ll L^{-3(1-\rho)/5}\) and
\(|\xi_U|\ll1\). Conditioning on an exact record \(n_{\mathcal P_y}=r\le E\)
leaves residual state \(U/r\), so Lemma~\ref{lem:rough-means}
centers its principal mean. Fixed divisor-growth record coefficients
cost only powers of \(L\), since
\(\sum_{r\ \mathcal P_y\text{-supported}}d_j(r)/r\ll L^{O_j(1)}\).
Rankin's bound with exponent \(c/\log y\) gives
\begin{equation}\label{eq:record-tail}
\#\{n\le T:n_{\mathcal P_y}>E\}\ll TL^{O(1)}e^{-cL^{1/10}}
 \quad(X^a\le T\le X^C).
\end{equation}
Indeed the Euler moment is \(L^{O(1)}\), whereas
\(E^{-c/\log y}\) gives the exponential saving. The same argument
permits fixed divisor weights. For \(H=G\), use the exact model
\(\mathcal D_U=\1_{n\bmod L_0\in U}\) and \(\xi_U=0\).


\section{Ordinary quadratic indicators and fixed-shift pairs}
\label{sec:quadratic-pairs}

The same comparison also treats ordinary quadratic norm indicators.
In this application inert primes may occur to an even exponent; they
are not excluded altogether. The rough Dirichlet series still has
two principal character branches, and their projection retains the
conditions at ramified primes. This gives a common statement for
the five pair conjectures considered by Smilansky
\cite[Sections 6.1 and 6.3]{smilansky}.

\subsection{A quadratic character and finitely many valuation rules}

Let \(\chi\) be a primitive nonprincipal quadratic Dirichlet
character of conductor \(q\). Fix a finite set of primes
\(B_0\) containing every prime divisor of \(q\), and masks
\(\alpha_p:\Z_{\ge0}\to\{0,1\}\), \(p\in B_0\), with
\(\alpha_p(0)=1\). Define the multiplicative indicator
\[
 f(n)=\prod_{p\in B_0}\alpha_p(v_p(n))
       \prod_{\substack{p\notin B_0\\\chi(p)=-1}}
                                      \1_{2\mid v_p(n)}
 \quad(n\ge1),\qquad f(n)=0\quad(n\le0).
\]
Thus split primes outside \(B_0\) have no valuation restriction.
The finite masks are fixed; no periodicity in their exponent is
required. Write \(F_p(s)=\sum_{e\ge0}\alpha_p(e)p^{-es}\)
for \(p\in B_0\).

The bad-prime factor must retain a joint residue condition. Let
\(x=(x_p)_{p\in B_0}\in\prod_{p\in B_0}\Z_p\), and put
\(e_p=v_p(x_p)\) when all these valuations are finite. For each
\(\ell\mid q\), define a unit modulo \(\ell^{v_\ell(q)}\) by
\[
 u_\ell(x)=\ell^{-e_\ell}x_\ell
            \prod_{\substack{p\in B_0\\p\ne\ell}}p^{-e_p}
                                  \pmod{\ell^{v_\ell(q)}}.
\]
These residues define \(u_{B_0}(x)\in(\Z/q\Z)^\times\) by CRT.
At an ordinary positive integer this is exactly
\(n/n_{B_0}\pmod q\). Set
\begin{align}
 \vartheta_{B_0}(x)
   &=\prod_{p\in B_0}\alpha_p(e_p)
                    \1_{\chi(u_{B_0}(x))=1},\notag\\
 W_{B_0}(h)
   &=\int_{\prod_{p\in B_0}\Z_p}
             \vartheta_{B_0}(x)\vartheta_{B_0}(x+h)\,dx.
                                                        \label{eq:quadratic-bad-mass}
\end{align}
Normalized product Haar measure is used. Define the integrand to
be zero on the null sets with an infinite valuation. In general
\(W_{B_0}\) is coupled: it is not a product of separate
primewise tests.

\begin{theorem}\label{thm:quadratic-indicator-pairs}
Assume the analytic inputs of Section~\ref{sec:inputs}. For the
fixed data above and every fixed \(h\ne0\),
\begin{equation}\label{eq:quadratic-pair-asymptotic}
 \sum_{1\le n\le X}f(n)f(n+h)
       =C_{\chi,\boldsymbol\alpha}(h)\frac X{\log X}
                            +o_{\chi,\boldsymbol\alpha,h}
                                      \left(\frac X{\log X}\right),
\end{equation}
where
\begin{align}
 C_{\chi,\boldsymbol\alpha}(h)
 &=\frac{4L(1,\chi)}\pi\,W_{B_0}(h)
       \prod_{p\in B_0}\frac{1-\chi(p)/p}{1-1/p}\notag\\
 &\quad\cdot
       \prod_{\substack{p\notin B_0,\ \chi(p)=-1\\p\mid h}}
                     \frac{1-p^{-v_p(h)-1}}{1-p^{-1}}.
                                                        \label{eq:quadratic-pair-constant}
\end{align}
If \(W_{B_0}(h)=0\), the pair count is identically zero.
Otherwise the constant is positive. The theorem concerns the
specified Euler indicator; it does not identify that indicator
with the values of an arbitrary principal quadratic form.
\end{theorem}

\begin{proof}
We separate the arithmetic model, its positive upper envelope, and
its correlation constant. Only the final assembly invokes the binary
comparison theorem.

\paragraph{Stage 1: model construction and rough means.}
Choose a fixed even \(L_0\) divisible by \(q\) and all primes
in \(B_0\), and use the prime class one modulo \(L_0\).
Put \(L=\log X\), \(y=e^{L^{2/5}}\), \(\mathcal P_y=\{p:p\le y\}\),
\(E=e^{\sqrt L}\), and \(V_y=\prod_{p\le y}(1-1/p)\).

The exact Euler identity for \(F(s)=\sum f(n)n^{-s}\) is
\begin{equation}\label{eq:quadratic-euler-square}
 F(s)^2=\zeta(s)L(s,\chi)
 \prod_{\substack{p\notin B_0\\\chi(p)=-1}}(1-p^{-2s})^{-1}
 \prod_{p\in B_0}F_p(s)^2(1-p^{-s})(1-\chi(p)p^{-s}).
\end{equation}
The first correction converges absolutely for \(\Re s>1/2\).
Every finite factor is analytic near one and positive there on
the real axis. With this branch convention, the one-variable
constant \(c_f\) satisfies
\begin{equation}\label{eq:quadratic-marginal}
 c_f^2=\frac{L(1,\chi)}\pi
 \prod_{\substack{p\notin B_0\\\chi(p)=-1}}(1-p^{-2})^{-1}
 \prod_{p\in B_0}F_p(1)^2(1-p^{-1})(1-\chi(p)p^{-1}).
\end{equation}
Here \(F_p(1)\) denotes the local factor evaluated at \(s=1\),
not at its series variable. The factor \(\pi=\Gamma(1/2)^2\)
is the Selberg--Delange normalization.

Remove all factors at primes in \(\mathcal P_y\), and write the resulting
rough series as \(F_R(s)=(s-1)^{-1/2}H_R(s)\). Its leading
coefficient is
\[
 H_R(1)/\Gamma(1/2)
          =c_f\prod_{p\le y}F_{f,p}(1)^{-1},
\]
where \(F_{f,p}\) is the actual local factor of \(f\).
The Cauchy bound
\(H_R^{(j)}(1)/j!\ll_j H_R(1)(C\log y)^j\) permits the
fixed-depth density
\[
 M_K(t)=\sum_{j=0}^K\frac{H_R^{(j)}(1)}{j!}
              \frac{(\log t)^{-1/2-j}}{\Gamma(1/2-j)}.
\]
Use the cutoff \(\eta_0(\log t/L)\) from Section~\ref{sec:model-preparation},
setting the model to zero for \(t\le X^{1/4}\). It equals one on the
rough range in (C3), where \(Y/r\ge X^{1/2}/E>X^{1/3}\).
For \(n=av\), \(a=n_{\mathcal P_y}\), define
\begin{equation}\label{eq:quadratic-model}
 \mathcal D(n)=f(a)(1+\chi(v))M_K(v)/V_y.
\end{equation}
For every fixed depth, \(\mathcal D\ll L^{-3/10}\) on the relevant
polynomial ranges, so (C1) holds eventually with a depth-independent
bound one. Multiplication by a split bank prime changes only
\(a\), preserving \(f(a)\) and \(v\) exactly. This proves (C2),
including repeated prime factors.

For (C3), inflate a test character \(\eta\) to a modulus
containing \(q\). The square of its rough Dirichlet series has
factors \(L(s,\eta)L(s,\eta\chi)\), an absolutely convergent
prime-square correction, and the removed small-prime factors.
All primes dividing the polylogarithmic test modulus lie in \(\mathcal P_y\).
Only \(\eta=1\) and \(\eta=\chi\) give a principal branch,
and their rough coefficients agree pointwise. The factor
\(1+\chi(v)\) in \eqref{eq:quadratic-model} reproduces both.
The proof of Lemma~\ref{lem:rough-means} now applies with exponent
\(1/2\): the contour is \(1-c/\log y\), its height is
\(e^{L^{1/10}}\), and the Taylor remainder has relative size
\(O_K(L^{-3(K+1)/5})\). The inverted inert-prime correction
in \eqref{eq:quadratic-euler-square} has the same absolute
convergence and derivative bounds as the correction in that lemma.
The finite masked factors have been removed completely from the
rough variable. Thus they introduce no further rough branches.
Conditioning on \(a\le E\) leaves only its fixed multiplier
\(f(a)\), outside the periodic rough sum. Choosing \(K\)
after the desired logarithmic accuracy proves (C3).

\paragraph{Stage 2: the prime-square sieve envelope.}
For (C4), keep exact bank records \(r,s\le E\); larger ones
have contribution \(XL^{O(1)}e^{-cL^{1/10}}\) by
\eqref{eq:record-tail}. The retained records are coprime because
bank primes dividing \(h\) have been deleted. Their divisibilities
give one progression with \(|J|/(rs)+O(1)\) terms.
Exactness contributes two excluded bank roots outside \(rs\),
and one at a prime of \(rs\). Together with \(1/(rs)\), these
are exactly the record masses \eqref{eq:bank-law}.

At an inert prime outside \(B_0\), a necessary condition is
\(v_p(n)\ne1\). It excludes \(p-1\) classes modulo \(p^2\).
For \(p\nmid h\), the two endpoints exclude \(2(p-1)\)
classes; bank-record slopes are units at these primes. Drop the
finitely many primes dividing \(2h\prod_{p\in B_0}p\) for
an upper bound. The nonbank sieve dimension is one, with a
convergent \(1+O(p^{-2})\) correction to the ordinary two-root
product. A Selberg upper sieve with prime-square atoms, level
\(z=X^{1/100}\), and \(w=z^{1/16}\), has endpoint error
\(O(z^{4+o(1)})\) per retained record. This follows by summing
the CRT root errors: a least-common-multiple atom has at most
\(2^{\omega(d)}d\) roots and modulus \(d^2\).
Summing the \(E^2=X^{o(1)}\) records remains power smaller
than \(X^{1/2}\).

For \(ff\), \(\mathcal D f\), and \(\mathcal D\mathcal D\), the nonbank sieve
denominators have orders respectively
\(\log w\), \((\log y\log w)^{1/2}\), and \(\log y\).
The bound \(\mathcal D\ll L^{-3/10}\) supplies the missing factor
on each model side. Explicitly, since \(\log y=L^{2/5}\)
and \(\log w\asymp L\), their factors are
\[
 (\log w)^{-1}\ll L^{-1},\qquad
 \frac{L^{-3/10}}{(\log y\log w)^{1/2}}\ll L^{-1},\qquad
 \frac{L^{-3/5}}{\log y}=L^{-1}.
\]
All three cases therefore give the common
scale \(\mathcal S_X=X/L\), with the exact bank masses.
This proves (C4), uniformly in the interval \(J\). The exponents
in the exceptional-record errors do not depend on the final
fixed depth.

\paragraph{Stage 3: model correlation and the arithmetic constant.}
We calculate the model main term separately from the analytic replacement. Let
\(I_y=\{p\le y:p\notin B_0,\ \chi(p)=-1\}\) and put
\[
 \Lambda_y=\frac{2c_f}{\sqrt L}
      \prod_{p\in B_0}\frac{F_p(1)^{-1}}{1-p^{-1}}
      \prod_{p\in I_y}(1+p^{-1}).
\]
Outside the negligible records \(n_{\mathcal P_y}>E\), the rough cofactor
has \(\log v=L+O(\sqrt L)\). If the small inert exponents
are even, its character is \(\chi(n/n_{B_0})\). Hence
\begin{equation}\label{eq:quadratic-model-leading}
 \mathcal D(n)=\Lambda_y\vartheta_{B_0}(n)
            \prod_{p\in I_y}\1_{2\mid v_p(n)}
                                      (1+O_K(L^{-1/2})).
\end{equation}
At an inert prime the local pair density is
\begin{equation}\label{eq:quadratic-inert-pair}
 \delta_p(h)=\frac p{p+1}(1-p^{-v_p(h)-1}).
\end{equation}
To see this, put \(e=v_p(h)\) and partition \(v_p(n)\)
below, equal to, and above \(e\). Below \(e\), both endpoints
have the same valuation, and the permitted even terms form a
geometric sum. If \(e\) is odd there are no further terms.
If \(e\) is even, the stratum with both valuations equal to
\(e\) has mass \(p^{-e}(p-2)/p\); the two possible higher
even tails have combined mass \(2p^{-e}/(p(p+1))\).
The two parity cases give \eqref{eq:quadratic-inert-pair}, also
for \(p=2\).

We justify the product of these densities at the growing cutoff.
Extract the inert square parts \(d^2,e^2\), with \(d,e\)
supported on \(I_y\), and retain \(d,e\le E\). Truncate
each bad-prime power at \(E\). The omitted integer count is
\(O_{B_0,h}(X/E+\sqrt X)\), and the omitted Haar mass is
\(O_{B_0,h}(E^{-1})\). Each retained joint record and residual
unit state gives an empty stratum or finitely many progressions
modulo \(X^{o(1)}\). On them the two residual forms must be
units at the primes in \(I_y\). Empty root cells are discarded;
the remaining sieve dimension is at most one. There are at most
\(O_{B_0,q}(E^2(\log E)^{2|B_0|})\) retained cells: two
square prefixes, two bounded valuation lists, and a fixed number
of residual unit choices. Their moduli are at most
\(O_{B_0,q}(E^{4+2|B_0|})\). Use upper and lower
fundamental-lemma weights of level \(D_0=X^{1/100}\).
Their parameter \(\log D_0/\log y\asymp L^{3/5}\)
gives relative error smaller than every log power. Since the
root count at a squarefree sieve divisor is at most
\(2^{\omega(d)}\), a conservative total endpoint bound is
\[
 E^{O_{B_0,q}(1)}D_0^2(\log D_0)^{O(1)}
                         =X^{1/50+o(1)}.
\]
This is negligible compared with \(X/L^A\) for every fixed
\(A\). Their main masses are summable: compatibility requires
\((d,e)^2\mid h\), and their progression densities are
\(\ll_h1/\operatorname{lcm}(d^2,e^2)\), with geometric
bad-prime tails. Summing the disjoint valuation cells proves
\[
 \sum_{X\le n\le2X}\mathcal D(n)\mathcal D(n+h)
   =\Lambda_y^2XW_{B_0}(h)\prod_{p\in I_y}\delta_p(h)
                                                 +o_h(X/L).
\]
For \(p\nmid h\),
\((1+p^{-1})^2\delta_p(h)=1-p^{-2}\); for \(p\mid h\)
the additional ratio is
\((1-p^{-v_p(h)-1})/(1-p^{-1})\). Substitution of
\eqref{eq:quadratic-marginal} cancels every \(F_p(1)\)
and the inert Euler product, leaving
\eqref{eq:quadratic-pair-constant}. Thus the model correlation
has the asserted constant on each dyadic interval.

\paragraph{Assembly.}
Stages 1 and 2 establish \textup{(C1)--(C4)} at scale \(X/L\).
The comparison theorem now gives
\begin{equation}\label{eq:quadratic-model-comparison}
 \sum_{X\le n\le2X}\{f(n)f(n+h)-\mathcal D(n)\mathcal D(n+h)\}
                                      =o_h(X/L).
\end{equation}
Combining this with Stage 3 proves the dyadic asymptotic for the
ordinary indicator. Summing dyadically gives
\eqref{eq:quadratic-pair-asymptotic}; the upper bound from (C4)
controls the part below \(\varepsilon X\), after which
\(\varepsilon\) tends to zero.

Finally, if an actual positive pair has \(f(n)f(n+h)=1\),
then both residual characters equal one. Its valuations and
residual units are preserved on a sufficiently deep joint
\(B_0\)-adic cylinder. That cylinder has positive Haar measure,
so \(W_{B_0}(h)>0\). This proves the exact zero assertion.
\end{proof}

\subsection{The five quadratic forms}

Let \(b_d(n)\) indicate the positive integers of the form
\(x^2+dy^2\), with \(x,y\in\Z\). Coordinates may vanish,
and no coprimality condition is imposed. Put \(b_d(n)=0\)
for \(n\le0\). Smilansky's
prime-factor criterion \cite[Theorem 10]{smilansky} identifies
these ordinary indicators with the preceding Euler indicators:
\[
\begin{array}{c|c|c|c}
d&\chi&B_0&\text{nontrivial mask in }B_0\\\hline
1&\chi_{-4}&\{2\}&\text{none}\\
2&\chi_{-8}&\{2\}&\text{none}\\
3&\chi_{-3}&\{3\}&\text{none}\\
4&\chi_{-4}&\{2\}&\alpha_2(e)=\1_{e\ne1}\\
7&\chi_{-7}&\{2,7\}&\alpha_2(e)=\1_{e\ne1}
\end{array}
\]
The corresponding order discriminants are \(-4,-8,-12,-16,-28\).
In particular the third, fourth and fifth forms are conductor-two
orders. For \(d=3\), the ordinary value set nevertheless has
exactly the inert-even rule for \(\chi_{-3}\), including the
prime two; no extra dyadic mask is imposed.



\begin{proof}[Proof of Corollary~\ref{cor:smilansky-five-pairs}]
Write \(Q_d=\{p:\chi(p)=-1\}\), except that no prime in
\(B_0\) is included; for the displayed five forms this exception
changes none of the inert sets. These are respectively
\[
 \{p\equiv3\ (4)\},\quad \{p\equiv5,7\ (8)\},\quad
 \{p\equiv2\ (3)\},\quad \{p\equiv3\ (4)\},\quad
 \{p\equiv3,5,6\ (7)\}.
\]
Let \(R_1=R_2=R_4=\{2\}\), \(R_3=\{3\}\), and
\(R_7=\{2,7\}\). The bad-prime integrand factors in these
cases. For \(d=7\), this uses \(\chi_{-7}(2)=1\), which
removes the cross-valuation phase from its ramified unit test.
Put \(a_p=v_p(|h|)\). The local densities are
\begin{align*}
 W_{1,2}(h)&=\begin{cases}
  1/4,&a_2=0,\\ 1/2-3/2^{a_2+2},&a_2\ge1,
 \end{cases}
 &W_{2,2}(h)&=\begin{cases}
  1/4,&a_2=0,1,\\ 1/2-3/2^{a_2+1},&a_2\ge2,
 \end{cases}\\[2mm]
 W_{3,3}(h)&=\tfrac12(1-2/3^{a_3+1}),
 &W_{7,7}(h)&=\tfrac12(1-4/7^{a_7+1}),\\[2mm]
 W_{4,2}(h)&=\begin{cases}
  1/8,&a_2=0,\\0,&a_2=1,\\5/16,&a_2=2,\\
  3/8-3/2^{a_2+2},&a_2\ge3,
 \end{cases}
 &W_{7,2}(h)&=\begin{cases}
  1/2,&a_2=0,1,\\3/4,&a_2\ge2.
 \end{cases}
\end{align*}
For completeness, the odd ramified unit predicate is a square
modulo \(p=3\) or \(7\). At \(a_p=0\), there are
\((p-3)/4\) nonzero square pairs modulo \(p\), and exactly
one zero neighborhood contributes \(1/(2p)\). The result is
\((p-1)/(4p)\). For \(a_p\ge1\), the recurrence
\(W(a_p)=(p-1)/(2p)+W(a_p-1)/p\) gives the displayed formulas.
At two the ramified unit predicates are respectively one modulo
four and one or three modulo eight. Their initial pair masses
are \(1/4\), and, for the latter, \(1/4\) also when \(a_2=1\).
The former has mass \(1/8\) at \(a_2=1\); thereafter its
recurrence is \(W(a)=1/4+W(a-1)/2\). The latter has mass
\(1/8\) at \(a_2=2\) and obeys that recurrence for \(a\ge3\).
Removing valuation one gives the \(d=4\) row: for \(a\ge2\)
its mass is \(1/4+W_{1,2}(h/4)/4\). The \(d=7\) dyadic
row counts only the two valuation masks, giving its stated masses.

Using primitive characters, the prefactors are
\[
 c_d=\frac{4L(1,\chi)}\pi
                 \prod_{p\mid\operatorname{cond}(\chi)}(1-p^{-1})^{-1},
 \qquad
 (c_1,c_2,c_3,c_4,c_7)
       =(2,2\sqrt2,2/\sqrt3,2,2\sqrt7/3).
\]
For \(d=7\) the extra bad prime two contributes ratio one
in \eqref{eq:quadratic-pair-constant}. Consequently
\[
 \mathscr T_{d,h}=c_d\prod_{p\in R_d}W_{d,p}(h)
        \prod_{\substack{p\in Q_d\\p\mid h}}
                    \frac{1-p^{-a_p-1}}{1-p^{-1}},
\]
as required. The primitive values are
\(L(1,\chi_{-3})=\pi/(3\sqrt3)\) and
\(L(1,\chi_{-7})=\pi/\sqrt7\); Smilansky's series in
his equation (6.3) omits the Euler factor at two.
At \(h=1\) the five constants are
\[
 (\mathscr T_{1,1},\mathscr T_{2,1},\mathscr T_{3,1},
      \mathscr T_{4,1},\mathscr T_{7,1})
   =\left(\frac12,\frac1{\sqrt2},\frac1{3\sqrt3},
                        \frac14,\frac1{2\sqrt7}\right).
\]
Only \(d=4\) has a zero coefficient for nonzero shifts:
\(v_2(h)=1\). Its pair count is exactly zero because its
values are zero or one modulo four. All other displayed
coefficients are positive.
\end{proof}

\section{Splitting semigroups and prime cosets}\label{sec:semigroup-apps}
\subsection{Subgroup sums and congruence admissibility}
Throughout, \(L_0\ge2\) is even,
\(G=(\Z/L_0\Z)^\times\), and \(H_1,H_2\le G\) are fixed.
Define
\[
 \mathcal P_i=\{p\nmid L_0:p\bmod L_0\in H_i\},\quad
 C_i=\langle\mathcal P_i\rangle,\quad
 \rho_i=\frac{|H_i|}{\varphi(L_0)},\quad
 \kappa=2-\rho_1-\rho_2.
\]
An empty product is one. Fix \(U\subset H_1\), \(V\subset H_2\),
and coprime positive \(L_0\)-supported integers \(A,B\).
Write
\[
 R_{U,V}^{A,B}(N;I)=\#\{Ax+By=N:x\in C_1,\ y\in C_2,
      \ x\bmod L_0\in U,\ y\bmod L_0\in V,\ Ax/N\in I\}.
\]
Omit \(I\) when \(I=[0,1]\). Put
\begin{align}
 c_i&=\frac{1}{\Gamma(\rho_i)}
    \lim_{s\downarrow1}(s-1)^{\rho_i}
              \prod_{p\in\mathcal P_i}(1-p^{-s})^{-1},\label{eq:marginal}\\
 \Gamma_N&=\frac{L_0}{|H_1||H_2|}
    \#\{(u,v)\in U\times V:Au+Bv\equiv N\pmod{L_0}\},\label{eq:finite}\\
 \mathcal E&=\{p\nmid L_0:p\notin\mathcal P_1,
                                  \ p\notin\mathcal P_2\},\notag\\
 C_{\mathcal E}&=\prod_{p\in\mathcal E}
               \left(1-\frac{1}{(p-1)^2}\right),\qquad
 J_{\mathcal E}(N)=\prod_{\substack{p\mid N\\p\in\mathcal E}}
                              \frac{p-1}{p-2}.\label{eq:singular}
\end{align}
The infinite product in \eqref{eq:singular} converges absolutely and
is positive, since \(2\mid L_0\). The constants \(c_i\) are positive.
For an explicit convergent expression, let \(r_i(p)>1\) denote the
order of \(pH_i\) in \(G/H_i\) at a forbidden good prime. Then
\begin{equation}\label{eq:marginal-product}
 c_i=\frac{(\varphi(L_0)/L_0)^{\rho_i}}{\Gamma(\rho_i)}
 \prod_{\substack{\eta\in H_i^\perp\\\eta\ne1}}L(1,\eta)^{\rho_i}
 \prod_{\substack{p\nmid L_0\\p\notin\mathcal P_i}}
                     (1-p^{-r_i(p)})^{1/r_i(p)}.
\end{equation}
Characters here are literal modulo \(L_0\). Continuation along the
positive real axis from \(s>1\) fixes the powers; conjugate factors
are paired. The Euler identity proving this formula is given in
Section~\ref{sec:model}.

\begin{theorem}\label{thm:general}
Under the analytic inputs of Section~\ref{sec:inputs}, with all these
data fixed, there is an explicit \(\delta>0\)
depending only on \(L_0\) such that, uniformly in \(N\) and every interval
\(I\subset[0,1]\),
\begin{align*}
 R_{U,V}^{A,B}(N;I)
 &=c_1c_2\Gamma_N C_{\mathcal E}|I|
       \frac{NJ_{\mathcal E}(N)}{AB(\log N)^\kappa}\\
 &\quad+O\left(\frac{NJ_{\mathcal E}(N)}
                         {AB(\log N)^{\kappa+\delta}}\right).
\end{align*}
If \(\Gamma_N=0\), the count is identically zero.
For \(\Gamma_N>0\) and \(I=[0,1]\), the error is a relative
\(O((\log N)^{-\delta})\). More generally it is relative \(o(1)\)
when \(|I|(\log N)^\delta\to\infty\).
\end{theorem}

The implicit constant and threshold may depend on all fixed data.
An admissible numerical choice of \(\delta\) is recorded in
Section~\ref{sec:assembly}, together with the localization budgets.
No numerical threshold, uniformity in a growing modulus, or
power-short interval assertion is included. For noncoprime fixed
coefficients, first divide by their gcd when it divides \(N\).
The even-modulus convention makes the prime two a bad prime throughout
the analytic theorem. Passing from an odd modulus to its even multiple
requires retaining this convention; freely admitting the prime two is
a separate valuation-profile operation, not a change of notation.

\begin{theorem}\label{thm:c12}
Assume the analytic inputs of Section~\ref{sec:inputs}.
Let \(\mathcal E_{12}=\{p\nmid6:p\not\equiv1\pmod{12}\}\), and set
\[
 A_{12}=\frac{4\pi}{\Gamma(1/4)^2}
       \sqrt{\frac{\log(2+\sqrt3)}{6}}
       \prod_{p\in\mathcal E_{12}}
                        \left(1-\frac{2}{p(p-1)}\right),\qquad
 J_{12}(N)=\prod_{\substack{p\mid N\\p\in\mathcal E_{12}}}
                                      \frac{p-1}{p-2}.
\]
There is an explicit \(\delta_{12}>0\) such that, for all sufficiently
large \(N\equiv2\pmod{12}\),
\[
 R_{12}(N)=A_{12}J_{12}(N)\frac{N}{(\log N)^{3/2}}
                 \left(1+O((\log N)^{-\delta_{12}})\right).
\]
All other residue classes have zero count.
\end{theorem}
\begin{proof}[Deduction from Theorem~\ref{thm:general}]
Take \(L_0=12\), \(H_1=H_2=U=V=\{1\}\), \(A=B=1\).
Then \(\rho_i=1/4\) and \(\Gamma_N=12\) on the indicated class.
The four characters give the exact identity
\begin{align*}
 F_{12}(s)^4={}&\zeta(s)L(s,\chi_{-4})L(s,\chi_{-3})L(s,\chi_{12})\,
        (1-2^{-2s})(1-3^{-2s})
        \prod_{p\in\mathcal E_{12}}(1-p^{-2s})^2.
\end{align*}
Here the three nonprincipal characters displayed are primitive.
Use \(L(1,\chi_{-4})=\pi/4\),
\(L(1,\chi_{-3})=\pi/(3\sqrt3)\), and
\[
 L(1,\chi_{12})=\int_0^1\frac{1-x^4}{1+x^6}\,dx
 =\left[\frac{1}{2\sqrt3}
 \log\frac{x^2+\sqrt3x+1}{x^2-\sqrt3x+1}\right]_0^1
 =\frac{\log(2+\sqrt3)}{\sqrt3}.
\]
Consequently
\[
 c_{12}=\frac{1}{\Gamma(1/4)}
 \left(\frac{\pi^2\log(2+\sqrt3)}{54}
       \prod_{p\in\mathcal E_{12}}(1-p^{-2})^2\right)^{1/4}.
\]
In \(12c_{12}^2C_{\mathcal E_{12}}\), combine the factors by
\((1-p^{-2})(1-(p-1)^{-2})=1-2/(p(p-1))\).
This is exactly \(A_{12}\). The real density integrates to one;
there is no extra beta-integral factor.
\end{proof}

The uniformity in the target also permits a coprimality restriction.
Only the primes allowed in the summands can divide their gcd.
Removing them changes the leading term by the following finite product.
\begin{corollary}\label{cor:coprime-c12}
Assume the analytic inputs of Section~\ref{sec:inputs}.
Put
\[
 R_{\mathrm{cop}}(N)=\#\{(a,b)\in C_{12}^2:a+b=N,\ (a,b)=1\},
 \qquad
 P_{12}(N)=\prod_{\substack{p\mid N\\p\equiv1\ (12)}}(1-p^{-1}).
\]
Uniformly over sufficiently large \(N\equiv2\pmod{12}\),
\[
 R_{\mathrm{cop}}(N)
 =A_{12}J_{12}(N)P_{12}(N)\frac{N}{(\log N)^{3/2}}
       \left(1+O((\log N)^{-\delta_{12}/2})\right).
\]
Every other residue class has zero count. In particular, every
sufficiently large admissible target has coprime summands.
\end{corollary}
\begin{proof}
M\"obius inversion gives the exact identity
\begin{equation}\label{eq:coprime-inversion}
 R_{\mathrm{cop}}(N)
 =\sum_{\substack{d\mid N\\d\in C_{12}}}\mu(d)R_{12}(N/d).
\end{equation}
Indeed, division by any common divisor preserves both prime supports,
and every common divisor lies in \(C_{12}\). Moreover \(d\equiv1\pmod{12}\),
so \(N/d\equiv2\pmod{12}\), and
\(J_{12}(N/d)=J_{12}(N)\).

Write \(L=\log N\), \(\ell=\log\log(3N)\),
\(P=P_{12}(N)\), and \(\delta=\delta_{12}\).
We use the elementary divisor estimates
\begin{align}
 S_0&=\sum_{\substack{d\mid N\\d\in C_{12}}}\frac{|\mu(d)|}{d}
       =\prod_{\substack{p\mid N\\p\equiv1\ (12)}}(1+p^{-1})
       \ll\ell,\qquad P^{-1}\ll\ell,\label{eq:coprime-products}\\
 S_1&=\sum_{\substack{d\mid N\\d\in C_{12}}}
                          \frac{|\mu(d)|\log d}{d}
       =S_0\sum_{\substack{p\mid N\\p\equiv1\ (12)}}
                                     \frac{\log p}{p+1}
       \ll\ell^3.\label{eq:coprime-logmoment}
\end{align}
For \eqref{eq:coprime-products}, split the prime products at \(L\).
Mertens' product bound controls primes at most \(L\) by \(O(\log L)\);
the remaining factors have bounded product, since
\(\sum_{p\mid N,\ p>L}p^{-1}\le1/\log L\).
For \eqref{eq:coprime-logmoment}, the primes at most \(L\) contribute
at most \(\sum_{n\le L}(\log n)/n\ll(\log L)^2\), and the others
contribute at most \(L^{-1}\sum_{p\mid N}\log p\le1\).

Split \eqref{eq:coprime-inversion} at \(d=\sqrt N\).
The bound \(R_{12}(M)\le M\) shows that its large-divisor part is
\(O(\sqrt N\,\tau(N))\), where \(\tau\) is the divisor function.
For \(d\le\sqrt N\), Theorem~\ref{thm:c12} applies uniformly, and
\[
 (L-\log d)^{-3/2}
       =L^{-3/2}+O(L^{-5/2}\log d).
\]
Consequently, using \(\sum_{d\mid N,\ d\in C_{12}}\mu(d)/d=P\)
and extending the main sum back to all divisors,
\begin{align*}
 R_{\mathrm{cop}}(N)
 ={}&A_{12}J_{12}(N)\frac{N}{L^{3/2}}
       \left(P+O(L^{-\delta}S_0+L^{-1}S_1
                            +N^{-1/2}\tau(N))\right)\\
 &\quad+O(\sqrt N\,\tau(N)).
\end{align*}
After division by the main term, the error is
\[
 O\!\left(L^{-\delta}\ell^2+L^{-1}\ell^4
                  +N^{-1/2}\tau(N)L^{3/2}\ell\right).
\]
Here \(J_{12}(N)\ge1\) and \(\tau(N)=N^{o(1)}\).
The displayed error is \(O(L^{-\delta/2})\), proving the claim.
\end{proof}

\begin{remark}[Comparison with the quartic-norm target weight]
The weight in \cite[Theorem 2]{iwaniecpomykala} is
\(W_{12}(N)=\prod_{p\mid N}(1-\varepsilon_p/p)\), with
\(\varepsilon_p=1\) on allowed primes and \(-1\) otherwise.
On the admissible class,
\[
 \frac{J_{12}(N)P_{12}(N)}{W_{12}(N)}
 =\frac23\prod_{\substack{p\mid N\\p\equiv5,7,11\ (12)}}
                 \frac{p(p-1)}{(p-2)(p+1)}\asymp1
\]
uniformly in \(N\), since the factors are \(1+O(p^{-2})\).
The factor \(2/3\) comes from \(2\mid N\); the prime three does
not divide an admissible target. Thus the coprime asymptotic has
the same order as the earlier conditional coprime bounds.
\end{remark}

We next clarify the congruence criterion without analytic inputs.
Write \(C_H(U)=\{n\in C_H:n\bmod L_0\in U\}\).
\begin{lemma}\label{lem:finite-images}
A residue \(r\bmod M\) is represented by \(C_H(U)\) if and only if
its reduction modulo \(\gcd(M,L_0)\) is that of some \(h\in U\),
and every forbidden prime \(p\mid M\), \(p\nmid L_0\), satisfies
\(p\nmid r\).
\end{lemma}
\begin{proof}
Necessity follows from prime support. Choose such an \(h\), and for
each allowed \(p^k\Vert M\) put \(e_p=\min(v_p(r),k)\).
Set \(D=\prod p^{e_p}\), so \(D\in C_H\) and \((D,L_0)=1\).
We seek a prime \(\ell\equiv hD^{-1}\pmod{L_0}\).
Where \(e_p<k\), impose
\[
 \ell\equiv(r/p^{e_p})(D/p^{e_p})^{-1}\pmod{p^{k-e_p}}.
\]
Where \(e_p=k\), impose \(\ell\equiv1\pmod p\).
At primes of \(M\) not dividing \(D\), impose the unit condition
\(\ell\equiv rD^{-1}\pmod{p^k}\). These congruences are compatible
with the \(L_0\)-condition precisely by the stated hypotheses.
They give a reduced CRT class. Dirichlet's theorem supplies a prime
\(\ell\) outside the finite exceptional set; then \(D\ell\) has
the required residue and state.
\end{proof}

\begin{theorem}\label{thm:local}
The condition \(\Gamma_N>0\) is equivalent to solvability of
\(Ax+By\equiv N\pmod M\) with actual
\(x\in C_{H_1}(U),y\in C_{H_2}(V)\), for every \(M\ge1\).
Under the analytic inputs of Section~\ref{sec:inputs}, for all
sufficiently large \(N\), it is also equivalent to positive additive
representability.
\end{theorem}
\begin{proof}
Taking \(M=L_0\) proves necessity. For sufficiency choose a root
counted by \(\Gamma_N\) and work modulo \(K=\lcm(M,L_0)\).
At a prime dividing \(L_0\), at least one coefficient is a unit;
adjust its variable to lift the equation while preserving the fixed
unit states. At a prime outside \(L_0\), both coefficients are units
and \(p\ge3\). Of the \(p\) solutions modulo \(p\), requiring a
unit on each forbidden side removes at most two. A surviving pair
lifts to every higher power. CRT and Lemma~\ref{lem:finite-images}
realize the required two residues by actual semigroup values.
Finally Theorem~\ref{thm:general} is positive for all sufficiently
large targets with \(\Gamma_N>0\).
\end{proof}

\begin{remark}[A joint congruence obstruction]\label{rem:coupled}
Take \(L_0=12\), \(H=\{1,11\}=\ker\chi_{12}\).
Then \(H+H=\{0,2,10\}\pmod{12}\), so \(N\equiv4\pmod{12}\)
is never a sum. Nevertheless \(H\) projects onto the full unit groups
modulo four and three. The separate closures are \(\Z_2^\times\)
and \(\Z_3^\times\); at other primes they are \(\Z_p\) or
\(\Z_p^\times\), according as that prime is allowed or forbidden.
Two units sum to every element of \(\Z_p\) for odd \(p\), and to
every even element for \(p=2\). Thus these targets pass every
separate \(p\)-adic test. The missing condition is the simultaneous
residue state modulo twelve.
\end{remark}

\subsection{Subgroup applications}\label{sec:model}

Use the separate models of Section~\ref{sec:model-preparation} for
the two subgroup restrictions, with outer scale \(X=N\).
Thus \(L=\log N\), \(y=\exp(L^{2/5})\), and
\(E=\exp(\sqrt L)\). We evaluate their correlation and verify the
positive record envelope.

\subsubsection{The model correlation}

Use separate constructions for the two subgroups with common \(\mathcal P_y\).
\begin{proposition}\label{prop:model-correlation}
For a smooth \(0\le\Phi\le1\) supported away from zero and one,
\begin{align*}
 &\sum_{Ax+By=N}\Phi(Ax/N)\mathcal D_U(x)\mathcal D_V(y)\\
 &\quad=c_1c_2\Gamma_N C_{\mathcal E}J_{\mathcal E}(N)
       \frac{N}{AB L^\kappa}
       \left(\int_0^1\Phi(t)\,dt+O(L^{-1/2})\right)+O(N^{1/2}).
\end{align*}
The statement is uniform in target valuations. Endpoint strips of
relative length \(\epsilon\) have upper bound
\(O(\epsilon NJ_{\mathcal E}(N)L^{-\kappa}+N^{1/2})\).
For every interval \(I\subset[0,1]\), the unsmoothed model correlation
has the same leading coefficient multiplied by \(|I|\), with absolute
error \(O(NJ_{\mathcal E}(N)L^{-\kappa-\delta_0})\), where
\(\delta_0=1/(8192\varphi(L_0))\).
\end{proposition}
\begin{proof}
Discard records larger than \(E\) using \eqref{eq:record-tail}.
The rough logarithms on the rest are \(L+O(\sqrt L)\), giving
\[
 \frac{M_{R,K}(n_{>y})}{V_y}
 =\frac{c_H L^{\rho-1}}{\alpha_0E_-(y)}(1+O_K(L^{-1/2})).
\]
Count the remaining small-prime sieve set in the fixed residues.
The fundamental lemma applies at level \(N^{\epsilon_0}\), for
small fixed \(\epsilon_0>0\), with
\(u=\log(N^{\epsilon_0})/\log y\asymp L^{3/5}\).
Its relative error is \(e^{-cu\log u}\), and CRT gives endpoint
error \(O(\nu(d))\) at each squarefree sieving modulus \(d\).
Summing those errors costs \(N^{\epsilon_0}L^{O(1)}\).
This is the standard fixed-dimensional fundamental-lemma application
\cite[Theorem 11.7]{richert}; target-root mergers
decrease its dimension and preserve uniformity.

The positive integer line is \((x,y)=(x_0+Bt,y_0-At)\), with
parameter length \(N/(AB)\). Modulo \(L_0\), this parameter bijects
onto all roots: at each prime at least one coefficient is a unit.
If \(R\) roots have the prescribed states, their raw probability is
\(R/L_0\). Multiplying the normalized marginal projections gives
\[
 \frac{R}{L_0}\frac{L_0c_1}{|H_1|}\frac{L_0c_2}{|H_2|}
       =c_1c_2\Gamma_N.
\]
A prime forbidden on only one side has relative factor one.
A prime forbidden on both sides has two excluded roots, merging
exactly when \(p\mid N\), and relative factor
\[
 \frac{1-(2-\1_{p\mid N})/p}{(1-1/p)^2}.
\]
Their product up to \(y\) is
\(C_{\mathcal E}(y)J_{\mathcal E,y}(N)\). The ordinary tail is
\(O(1/y)\), and the target-prime tail is \(O(L/(y\log y))\).
Partial summation inserts \(\Phi\), costing its total variation.
The same upper sieve on an endpoint strip gives its length times
the sparse density, with a power-smaller endpoint error.
To obtain the uniform interval assertion, put
\(\epsilon=L^{-\delta_0}\) and remove strips of this width near
zero and one. On the rest the rough logarithms are
\(L+O(\sqrt L+\delta_0\log L)\), giving the same relative
\(O(L^{-1/2})\) replacement. The fundamental-lemma count in any
remaining interval has the same CRT endpoint error; partial summation
costs the total variation, bounded independently of the interval.
Restoring the strips costs \(O(\epsilon)\) of the sparse scale.
All other tails are superlogarithmic or power-smaller, proving the claim.
\end{proof}

\subsubsection{Verification of the comparison hypotheses}
The indicators lie between zero and one, and the model bound after
\eqref{eq:corrected-model} gives \(\mathcal D_U\le1\) eventually for each
fixed depth. Thus (C1) holds, with a threshold allowed to depend on
that depth. The full-unit-group exact model also lies between zero
and one.

Choose banks of primes \(p\equiv1\pmod{L_0}\), with small banks in
\([\exp(L^{1/8}),\exp(L^{1/4})]\) and big banks in
\([\exp(L^{1/3}),\exp(L^{2/5})]\). Delete primes dividing \(N\).
Every bank prime is a unit at \(A,B\), belongs to both allowed
sets, and preserves the fixed states. Hence
\begin{equation}\label{eq:exact-invariance}
 f_U(pn)=f_U(n),\quad \mathcal D_U(pn)=\mathcal D_U(n),\quad \xi_U(pn)=\xi_U(n)
\end{equation}
exactly, including higher powers. The deletion changes reciprocal
mass by at most \(O(L/(P_{\min}\log P_{\min}))\), which is smaller
than every fixed inverse logarithmic power. This verifies (C2).

For (C3), fix an \(\mathcal P_y\)-smooth record \(r\le E\). If \(f(r)=0\),
both the indicator and model on that record vanish. Otherwise
\(r\bmod L_0\in H\), and insert the residual state \(U/r\)
into the periodic test of Lemma~\ref{lem:rough-means}. Its smooth
scale is \(Y/r\), so the error is \((Y/r)L^{-D}\), as required.
An additional periodic test of period at most \(L^B\) increases the
period by only a fixed factor. For every required accuracy, the
Hankel remainder works for every sufficiently large fixed depth;
constants and thresholds, rather than earlier graph exponents, may
depend on that depth. This proves (C3) with all surviving character
branches retained.

It remains to check the positive envelope (C4). Write
\(\mu_{\mathcal B}\) for the law \eqref{eq:bank-law}.

\begin{lemma}\label{lem:envelope}
For a nonnegative record function \(F(r,s)\le L^B\), and either
indicator or model on each side, its weighted correlation on the
positive line restricted by \(Ax\in J\subset[0,N]\) is at most
\[
 C\frac{|J|J_{\mathcal E}(N)}{AB L^\kappa}\E_{\mu_{\mathcal B}}F
       +NL^{O(1)}e^{-cL^{1/10}}+O(N^{1/2}).
\]
The constant may depend on the fixed coefficients.
\end{lemma}
\begin{proof}
Discard records \(r,s>E\) using \eqref{eq:record-tail} and bounded
weights. Their total contribution is
\(NL^{O(1)}e^{-cL^{1/10}}\). There are at most \(E^2\) retained
record pairs. A shared bank prime makes a stratum empty, since every
bank prime avoids \(N\). For coprime \(r,s\), their divisibilities
select one progression of length \(|J|/(ABrs)\).

At \(p^e\Vert r\), write \(x=r\ell\). Exact valuation excludes
\(\ell\equiv0\pmod p\), while \(y\) is automatically a unit because
\(p\nmid N\). The same argument applies to \(s\). At an unused bank
prime there are two distinct excluded roots. Thus the complete bank
factor, including the progression density, is
\[
 \frac1{rs}\prod_{\substack{p\in\mathcal B\\p\nmid rs}}(1-2/p)
             \prod_{p\mid rs}(1-1/p)
       =\mu_{\mathcal B}(r,s).
\]
All forbidden-prime slopes remain units. Apply the upper Selberg
sieve with \(z=N^{1/10}\) and primes up to \(w=z^{1/16}\),
which contains every bank for sufficiently large \(N\).
The nonbank sieve denominators are bounded below, up to fixed positive
constants, by
\[
\begin{array}{c|c}
 \text{weights}&\text{sieve denominator}\\ \hline
 f_1,f_2&(\log w)^\kappa/J_{\mathcal E}(N),\\
 \mathcal D_1,f_2&
   (\log y)^{1-\rho_1}(\log w)^{1-\rho_2}/J_{\mathcal E}(N),\\
 \mathcal D_1,\mathcal D_2&
   (\log y)^\kappa/J_{\mathcal E}(N).
\end{array}
\]
The middle row also applies with the sides interchanged.
A model excludes its forbidden primes up to \(y\); an indicator
excludes them up to \(w\). A prime forbidden on only one side enters
only that side's denominator. Target-root mergers cost at most the
displayed factor \(J_{\mathcal E}(N)\). Each model bound contributes
\(L^{-3(1-\rho_i)/5}\), so all three rows give the same sparse
factor \(L^{-\kappa}J_{\mathcal E}(N)\).

Consequently every retained stratum, for \(w_i\in\{f_i,\mathcal D_i\}\),
has the single uniform estimate
\[
 \sum_{\substack{Ax+By=N,\ Ax\in J\\
                  x_{\mathcal B}=r,\ y_{\mathcal B}=s}}
          w_1(x)w_2(y)
 \le C_K\frac{|J|}{AB}\mu_{\mathcal B}(r,s)
             \frac{J_{\mathcal E}(N)}{L^\kappa}
        +O_K(z^{4+\epsilon}),
\]
where \(\epsilon>0\) is arbitrarily small and fixed.
The remainder follows from the squared-weight CRT root counts.
Multiplying by \(F(r,s)\le L^B\) and summing at most \(E^2\)
retained records leaves
\[
 L^BE^2z^{4+\epsilon}=O(N^{1/2})
\]
after fixing, for example, \(\epsilon=1/2\).
Summing the main terms gives \(\mathbb E_{\mu_{\mathcal B}}F\).
The endpoint exponent and sieve level are independent of the final
model depth; constants and thresholds may depend on it.
\end{proof}

\begin{proof}[Proof of Theorem~\ref{thm:general}]
The preceding checks verify (C1)--(C4), with
\(\mathcal S_N=NJ_{\mathcal E}(N)/(AB L^\kappa)\).
This scale is between \(NL^{-C}\) and \(NL^C\) for a fixed \(C\),
by \(J_{\mathcal E}\ge1\) and Mertens' upper product bound.
Theorem~\ref{thm:model-comparison} replaces both indicators by their
models with error \(O(\mathcal S_NL^{-\delta_0})\), uniformly in
the interval. Proposition~\ref{prop:model-correlation} evaluates that
model count with exactly the coefficient in the theorem.
If \(\Gamma_N=0\), the residue restrictions give zero identically.
Otherwise its finitely many positive values have a fixed positive
minimum, yielding the relative assertion.
\end{proof}


\subsection{Self-sums with prime support in a fixed coset}\label{sec:cosets}

The comparison theorem also covers a prime coset which contains no
identity class. The additional arithmetic issue is that its character
twists can have negative or genuinely complex singular exponents.
Those branches enter the model used for centering, even though they
are smaller than its positive leading part.

Fix an even integer \(m\), put \(G=(\Z/m\Z)^\times\), and let
\(K\le G\), \(a\in G\). Write
\[
 \mathscr A=aK,\qquad H=\langle K,a\rangle,\qquad
 r=[G:K],\qquad d=\operatorname{ord}_{G/K}(aK),\qquad \rho=1/r.
\]
Define
\[
 \mathcal P_{\mathscr A}=\{p:p\nmid m,\ p\bmod m\in\mathscr A\},\qquad
 C_{\mathscr A}=\langle\mathcal P_{\mathscr A}\rangle,\qquad f=\1_{C_{\mathscr A}}.
\]
Thus \(1\in C_{\mathscr A}\). The generated residue group is \(H\), whereas
the prime density is \(\rho\); these are different data. Put
\begin{align}
 c_{\mathscr A}&=\frac1{\Gamma(\rho)}\lim_{s\downarrow1}(s-1)^\rho
                        \prod_{p\in\mathcal P_{\mathscr A}}(1-p^{-s})^{-1},
                         \label{eq:coset-marginal}\\
 \Gamma_{\mathscr A}(N)&=\frac{m}{|H|^2}
       \#\{(u,v)\in H^2:u+v\equiv N\pmod m\},\notag\\
 \mathcal E_{\mathscr A}&=\{p:p\nmid m,\ p\bmod m\notin\mathscr A\},\notag\\
 C_{\mathcal E_{\mathscr A}}&=\prod_{p\in\mathcal E_{\mathscr A}}
                 \left(1-\frac1{(p-1)^2}\right),\qquad
 J_{\mathscr A}(N)=\prod_{\substack{p\mid N\\p\in\mathcal E_{\mathscr A}}}
                                      \frac{p-1}{p-2}.\notag
\end{align}
The constant \(c_{\mathscr A}\) is positive; an explicit expression is given below.
For \(d>1\), set
\[
 g_{\mathscr A}=\rho\bigl(1-\cos(2\pi/d)\bigr),\qquad
 \delta_{\mathscr A}=\min\left\{\frac1{8192\varphi(m)},\frac{3g_{\mathscr A}}{10}\right\}.
\]
For \(d=1\), use \(g_{\mathscr A}=+\infty\) and
\(\delta_{\mathscr A}=1/(8192\varphi(m))\).

\begin{theorem}\label{thm:coset}
Under the analytic inputs of Section~\ref{sec:inputs}, uniformly in
all sufficiently large targets with \(\Gamma_{\mathscr A}(N)>0\),
\[
 \#\{(x,y)\in C_{\mathscr A}^2:x+y=N\}
 =c_{\mathscr A}^2\Gamma_{\mathscr A}(N)C_{\mathcal E_{\mathscr A}}J_{\mathscr A}(N)
       \frac{N}{(\log N)^{2-2\rho}}
       \left(1+O((\log N)^{-\delta_{\mathscr A}})\right).
\]
If \(\Gamma_{\mathscr A}(N)=0\), the count is identically zero.
The condition \(\Gamma_{\mathscr A}(N)>0\) is equivalent to solvability in
every finite quotient. Consequently it characterizes eventual
representability. The modulus and coset remain fixed, and the error
is uniform in all target prime valuations.
\end{theorem}

The exponent incorporates both discarded banks and the gap between
the leading and secondary character branches. In particular a fixed
positive exponent works for every fixed coset. For \(d\ge2\),
\(1-\cos(2\pi/d)=2\sin^2(\pi/d)\ge8/d^2\), so one may always
replace \(\delta_{\mathscr A}\) by the simpler, weaker
\(1/(8192\varphi(m)^3)\). There is no uniformity in growing \(m\).
The new assertion concerns unfiltered self-sums. A filter on the
whole numerical value need not be preserved by multiplication by a
prime in a nonidentity coset.

\subsubsection{Every fixed-character branch}

For a character \(\psi\), define
\(F_\psi(s)=\sum_n\psi(n)f(n)n^{-s}\), with all characters inflated
to a common modulus containing \(m\). Fourier expansion of the prime
coset gives
\begin{equation}\label{eq:coset-character-factorization}
 F_\psi(s)=G_\psi(s)
        \prod_{\chi\in K^\perp}L(s,\psi\chi)^{\rho\overline{\chi(a)}}.
\end{equation}
The correction is uniformly absolutely convergent for \(\Re s>1/2\).
Indeed the coefficient of \(p^{-s}\) in the logarithm is
\[
 \rho\psi(p)\sum_{\chi\in K^\perp}\overline{\chi(a)}\chi(p)
                =\psi(p)\1_{p\bmod m\in\mathscr A};
\]
the remaining logarithmic terms are bounded by a fixed multiple of
\(\sum_p p^{-2\Re s}\). Literal Euler factors include all conductor
deletions. Analytic logarithms continued from the Euler half-plane
specify all complex powers.

A zeta factor occurs precisely for a fixed character
\(\psi\in K^\perp\); its exponent is
\begin{equation}\label{eq:coset-exponents}
                       \alpha_\psi=\rho\psi(a).
\end{equation}
All other twists have no singular branch at one. The exponents
\(\rho\) occur exactly for \(\psi\in H^\perp\). The remaining
exponents are \(\rho\) times nontrivial \(d\)-th roots of unity,
each with the same multiplicity. In particular,
\(\rho-\Re\alpha_\psi\ge g_{\mathscr A}\) outside \(H^\perp\).

For the untwisted series, an explicit positive correction is
\[
 \log G_{1,p}(s)=\sum_{j\ge2}
 \frac{\1_{p\bmod m\in\mathscr A}-\1_{(p\bmod m)^j\in\mathscr A}}{j p^{js}}
 \qquad(p\nmid m).
\]
Thus
\begin{equation}\label{eq:coset-constant}
 c_{\mathscr A}=\frac{(\varphi(m)/m)^\rho G_1(1)}{\Gamma(\rho)}
        \prod_{\substack{\chi\in K^\perp\\\chi\ne1}}
                     L(1,\chi)^{\rho\overline{\chi(a)}}.
\end{equation}
Here \(G_1(1)\) is an absolutely convergent positive Euler product.
Conjugate characters are paired. Continuation along \(s>1\) makes
the expression positive and unambiguous. The ordinary one-variable
asymptotic, and its complex twisted versions, also follow from
\cite[Theorem 1.2 and (1.9)--(1.10)]{delabretechetenenbaum}.
That fixed-function theorem is not used in place of the moving-cutoff
estimate which follows.

\subsubsection{The arithmetic model and its one-variable means}

If \(K=G\), then \(f\) is the indicator of integers coprime to
\(m\), and we use the exact model \(\mathcal D=f\). Counting the finite
residue classes on the line proves Theorem~\ref{thm:coset} with
\(O_m(1)\) error. In the model construction below assume
\(K\ne G\), so \(\rho<1\) and the model upper bound tends to zero.

Keep an outer scale \(X\), with \(L=\log X\), and let
\[
 y=e^{L^{2/5}},\quad \mathcal P_y=\{p:p\le y\},\quad
 V_y=\prod_{p\le y}(1-1/p),\quad E=e^{\sqrt L}.
\]
For \(n=A v\), where \(A=n_{\mathcal P_y}\), the integer \(v\) is
\(\mathcal P_y\)-rough. The outer scale and cutoff are held fixed
under multiplication by bank primes.
For \(\psi\in K^\perp\), write
\[
 F_{\psi,R}(s)=F_\psi(s)
        \prod_{\substack{p\le y\\p\in\mathcal P_{\mathscr A}}}
                         (1-\psi(p)p^{-s})
          =(s-1)^{-\alpha_\psi}H_{\psi,R}(s).
\]
Since \(\psi(p)=\psi(a)\) on the allowed prime set, fixed-character
Mertens and a disk of radius \(c/\log y\) give
\begin{equation}\label{eq:coset-rough-regular}
 |H_{\psi,R}(1)|\asymp(\log y)^{-\Re\alpha_\psi},\qquad
 \frac{|H_{\psi,R}^{(j)}(1)|}{j!}
       \ll_j|H_{\psi,R}(1)|(C\log y)^j.
\end{equation}
For the derivative bound, the logarithm of each finite Euler product
relative to its value at one is
\(O(|s-1|\sum_{p\le y}(\log p)/p)=O(1)\) on that disk;
the fixed regular factor is analytic and nonzero there.

Choose a sufficiently large fixed depth \(J\), and put
\begin{equation}\label{eq:coset-hankel}
 M_\psi(t)=\sum_{j=0}^J\frac{H_{\psi,R}^{(j)}(1)}{j!}
              \frac{(\log t)^{\alpha_\psi-j-1}}
                   {\Gamma(\alpha_\psi-j)}.
\end{equation}
Write \(M_0=M_1\) for the positive principal branch. On any fixed
power range, \(M_0(t)>0\), and
\begin{equation}\label{eq:coset-secondary-gap}
 |M_\psi(t)|\ll L^{-3g_{\mathscr A}/5}M_0(t)
                     \qquad(\psi\notin H^\perp).
\end{equation}
In fact its individual exponent is
\(3(\rho-\Re\alpha_\psi)/5\): the factors \(L\) and \(\log y\)
in \eqref{eq:coset-rough-regular} and \eqref{eq:coset-hankel}
give exactly this ratio. Conjugate branches are conjugate functions.

Every \(\eta\in H^\perp\) is one on the support of \(f\), so
\(F_{\psi\eta,R}=F_{\psi,R}\) and \(M_{\psi\eta}=M_\psi\)
exactly. Define
\begin{equation}\label{eq:coset-model}
 \mathcal D_J(n)=\frac{f(A)}{V_y}
              \sum_{\psi\in K^\perp}\overline{\psi(v)}M_\psi(v),
 \qquad \xi_J(n)=f(n)-\mathcal D_J(n).
\end{equation}
Grouping characters modulo \(H^\perp\) factors the bracket by
\(|G/H|\1_{v\bmod m\in H}\). Its leading term is positive and
all other terms satisfy \eqref{eq:coset-secondary-gap}. Therefore
\(\mathcal D_J\) is real and nonnegative for large \(X\), and
\begin{equation}\label{eq:coset-model-domination}
 \mathcal D_{0,J}(n)=|G/H|\1_{v\bmod m\in H}\frac{f(A)M_0(v)}{V_y},
 \qquad |\mathcal D_J-\mathcal D_{0,J}|\ll L^{-3g_{\mathscr A}/5}\mathcal D_{0,J}.
\end{equation}
For \(d=1\), the difference is zero.

For a precise extension, use the cutoff from
Section~\ref{sec:model-preparation}, immediately after
\eqref{eq:hankel-model}, with the present outer scale \(X\):
set every \(M_\psi(t)=0\) for \(t\le X^{1/4}\), and multiply
\eqref{eq:coset-hankel} by \(\eta_0(\log t/L)\) above that point.
It equals the displayed formula for \(t\ge X^{1/3}\).
The common real cutoff preserves conjugacy,
\(M_{\psi\eta}=M_\psi\), positivity of the combined model,
and \eqref{eq:coset-secondary-gap} on the transition range.
Condition (C3) only uses \(t\gg X^{1/2}/E>X^{1/3}\).
Consequently, on
\(1\le n\le X^3\),
\[
 0\le \mathcal D_J(n)\ll L^{-3(1-\rho)/5},\qquad |\xi_J(n)|\ll1.
\]
These exponents are independent of the fixed depth. For an allowed
\(p\le y\), multiplication changes \(A\) to \(pA\) and leaves
\(v\) unchanged. Hence, exactly,
\begin{equation}\label{eq:coset-invariance}
 f(pn)=f(n),\qquad \mathcal D_J(pn)=\mathcal D_J(n),\qquad \xi_J(pn)=\xi_J(n).
\end{equation}
The characters in \eqref{eq:coset-model} are evaluated at the rough
part \(v\), rather than at the whole integer \(n\).

\begin{lemma}\label{lem:coset-centering}
For every fixed \(B,A_0>0\), a sufficiently large fixed depth \(J\)
gives the recordwise estimate \eqref{eq:comparison-rough-means} for
\(\xi_J\). It holds uniformly for every \(\mathcal P_y\)-supported
record \(r_0\le E\), every \(Y\in[X^{1/2},X^2]\), bounded periodic
tests of period at most \(L^B\), and the smooth weights permitted
there. In particular it includes every fixed-character contribution.
\end{lemma}
\begin{proof}
First consider an \(\mathcal P_y\)-rough variable at scale
\(T\in[X^{2/5},X^2]\). Expand a periodic test in characters after
refining its period to \(Q=\operatorname{lcm}(q,m)\). All primes
dividing \(Q\) lie below \(y\). For a character \(\vartheta\bmod Q\),
\eqref{eq:coset-character-factorization} holds literally, with every
small allowed Euler factor removed. A principal constituent occurs
only if \(\vartheta\) is inflated from \(K^\perp\). In that case
the rough series is precisely \(F_{\vartheta,R}\) defined above;
it does not acquire a new conductor-dependent regular part.

Complex powers require a bound on the analytic logarithm, rather
than merely on the modulus of an \(L\)-function. For a principal
character put
\(L^\#(s)=(s-1)L(s)/(s+1)\), and for a nonprincipal character put
\(L^\#=L\). The zero-free input makes \(L^\#\) analytic and
nonzero on \(\Re s>11/12\). Its logarithm is fixed by the Euler
logarithm on \(\Re s=2\), together with the displayed pole-removing
factor. On a fixed disk centered at \(2+it\), contained in that
half-plane, convexity gives
\(\Re\log L^\#\ll\log(Q(2+|t|))\), while the center logarithm is
bounded by the same quantity. Borel--Carath\'eodory on a smaller
fixed disk therefore gives
\begin{equation}\label{eq:coset-analytic-log}
 |\log L^\#(\sigma+it)|\ll\log(Q(2+|t|))
                  \qquad(23/24\le\sigma\le2).
\end{equation}
The disks can have fixed radii because their left boundary remains
strictly to the right of \(11/12\). The consistent center logarithm
is uniformly controlled by its Euler series. Thus every fixed
complex power in \eqref{eq:coset-character-factorization} has
polynomial growth in conductor and height; its imaginary exponent
causes no uncontrolled argument factor.

Use smooth Mellin inversion to height \(\exp(L^{1/10})\), and move
the line to \(1-c/\log y\). This lies inside the preceding strip
for large \(X\). The finite Euler factors cost only \(L^{O(1)}\),
and \eqref{eq:coset-analytic-log} costs \(\exp(O(L^{1/10}))\).
The factor \(T^{s-1}\) supplies \(\exp(-cL^{3/5})\), which dominates
both. Sufficiently many fixed derivatives of the smooth weight
dispose of the Mellin tails. A twist without a principal constituent
therefore has mean \(O_A(TL^{-A})\) for every fixed \(A\).

For \(\vartheta\in K^\perp\), Hankel integration around one gives
exactly the density \(M_\vartheta\) in \eqref{eq:coset-hankel}.
The derivative bound \eqref{eq:coset-rough-regular} gives relative
Taylor remainder \(O_J(L^{-3(J+1)/5})\) for that branch. All branch
magnitudes are at most a fixed multiple of the principal magnitude.
This argument applies without change to negative and complex
\(\alpha_\vartheta\); the reciprocal Gamma factor is the Hankel
integral of \((s-1)^{-\alpha_\vartheta+j}\).

On the model side the unrestricted rough population has series
\(\zeta(s)\prod_{p\le y}(1-p^{-s})\), principal density \(V_y\),
and no other principal branch. Its twist by
\(\vartheta\overline\psi\) is principal exactly when
\(\vartheta=\psi\) on the units modulo \(Q\). Hence
\(V_y^{-1}\sum_{\psi\in K^\perp}\overline\psi(v)M_\psi(v)\)
matches each of the original character means with its exact
coefficient. This proves arbitrary logarithmic centering after
choosing \(J\) beyond the periodic-character and smooth-derivative
costs. A factor \(v^{iu}\), \(|u|\le L^B\), only enlarges those
fixed costs.

Finally, for a fixed record \(r_0\le E\), multiplicativity gives
\(f(r_0v)=f(r_0)f(v)\), and \eqref{eq:coset-model} has the same
factor \(f(r_0)\). The remaining scale is \(T=Y/r_0\), within
the preceding power range for large \(X\). If the record is
forbidden both terms vanish. Otherwise the just-proved estimate
is \(O((Y/r_0)L^{-A_0})\), uniformly in the record and test.
This is the required recordwise assertion, with the depth chosen
after the requested accuracy.
\end{proof}

\subsubsection{Common banks and the independent sieve check}

Choose one residue \(a_0\in aK\) and take all small and big banks
from primes \(p\equiv a_0\pmod m\), deleting \(p\mid N\).
Their fixed density is \(1/\varphi(m)\), irrespective of \(\rho\).
Equation \eqref{eq:coset-invariance} verifies condition (C2) of
Theorem~\ref{thm:model-comparison}; the bounds above verify (C1),
and Lemma~\ref{lem:coset-centering} verifies (C3).

Apply Lemma~\ref{lem:gr-ideal} with \(L_0=m\) and \(c=a_0\).
Its quotient character frame already retains the common unit phase
on products with the same slot type. Thus these banks require no
new graph construction or mark-length loss.

We verify the remaining envelope assumption independently of the
comparison. Put \(\kappa_{\mathscr A}=2-2\rho\) and
\(\Sigma_N=NJ_{\mathscr A}(N)L^{-\kappa_{\mathscr A}}\). This satisfies
\eqref{eq:comparison-scale}, for example with \(\kappa_0=3\), since
\(0\le\kappa_{\mathscr A}<2\), \(J_{\mathscr A}(N)\ge1\), and the prime-product bound
\(J_{\mathscr A}(N)\ll(\log\log(3N))^2\) is uniform in \(N\).
At a bank prime \(p\nmid N\),
the exact two-coordinate valuation law is \eqref{eq:bank-law}.
For coprime bank records \(r_1,r_2\le E\), their divisibilities
select one interval progression of length
\(|\mathcal I|/(r_1r_2)\) on \(x+y=N\).
Exactness excludes two roots at every unused bank prime, and one
remaining root at a prime dividing \(r_1r_2\). Together with
\(1/(r_1r_2)\), these factors are exactly its probability under
the independent bank law. Shared bank-prime divisibility is impossible.

Use the upper Selberg sieve at \(z=N^{1/10}\), with primes up to
\(w=z^{1/16}\). This contains all banks. Every good forbidden prime
has two roots, merging precisely when it divides \(N\). Fixed-coset
prime Mertens gives the following lower bounds for the nonbank
sieve denominator, up to fixed positive constants:
\[
\begin{array}{c|c}
 \text{weights}&\text{sieve denominator}\\ \hline
 f,f&(\log w)^{2(1-\rho)}/J_{\mathscr A}(N),\\
 \mathcal D_J,f&(\log y)^{1-\rho}(\log w)^{1-\rho}/J_{\mathscr A}(N),\\
 \mathcal D_J,\mathcal D_J&(\log y)^{2(1-\rho)}/J_{\mathscr A}(N).
\end{array}
\]
For a model factor use its nonnegative domination by \(2D_{0,J}\),
its bound \(L^{-3(1-\rho)/5}\), and its exclusion of every forbidden
prime at most \(y\). These bounds turn each row into exactly
\(L^{-\kappa_{\mathscr A}}J_{\mathscr A}(N)\). Fixed residue conditions can be discarded
for this upper estimate. Its CRT endpoint error per retained record
pair is \(O(z^{4+\epsilon})\); summing at most \(E^2\) pairs gives
\(N^{0.45+o(1)}\). The exponent is independent of the model depth.

Rankin's bound with parameter \(c/\log y\) gives
\[
 \#\{n\le N:n_{\mathcal P_y}>E\}
                 \ll NL^{O(1)}e^{-cL^{1/10}}.
\]
It applies to every actual/model pair because the weights are
bounded. Multiplying the record-stratum estimate by an arbitrary
nonnegative record function bounded by \(L^B\), and summing, proves
\eqref{eq:comparison-envelope}, including its interval-length factor
and its errors. This is condition (C4). Its fixed logarithmic error
exponents do not depend on the eventual density depth.

All four hypotheses of Theorem~\ref{thm:model-comparison} have now
been checked by invariance, one-variable estimates, or a positive
upper sieve. No additive cancellation or residue-energy conclusion
has been included as an arithmetic assumption.

\subsection{The main term and the exact congruence criterion}

\begin{lemma}\label{lem:coset-finite-images}
A residue \(t\bmod M\) is represented by \(C_{\mathscr A}\) if and only if
its reduction modulo \(\gcd(M,m)\) is that of some \(h\in H\),
and \(p\nmid t\) at each forbidden prime \(p\mid M\), \(p\nmid m\).
\end{lemma}
\begin{proof}
Necessity follows from prime support. For sufficiency extract an
allowed-prime prefix \(D\) with valuations
\(\min(v_p(t),v_p(M))\) at allowed primes dividing \(M\).
It belongs to \(C_{\mathscr A}\), and its residue belongs to \(H\).
The desired unit tail has state \(hD^{-1}\in H\).

The positive monoid generated by \(\mathscr A\) in the finite group is \(H\),
so choose a nonempty word \(s_1\cdots s_k=hD^{-1}\), with
\(s_i\in\mathscr A\). An identity state also has a nonempty word, by taking
a positive power of any element of \(\mathscr A\). Realize all but the last
letter by primes in their reduced classes modulo \(m\), avoiding
\(M\). For the last letter, impose its class modulo \(m\), and
the remaining unit-tail congruences modulo the prime powers of \(M\).
At an allowed prime with prefix valuation \(e<v_p(M)\), this is
the quotient of \(t/p^e\) by the other unit factors modulo
\(p^{v_p(M)-e}\). At a saturated prefix no further congruence is
needed, beyond avoiding that prime. All other good primes require
a unit tail by the hypothesis. At primes dividing \(m\), compatibility
is precisely the prescribed reduction of \(h\).
The resulting CRT class is reduced, and Dirichlet's theorem supplies
a prime in it. Multiplying all chosen primes and \(D\) realizes
the required residue. A single prime would not suffice in general:
the tail state need not itself lie in \(\mathscr A\).
\end{proof}

\begin{proof}[Proof of Theorem~\ref{thm:coset}]
The full-group case was handled by its exact model. Suppose \(K\ne G\).
Theorem~\ref{thm:model-comparison}, applied using the checks above,
replaces both copies of \(f\) by \(\mathcal D_J\), at the sparse scale
\(NJ_{\mathscr A}(N)L^{-\kappa_{\mathscr A}}\), with relative error
\(O(L^{-1/(8192\varphi(m))})\). Its interval version and the
upper sieve handle the endpoint strips.

On the retained records and power bulk,
\eqref{eq:coset-model-domination} replaces \(\mathcal D_J\mathcal D_J\) by
\(\mathcal D_{0,J}\mathcal D_{0,J}\) with relative error \(O(L^{-3g_{\mathscr A}/5})\).
Put
\[
 \alpha_0=\varphi(m)/m,\qquad
 E_-(y)=\prod_{\substack{p\le y\\p\nmid m,\ p\notin\mathcal P_{\mathscr A}}}
                        (1-p^{-1}).
\]
Since \(\log v=L+O(\sqrt L+\log L)\), the leading model is
\[
 |G/H|\1_{n\bmod m\in H}\,f(n_{\mathcal P_y})\lambda_X,
 \qquad \lambda_X=\frac{c_{\mathscr A}L^{\rho-1}}{\alpha_0E_-(y)},
\]
with relative error \(O(L^{-1/2})\). Here
\(n_{\mathcal P_y}\bmod m\in H\), so the rough and whole-value
state tests are equivalent for this leading model.

Let \(R_N\) be the number of pairs counted in \(\Gamma_{\mathscr A}(N)\).
The raw fixed-residue probability on the integer line is \(R_N/m\).
Its two projections and marginal normalization give
\[
 \frac{R_N}{m}\frac{|G/H|^2}{\alpha_0^2}
              =\frac{mR_N}{|H|^2}=\Gamma_{\mathscr A}(N).
\]
This calculation uses \(H\) for the residue support and \(\rho\)
for the logarithmic density; it never identifies \(\rho\) with
\(|H|/|G|\).
At a good forbidden prime the normalized pair factor is
\[
 \frac{1-(2-\1_{p\mid N})/p}{(1-1/p)^2}.
\]
The fundamental sieve lemma at a fixed power level, as in
\cite[Theorem 11.7]{richert}, therefore evaluates the leading model
as
\[
 c_{\mathscr A}^2\Gamma_{\mathscr A}(N)C_{\mathcal E_{\mathscr A}}J_{\mathscr A}(N)
                       \frac{N}{L^{\kappa_{\mathscr A}}}.
\]
Its sieve error is smaller than every fixed logarithmic power;
CRT endpoints are \(N^{0.45+o(1)}\). The omitted ordinary Euler
tail is \(O(1/y)\), and its target-prime tail is
\(O(L/(y\log y))\), uniformly in all target valuations.
The density replacement and secondary-branch errors are bounded
by the asserted \(L^{-\delta_{\mathscr A}}\). This proves the asymptotic.

For the local assertion, testing modulo \(m\) proves necessity of
\(\Gamma_{\mathscr A}(N)>0\). Conversely fix such a root. At primes dividing
\(m\), lift the equation while preserving its unit states.
At every other prime, both coefficients are units and \(p\ge3\).
The equation modulo \(p\) has \(p\) pairs; requiring a unit on
each forbidden coordinate removes at most two. A surviving pair
lifts, and CRT combines the choices. Lemma~\ref{lem:coset-finite-images}
realizes each coordinate by an actual element of \(C_{\mathscr A}\).
This proves solvability in every finite quotient. Finally the
positive asymptotic proves eventual global representability.
\end{proof}

\subsubsection{Two cosets which test the branches}

For \(m=4\), \(K=\{1\}\), \(a=3\), the generated group is
\(H=G\), while \(\rho=1/2\). The exponents are \(1/2,-1/2\).
Let \(C^-_4\) contain the positive integers supported on primes
\(3\pmod4\). With
\(P_3=\prod_{p\equiv3(4)}(1-p^{-2})^{-1}\), its marginal is
\(c_{\mathscr A}^2=2P_3/\pi^2\). Indeed
\[
 F(s)^2=(1-2^{-s})\frac{\zeta(s)}{L(s,\chi_{-4})}F(2s).
\]
Here \(\Gamma_{\mathscr A}(N)=2\) on every even target. Since
\(\prod_{p\text{ odd}}(1-p^{-2})^{-1}=\pi^2/8\),
\[
 2c_{\mathscr A}^2C_{\mathcal E_{\mathscr A}}
       =\frac12\prod_{p\equiv1(4)}
                    \left(1-\frac2{p(p-1)}\right)=:A_-.
\]
Thus Theorem~\ref{thm:coset} gives
\[
 \#\{x+y=N:x,y\in C^-_4\}
 = A_-\prod_{\substack{p\mid N\\p\equiv1(4)}}\frac{p-1}{p-2}
       \frac N{\log N}
       \left(1+O((\log N)^{-1/16384})\right)
\]
for every sufficiently large even \(N\). Odd targets have zero
count. The negative branch is needed for centering, although it
does not change this leading constant.

For a genuinely complex example, take \(m=10\), \(K=\{1\}\),
\(a=3\), and let \(C_{3,10}\) be supported on primes
\(3\pmod{10}\). Choose the literal quartic character \(\chi\bmod10\)
with \(\chi(3)=i\). Then \(H=G\), \(\rho=1/4\), \(d=4\), and
the exponents of \(1,\chi,\chi^2,\overline\chi\) are
\[
                       \frac14,\quad\frac i4,\quad-\frac14,
                                      \quad-\frac i4.
\]
The rough model is explicitly real:
\[
 \mathcal D_J(n)=\frac{f(A)}{V_y}
   \left\{M_1(v)+2\Re\bigl(\overline{\chi(v)}M_\chi(v)\bigr)
                       +\chi(v)^2 M_{\chi^2}(v)\right\}.
\]
The complex terms have relative size \(O(L^{-3/20})\), and the
negative branch has relative size \(O(L^{-3/10})\). They cannot
be omitted from an arbitrarily accurate centered model.
The leading constant is given by \eqref{eq:coset-marginal}; equivalently,
\[
 c_{3,10}=\frac{(2/5)^{1/4}G_1(1)}{\Gamma(1/4)}
       L(1,\chi)^{-i/4}L(1,\chi^2)^{-1/4}
                              L(1,\overline\chi)^{i/4},
\]
with the positive determination and convergent correction specified
in \eqref{eq:coset-constant}. The forbidden good residues are
\(1,7,9\pmod{10}\). For even \(N\),
\[
 \Gamma_{\mathscr A}(N)=
 \begin{cases}5/2,&5\mid N,\\15/8,&5\nmid N.\end{cases}
\]
This follows by counting pairs of nonzero residues modulo five:
there are four at a zero target and three otherwise. Consequently
every even target is congruence-admissible, and
\[
 \#\{x+y=N:x,y\in C_{3,10}\}
 =c_{3,10}^2\Gamma_{\mathscr A}(N)C_{\mathcal E_{\mathscr A}}J_{\mathscr A}(N)
       \frac N{(\log N)^{3/2}}
       \left(1+O((\log N)^{-1/32768})\right).
\]
This includes all sufficiently large even targets, uniformly in
their prime valuations. The example tests complex character
contributions without changing the comparison theorem or its inputs.

The common-bank condition is substantive. Two different prime
prescriptions can be disjoint and supply no common allowed prime
class. Theorem~\ref{thm:coset} makes no assertion about arbitrary
pairs of prime cosets.

\section{Primitive forms and invariance within a genus}\label{sec:forms}

Fix a negative fundamental discriminant \(D\), put \(K=\Q(\sqrt D)\),
and let
\[
 \mathcal O_i=\Z+f_i\mathcal O_K,\qquad
 \Delta_i=f_i^2D,\qquad G_i=\operatorname{Pic}(\mathcal O_i)
 \quad(i=1,2).
\]
Let \(Q_i\) be a primitive positive definite integral binary form of
discriminant \(\Delta_i\). For any such form \(Q\), define
\[
 \begin{aligned}
 P_Q&=\{Q(x,y):x,y\in\Z,\ \gcd(x,y)=1\},\\
 \mathscr V_{p,Q}&=Q(\Z_p^2\setminus p\Z_p^2),\qquad
 \Lambda_Q=\{n\ge1:n\in\mathscr V_{p,Q}\text{ for every }p\}.
 \end{aligned}
\]
Only positive values are included. Coordinates may vanish, and each
numerical value is counted once. Write
\[
 R_{T_1,T_2}(N)=\#\{(a,b)\in T_1\times T_2:a+b=N\},\qquad
 R_{Q_1,Q_2}=R_{P_{Q_1},P_{Q_2}}.
\]
Two forms are in the same genus when they are equivalent over every
\(\Z_p\) and over \(\mathbb R\). These equivalences preserve primitive
vectors and therefore the primitive local value sets.

Put \(\chi=\chi_D\) and
\[
 S=\{p:p\mid2\Delta_1\Delta_2\},\qquad
 J_S(N)=\prod_{\substack{p\mid N\\p\notin S,\ \chi(p)=-1}}
                              \frac{p-1}{p-2}.
\]
A target is locally admissible for \(Q_1,Q_2\) if
\(N\in\mathscr V_{p,Q_1}+\mathscr V_{p,Q_2}\) for every prime \(p\).
For a finite abelian group \(G\), put
\[
 \eta_*(G)=
 \begin{cases}
 +\infty,&2G=0,\\
 \displaystyle\min_{\substack{\ell\mid |2G|\\\ell\ {\rm prime}}}
                   \frac{\ell-1}{2\ell},&2G\ne0.
 \end{cases}
\]
The mixed local value densities and universal factor are
\[
 d_{p;1,2}(N)=\int_{\Z_p}
       \1_{\mathscr V_{p,Q_1}}(a)\1_{\mathscr V_{p,Q_2}}(N-a)\,da
       \quad(p\in S),
\]
\begin{equation}\label{eq:form-universal-factor}
 \mathfrak b_{D,S}=\frac{4L(1,\chi)}\pi
  \prod_{p\in S}\frac{1-\chi(p)/p}{1-1/p}
  \prod_{\substack{p\notin S\\\chi(p)=-1}}
                      \left(1-\frac2{p(p-1)}\right).
\end{equation}
Haar measure gives \(\Z_p\) mass one.

\begin{theorem}[Common-field primitive form values]\label{cor:common-field-form-pairs}
For the fixed data above, the following assertions hold.
\begin{enumerate}
\item\label{part:form-classical}\label{thm:genus-binary-invariance}
\emph{Classical order bounds and genus comparison.}
Uniformly on sufficiently large locally admissible targets,
\[
 R_{Q_1,Q_2}(N)\asymp R_{\Lambda_{Q_1},\Lambda_{Q_2}}(N)
                              \asymp\frac{NJ_S(N)}{\log N}.
\]
If \(Q'_i\) belongs to the genus of \(Q_i\), then, for every fixed
\(0<\eta<\min_i\eta_*(G_i)\), uniformly on their common sufficiently
large locally admissible targets,
\[
 R_{Q_1,Q_2}(N)=R_{Q'_1,Q'_2}(N)
                         \left(1+O((\log N)^{-\eta})\right).
\]
If both doubled class groups vanish, the counts agree for every target.
Otherwise \(\eta=1/5\) is always permissible. These assertions use
Appendix~\ref{sec:older-route} and the class-error sieve, independently
of the recent zero-free and weighted-dilation inputs.
\item\label{part:form-leading}\label{thm:form-binary-leading}
\emph{Explicit leading term.}
Under the analytic inputs of Section~\ref{sec:inputs}, uniformly
through sufficiently large locally admissible targets,
\begin{equation}\label{eq:form-binary-leading}
 R_{Q_1,Q_2}(N)\sim
 \mathfrak b_{D,S}\left(\prod_{p\in S}d_{p;1,2}(N)\right)
                              J_S(N)\frac N{\log N}.
\end{equation}
Each bad-prime density is a computable rational periodic function.
The leading coefficient depends only on the respective genera,
with no class-number divisor.
\end{enumerate}
Nonadmissible targets have zero count. Constants and thresholds may
depend on the fixed forms and, in the comparison, on \(\eta\).
The two conductors may differ; the quadratic field remains fixed.
\end{theorem}

A saturated numerical value can occur in every class of its compatible
genus; the classes do not partition the values. Taking \(f_1=f_2=f\)
gives the one-discriminant case, with \(\Delta=f^2D\) and \(J_\Delta=J_S\).
Write \(\mathcal F=(Q_1,Q_2)\) for the fixed pair in the proof.

\subsection{Proper ideals and the genus map}

Fix either order temporarily, write it as \(\mathcal O\), and put
\(G=\operatorname{Pic}(\mathcal O)\). The set \(S\) may contain
additional primes besides those dividing \(2\operatorname{disc}\mathcal O\).
We use proper invertible ideals of the order.  If $I$ is an invertible
fractional ideal, write $I_p=i_p\mathcal O_p$ and let $n(I)>0$ be its
covolume norm.  Local determinants give
\begin{equation}\label{eq:norm-ideal-identities}
 v_p(n(I))=v_p(\operatorname N(i_p)),\qquad
 n(IJ)=n(I)n(J),\qquad I\overline I=n(I)\mathcal O.
\end{equation}
These statements hold at conductor primes as well.  Call an integral
ideal primitive if it is not contained in $p\mathcal O$ for any rational
prime $p$.  Finitely many prescribed locally principal lattices, equal
to $\mathcal O_p$ elsewhere, glue to an invertible fractional ideal:
place them between $m\mathcal O_p$ and $m^{-1}\mathcal O_p$, glue the
residue submodules by the Chinese remainder theorem, and glue their
inverses in the same way.

The following principal-genus statement is classical for arbitrary
orders; compare \cite[Appendix~A, Proposition~A.1]{khayutin}.  We give
the norm argument because retaining conductor-local ideal types is
necessary here.

\begin{lemma}\label{lem:norm-genus}
Put $U_p=\operatorname N(\mathcal O_p^\times)$.  The homomorphism
\[
 \Theta:G\longrightarrow
       \prod_{p\in S}\Z_p^\times/U_p,\qquad
 \Theta_p([I])=n(I)/\operatorname N(i_p)\pmod{U_p}
\]
is well defined and has kernel $2G$.
\end{lemma}

\begin{proof}
The quotient in $\Theta_p$ is a unit by
\eqref{eq:norm-ideal-identities}.  A change of local generator changes
it by an order-unit norm.  Multiplication of $I$ by a principal ideal
changes numerator and denominator by the same positive element norm.
Every $U_p$ contains the rational unit squares, so the target has
exponent two and $2G$ is in the kernel.  Outside $S$ every unit is a
norm from $\mathcal O_p^\times$.

Conversely, if $\Theta([I])=1$, choose units $u_p\in\mathcal O_p^\times$
such that $\operatorname N(i_pu_p)=n(I)$ at every finite prime.
Positivity handles the infinite place.  The cyclic Hasse norm theorem
\cite[Chapter~VIII, Theorem~3.1]{milnecft} supplies $\alpha\in K^\times$
with $\operatorname N(\alpha)=n(I)$.  Local Hilbert~90 writes
$i_pu_p/\alpha=z_p/\overline z_p$.  For a norm-one element $c\ne-1$
one may take $z=1+c$; for $c=-1$, take $z=\sqrt D$.  At almost all
primes a unit solution exists: in an unramified field normalize its
valuation by a rational uniformizer, and in a split algebra use
$z=(c_1,1)$.  The lattices $z_p\mathcal O_p$ therefore glue to an
invertible fractional ideal $J$.  They give
$I/(\alpha)=J/\overline J$, hence $[I]=2[J]$ by
\eqref{eq:norm-ideal-identities}.

\end{proof}

\subsection{The target ideal class and its local shadow}

For either form \(Q(x,y)=ax^2+bxy+cy^2\), use its own order,
discriminant and class group, and take
\[
 \beta=(b+\sqrt\Delta)/2,\qquad
 A=\Z a+\Z\beta\subset\mathcal O,\qquad c_Q=-[A]\in G.
\]
Here \(\mathcal O=\Z[\beta]\), \(n(A)=a\), and
\(N(ax+\beta y)/a=Q(x,y)\). The ideal is proper invertible:
\[
 A\overline A=a(a,\beta,\overline\beta,c)\mathcal O=a\mathcal O,
\]
since \(\beta+\overline\beta=b\) and \(\gcd(a,b,c)=1\).
This construction allows \(a\) to share primes with the conductor.
It also fixes the class sign in our convention. For the classical
normalized ideal-form description compare
\cite[Lemma 10.2]{khayutin}.

\begin{lemma}\label{lem:form-ideal-shadow}
For every positive integer \(n\),
\begin{align*}
 n\in P_Q
 &\ \Longleftrightarrow\
 \text{some primitive integral invertible ideal }I
       \text{ has }n(I)=n,\ [I]=c_Q,\\
 n\in\Lambda_Q
 &\ \Longleftrightarrow\
 \text{some primitive integral invertible ideal }I
       \text{ has }n(I)=n,\ [I]\in c_Q+2G.
\end{align*}
\end{lemma}
\begin{proof}
An element \(\alpha\in A\) gives \(I=(\alpha)A^{-1}\), with
\(n(I)=N(\alpha)/a\) and class \(-[A]\). It is integral because
\((\alpha)\subset A\). Multiplication by the invertible ideal \(A\)
shows \(I\subset p\mathcal O\) if and only if \(\alpha\in pA\).
This is exactly the coordinate-primitivity condition. Conversely,
if \([I]=-[A]\), then \(IA=(\alpha)\subset A\), and the same
identities give a primitive representation of \(n(I)\).

Write \(A_p=z_p\mathcal O_p\). The norm identities in
\eqref{eq:norm-ideal-identities} give the units
\[
 \lambda_p=N(z_p)/a\in\Z_p^\times,\qquad
 \mathscr V_{p,Q}=\lambda_p\mathscr V_p,
\]
where
\(\mathscr V_p=\{N(w):w\in\mathcal O_p\setminus p\mathcal O_p\}\)
is the primitive norm-value set for \(\mathcal O_p\). The change from
an integral basis of \(A\) to the local generator is in \(\operatorname{GL}_2(\Z_p)\).
Moreover \(\Theta_p(c_Q)=\lambda_p\pmod{U_p}\), with \(\Theta\)
and \(U_p\) from Lemma~\ref{lem:norm-genus}.

If \(n\in\Lambda_Q\), choose primitive local elements
\(\alpha_p=z_pw_p\in A_p\) with \(N(\alpha_p)=an\).
The ideals \(w_p\mathcal O_p\) are primitive integral invertible,
and equal \(\mathcal O_p\) outside finitely many primes. Glue them
as in Lemma~\ref{lem:norm-genus}. Their norm valuations are
\(v_p(n)\), so the resulting ideal \(I\) has norm \(n\), and
\[
 \Theta_p([I])=n/N(w_p)=\lambda_p\pmod{U_p}.
\]
The kernel of \(\Theta\) is \(2G\), proving the coset assertion.
Conversely, for such an \(I\), write \(I_p=i_p\mathcal O_p\).
Equality of the \(\Theta\)-classes supplies a unit \(u_p\) with
\(N(u_p)=n/(\lambda_pN(i_p))\). Then
\(\alpha_p=z_pi_pu_p\) is primitive in \(A_p\) and has norm
\(an\). Outside \(S\) the order-unit norm group is all of
\(\Z_p^\times\), so the same conclusion holds there.
This proves the local-shadow equivalence at conductor primes as well.
\end{proof}

The coset condition is existential. A conductor-divisible value need
not determine a unique genus. For example, at discriminant \(-32\),
\(x^2+8y^2\) and \(3x^2+2xy+3y^2\) have primitive odd values
\(1\pmod8\) and \(3\pmod8\), respectively, and hence different
genera. Both nevertheless primitively represent \(24\), at
\((4,1)\) and \((3,-1)\).

\subsection{Orientation saturation and maximal subgroups}

For either form \(Q\), work in its own group \(G\), with target
class \(c_Q=-[A]\), and write $h_{\mathcal O}=|G|$ and $B=2G$. For every inverse pair
$\{g,-g\}\subset G$ with $2g\ne0$, let
$\Omega_g$ be the rational split primes outside $S$ whose two prime
ideal classes are $g,-g$, and put
\[
 m_g=\operatorname{ord}(2g)-1,\qquad
 k_g(n)=\#\{p\in\Omega_g:v_p(n)=1\}.
\]
These prime families are disjoint. Enlarging \(S\) deletes only
finitely many primes. The fixed-discriminant prime-class
estimate \cite[Lemma~2, equation~(5.2)]{bourgainfuchs} gives
\begin{equation}\label{eq:norm-class-mertens}
 \sum_{\substack{p\le y\\p\in\Omega_g}}\frac1p
                =\frac1h_{\mathcal O}\log\log y+O_{\mathcal O,S}(1).
\end{equation}
The coefficient is $1/h_{\mathcal O}$ for a nonambiguous inverse pair.  The source
allows nonfundamental negative form discriminants.  For a fixed order
its prime-counting error is $O_{\mathcal O}(x\exp(-c\sqrt{\log x}))$;
partial summation gives the bounded remainder in
\eqref{eq:norm-class-mertens}.

Call the pair $\{g,-g\}$ saturated for $n$ if $k_g(n)\ge m_g$, and set
\[
 B_n=\langle2g:\{g,-g\}\text{ is saturated for }n\rangle\le B.
\]
If $n\in\Lambda_Q$ and $B_n=B$, then $n\in P_Q$. To see this,
start with a primitive ideal from Lemma~\ref{lem:form-ideal-shadow},
whose class belongs to $c_Q+B$. At a good split prime, primitivity puts
the entire prime power in one conjugate prime ideal. From each saturated family
reserve $m_g$ exponent-one factors, initially orienting each as $g$.
This changes the ideal class only by an element of $B$, so its class
still belongs to $c_Q+B$. Reversing $j$ reserved orientations changes
its class by $-2jg$, for $0\le j\le m_g$; these changes fill
$\langle2g\rangle$. Families have disjoint prime support, so their
choices are independent and their sums fill $B_n=B$. They therefore
cancel its difference from \(c_Q\). The resulting ideal is primitive,
has class \(c_Q\), and has the same numerical norm and every
conductor-local ideal unchanged. Hence $n\in P_Q$ by
Lemma~\ref{lem:form-ideal-shadow}.

Suppose $B\ne0$, and let $\mathcal H$ be its finite set of maximal
proper subgroups. For $M\in\mathcal H$ define
\[
 \begin{gathered}
 \ell_M=[B:M],\qquad
 \Omega_M=\bigcup_{\{g,-g\}:\,2g\notin M}\Omega_g,\\
 K_M=\sum_{\{g,-g\}:\,2g\notin M}(m_g-1),\qquad
 k_M(n)=\#\{p\in\Omega_M:v_p(n)=1\}.
 \end{gathered}
\]
Each inverse pair is counted once. Since $B$ is finite abelian,
$B/M$ is cyclic of prime order $\ell_M$. If $n\in\Lambda_Q\setminus P_Q$,
then $B_n$ is proper and lies in some $M\in\mathcal H$. Every pair
with $2g\notin M$ is consequently deficient, giving
\begin{equation}\label{eq:norm-deficiency}
 \Lambda_Q\setminus P_Q\subset
     \bigcup_{M\in\mathcal H}\{n:k_M(n)\le K_M\}.
\end{equation}
If $2G=0$, Lemma~\ref{lem:form-ideal-shadow} already gives $P_Q=\Lambda_Q$.

The doubling map $G\to B$ has equal fibres. Precisely
$h_{\mathcal O}(\ell_M-1)/\ell_M$ classes satisfy $2g\notin M$. None is
ambiguous, since $2g\ne0$, so they form
$h_{\mathcal O}(\ell_M-1)/(2\ell_M)$ inverse pairs. Summing
\eqref{eq:norm-class-mertens} over these disjoint families proves
\begin{equation}\label{eq:norm-maximal-mertens}
 \sum_{\substack{p\le y\\p\in\Omega_M}}\frac1p
     =d_M\log\log y+O_{\mathcal O,S}(1),\qquad
 d_M=\frac{\ell_M-1}{2\ell_M}\ge\frac14.
\end{equation}
Every prime divisor of $|B|$ occurs as the index of a maximal
subgroup: in the corresponding primary component, reduce a cyclic
factor modulo that prime. Thus $\min_{M\in\mathcal H}d_M=\eta_*(G)$.
The proper-subgroup mechanism is classical in the exceptional-value
argument of Bourgain--Fuchs
\cite[Theorem 1, alternative (4.2), and Section 2, part (5)]{bourgainfuchs}.
Here its use in \eqref{eq:norm-deficiency} makes an entire positive-density
collection deficient on the binary line. For $G\simeq\Z/8\Z$ or
$\Z/2\Z\times\Z/4\Z$ the threshold is $1/4$, while for
$G\simeq(\Z/3\Z)^2$ it is $1/3$.

\begin{proposition}[The class error]\label{prop:norm-class-error}
In the common-field setting above, for every fixed
\(0<\eta<\min_i\eta_*(G_i)\), uniformly for \(N\ge3\),
\begin{equation}\label{eq:form-class-error}\csname ltx@label\endcsname{eq:norm-class-error}
 0\le R_{\Lambda_{Q_1},\Lambda_{Q_2}}(N)-R_{Q_1,Q_2}(N)
 \ll_{\mathcal F,\eta}
              \frac{NJ_S(N)}{(\log N)^{1+\eta}}.
\end{equation}
If \(2G_i=0\), then \(P_{Q_i}=\Lambda_{Q_i}\), and that summand
contributes no class error. If both doubled groups vanish the displayed
difference is zero. In every case the exponent \(\eta=1/5\) is
permissible. Only this saving exponent is independent of the orders;
the implied constant remains dependent on the fixed forms.
The proposition is classical and uses neither recent analytic input.
\end{proposition}

\subsection{A sieve for the class error}

We prove Proposition~\ref{prop:norm-class-error} directly on the line
$a+(N-a)=N$. First charge a bad pair to its first summand
\(a\in\Lambda_{Q_1}\setminus P_{Q_1}\), and use the maximal-subgroup
families in \(G_1\). If \(2G_1=0\), there is no such error.
Otherwise fix \(M\in\mathcal H\) in this group. At every good inert
prime $p$, both values must be units,
so the forbidden residues of $a$ modulo $p$ are $0,N$, with
$\nu_N(p)=2-\mathbf1_{p\mid N}$ roots.  At a marked prime
$p\in\Omega_M$, forbid exactly $v_p(a)=1$: this removes the $p-1$
residues $p,2p,\ldots,(p-1)p$ modulo $p^2$.  Its forbidden density is
\[
 f_*(p)=\frac{p-1}{p^2}.
\]
Higher prime powers remain permitted, including when $p\mid N$.

Here are the needed uniform sieve details.  Given any mark set
$\mathcal M\subset\Omega_M\cap[1,y]$, use modulus $p$ at inert primes
and $p^2$ at marked primes. All sieve indices below are squarefree
and supported on the chosen inert and marked primes at most $y$.
For such an index $d$, let
$\ell(d)$ be the product of these moduli, $\nu(d)$ the product of
the root counts, and $f(d)=\nu(d)/\ell(d)$.  Exact interval counting
by the Chinese remainder theorem gives $Xf(d)+O(\nu(d))$ on an
interval of length $X$.  In all cases $\nu(d)\le d$ and
$0<f(p)\le2/p<1$.

For $d\le z$ the Selberg quadratic majorant has main term $X/G(z)$,
where
\[
 w(d)=\prod_{p\mid d}\frac{f(p)}{1-f(p)},\qquad
 G(z)=\sum_{d\le z}\mu^2(d)w(d).
\]
To see the normalization, set $x_d=\lambda_df(d)$ and
$y_u=\sum_{u\mid d}x_d$.  Then
\[
 \sum_{d,e\le z}\lambda_d\lambda_e f([d,e])
       =\sum_{u\le z}\frac{y_u^2}{w(u)},\qquad
 \lambda_1=\sum_{u\le z}\mu(u)y_u.
\]
The minimum with $\lambda_1=1$ is $1/G(z)$, attained at
$y_u=\mu(u)w(u)/G(z)$.  Finite inversion gives
\[
 \lambda_d=\mu(d)\prod_{p\mid d}(1-f(p))^{-1}
       \frac{G_d(z/d)}{G(z)},\qquad
 G_d(v)=\sum_{\substack{e\le v\\(d,e)=1}}\mu^2(e)w(e).
\]
Uniformly in the mark set and target, $|\lambda_d|\ll_\epsilon d^\epsilon$.
Since $\nu([d,e])\le de$, the total interval error is
$O_\epsilon(z^{4+\epsilon})$.

Take $y=z^{1/16}$.  The unrestricted Euler mass for the same prime
indices is
\[
 E(y)=\prod_{\substack{p\le y\\p\notin S,\,\chi(p)=-1}}
                   (1-\nu_N(p)/p)^{-1}
       \prod_{p\in\mathcal M}(1-f_*(p))^{-1}.
\]
In its normalized positive expansion, a prime is included with
probability $f(p)$.  Therefore
$\mathbb E\log d\le2\sum_{p\le y}(\log p)/p\le4\log y$
for sufficiently large $y$.  Markov's inequality retains at least
three quarters of the mass on $d\le z$, giving $G(z)\ge3E(y)/4$.
Fixed quadratic-character prime number theory gives
\[
 \prod_{\substack{p\le y\\p\notin S,\,\chi(p)=-1}}
                 (1-\nu_N(p)/p)^{-1}
 =\frac{\displaystyle\prod_{p\le y,\,p\notin S,\,\chi(p)=-1}
                           (1-2/p)^{-1}}
        {\displaystyle\prod_{p\le y,\,p\mid N,\,p\notin S,\,\chi(p)=-1}
                           (p-1)/(p-2)}
 \gg_{\mathcal F}\frac{\log y}{J_S(N)}.
\]
Consequently, the number surviving the marked sieve is at most
\begin{equation}\label{eq:norm-marked-sieve}
 \ll_{\mathcal F}
 \frac{XJ_S(N)}{\log y}
                \prod_{p\in\mathcal M}(1-f_*(p))
                +O_\epsilon(z^{4+\epsilon}).
\end{equation}
This also holds without marked primes.  Only root counts and exact
interval errors have been used, so the estimate is uniform in target
valuations and root positions.

Fix $M\in\mathcal H$. Independently mark each prime of
$\Omega_M\cap[1,y]$ with probability $t\in(0,1)$. An integer with
$k_M(a)\le K_M$ survives with probability at least $(1-t)^{K_M}$:
its exponent-one factors below $y$ are at most its total such factors.
This lower bound is a positive constant depending on the fixed order
and $t$. On taking the finite expectation in
\eqref{eq:norm-marked-sieve}, the marked factor becomes
\[
 \prod_{\substack{p\le y\\p\in\Omega_M}}(1-tf_*(p))
                 \ll_{\mathcal F,t}(\log y)^{-td_M},
\]
by \eqref{eq:norm-maximal-mertens}, since $f_*(p)=p^{-1}+O(p^{-2})$.
Take $z=N^{1/10}$, $y=z^{1/16}$ and $\epsilon=1/2$.
The number of pairs in \(\Lambda_{Q_1}\times\Lambda_{Q_2}\) whose first summand has
$k_M(a)\le K_M$ is at most
\[
 \ll_{\mathcal F,t}
 \frac{NJ_S(N)}{(\log N)^{1+td_M}}+N^{9/20}.
\]
Given \(\eta<\eta_*(G_1)\), choose \(t\in(0,1)\) with
\(\eta<t\eta_*(G_1)\). Sum over this group's finitely many maximal
subgroups. Repeat the argument for the second summand, using only its
own group \(G_2\) and prime families. Their Mertens densities are
unchanged by deleting the finitely many extra primes in \(S\).
The common good inert sieve is the same in both arguments, while
only one summand is marked at a time. Overlap between the two orders'
split-prime families is therefore irrelevant. Taking \(\eta\) below
both thresholds and adding the two bounds proves
\eqref{eq:form-class-error}; the power error is absorbed since
\(J_S(N)\ge1\). Increasing the constant covers all \(N\ge3\).
For the uniform choice $\eta=1/5$, one can take $t=9/10$, because
$td_M\ge9/40>1/5$. This argument uses only the existing fixed-order
prime-class estimate, CRT and the finite upper sieve, independently
of the recent zero-free and graph theorems.
The use of \(t<1\) explains the strict range \(\eta<\min_i\eta_*(G_i)\);
no endpoint or optimality assertion is made.

\subsection{Two local filters on a common split core}\label{sec:form-local-assembly}

The remaining calculation concerns the two local shadows. It is shared
by all pairs above and by the principal specialization below.

First the local sets have a finite description, even when a split
conductor prime permits unbounded valuations. Put
\[
 e_p=1+2\max_i v_p(\Delta_i)\qquad(p\in S).
\]
At a primitive vector \(v\), the identity
\[
 \operatorname{adj}(\operatorname{Hess}Q_i)\nabla Q_i(v)
                  =-\Delta_i v
\]
shows that some partial derivative has valuation
\(s\le v_p(\Delta_i)\). For any target congruent to \(Q_i(v)\)
modulo \(p^{e_p}\), strong one-variable Hensel lifting in that
coordinate applies because \(e_p>2s\). The lifted vector stays
primitive. Thus each \(\mathscr V_{p,Q_i}\) is a union of full
cosets modulo \(p^{e_p}\), found by enumerating primitive vectors
in that quotient, and
\begin{equation}\label{eq:norm-local-density-finite}
 d_{p;1,2}(N)=\frac1{p^{e_p}}
 \#\left\{a\bmod p^{e_p}:
 \begin{array}{l}
 a\in\mathscr V_{p,Q_1}\bmod p^{e_p},\\
 N-a\in\mathscr V_{p,Q_2}\bmod p^{e_p}
 \end{array}\right\}.
\end{equation}
In particular these densities are rational periodic functions.
Outside \(S\), the primitive value set of either form is \(\Z_p\)
at split primes and \(\Z_p^\times\) at inert primes. These primes
are odd, so every target is a sum of two permitted local values.
Local admissibility is therefore the finite criterion
\(\prod_{p\in S}d_{p;1,2}(N)>0\), and this product has a fixed
positive minimum on admissible targets.

The local genus biases must still be retained. Set
\[
 L=8\prod_{\substack{p\in S\\p\text{ odd}}}p,\qquad
 \alpha=\varphi(L)/L,\quad H=\ker(\chi\bmod L),\quad
 \mathcal C=\langle p:p\notin S,\ \chi(p)=1\rangle.
\]
For each \(S\)-supported \(r\), put
\[
 A_{r,i}=\{u\in(\Z/L\Z)^\times:
                      ru\in\mathscr V_{p,Q_i}\text{ for all }p\in S\}.
\]
These sets are well defined and are unions of unit-square states:
equal units modulo an odd \(p\), or modulo eight at two, differ
by a local unit square; multiplication of a primitive vector by a
rational unit scales its value by that square.

Crucially \(A_{r,i}\subset H\), although individual local genus
characters can be nonprincipal. To verify this, use an ideal \(A_i\)
as above, with norm \(a_i\) and local multipliers \(\lambda_{i,p}\).
At every prime outside \(S\), an ideal norm has even valuation at
an inert prime and has no field-norm obstruction at a split prime.
Positivity and Hilbert reciprocity therefore give
\[
 \prod_{p\in S}(\lambda_{i,p},D)_p
      =\prod_{p\in S}(a_i,D)_p=1.
\]
Here the numerator in \(\lambda_{i,p}=N(z_{i,p})/a_i\) is a
field norm. If \(u\in A_{r,i}\), then
\((ru,D)_p=(\lambda_{i,p},D)_p\) at every bad prime.
The positive \(S\)-supported number \(r\) has product of bad
symbols one. For an integer unit representative \(u\), reciprocity
identifies \(\prod_{p\in S}(u,D)_p\) with \(\chi(u)\).
Multiplication of the displayed local relations gives \(\chi(u)=1\).
This does not identify \(A_{r,i}\) with the principal-form filter.

Consequently the exact disjoint decomposition is
\begin{equation}\label{eq:form-prefix-decomposition}\csname ltx@label\endcsname{eq:norm-prefix-decomposition}
 \Lambda_{Q_i}=\coprod_{r\in\langle S\rangle}
                  r\{x\in\mathcal C:x\bmod L\in A_{r,i}\}.
\end{equation}
Let $c_{\mathcal C}$ denote the Gamma-normalized marginal constant
of this core.  Its Euler series satisfies
\[
 \left(\sum_{n\in\mathcal C}n^{-s}\right)^2
  =\zeta(s)L(s,\chi)
      \prod_{p\in S}(1-p^{-s})(1-\chi(p)p^{-s})
      \prod_{\substack{p\notin S\\\chi(p)=-1}}(1-p^{-2s}),
\]
so
\begin{equation}\label{eq:norm-core-constant}
 c_{\mathcal C}^2=\frac{L(1,\chi)}\pi
       \prod_{p\in S}(1-p^{-1})(1-\chi(p)p^{-1})
       \prod_{\substack{p\notin S\\\chi(p)=-1}}(1-p^{-2}).
\end{equation}
Also put
\[
 C_{\rm in}=\prod_{\substack{p\notin S\\\chi(p)=-1}}
                         \left(1-\frac1{(p-1)^2}\right).
\]

\subsection{Finite profiles, the classical bounds, and genus comparison}

For fixed prefixes \(r,s\in\langle S\rangle\), put \(g=(r,s)\).
When \(g\mid N\), write \(A=r/g\), \(B=s/g\), \(M=N/g\), and set
\[
 \Gamma_{r,s;1,2}(N)=\frac L{|H|^2}
  \#\{(u,v)\in A_{r,1}\times A_{s,2}:Au+Bv=M\bmod L\};
\]
otherwise set it to zero. The original two filters are retained
after dividing by \(g\).

A finite collection of profiles covers every locally admissible
target. Put \(Z=\prod_{p\in S}p^{e_p+1}\). Given local summands
\(a+b=N\), replace them by \(a+p^{e_p}t,b-p^{e_p}t\).
They remain in their full value cosets. A valuation below \(e_p\)
stays fixed; otherwise at most one residue \(t\bmod p^2\) makes that
summand divisible by \(p^{e_p+2}\). Since \(p^2>2\), both forbidden
residues can be avoided. The summands can therefore be chosen with
valuations at most \(e_p+1\). Their valuations determine \(r,s\mid Z\),
and CRT assembles their unit quotients into a root counted by
\(\Gamma_{r,s;1,2}(N)\). Conversely, any such root lifts to an exact
linear solution at each bad prime: one of \(A,B\) is a unit, and
fixing the other unit residue preserves the two unit-square states.
The local definitions of \(A_{r,1},A_{s,2}\) therefore remain valid.
Outside \(S\) there is already no additive obstruction. Hence
\[
 N\text{ locally admissible}\quad\Longleftrightarrow\quad
 \Gamma_{r,s;1,2}(N)>0\text{ for some }r,s\mid Z.
\]

\begin{proof}[Proof of Theorem~\ref{cor:common-field-form-pairs}, part~\ref{part:form-classical}]
Lemma~\ref{lem:older-quadratic-filter} allows two distinct filters
\(U,V_0\subset H\). In a positive fixed profile it supplies
\(\gg_{\mathcal F}MJ_S(M)/\log M\) core pairs on \(Ax+By=M\).
Multiplication by \(r,s\) is injective, and
\(J_S(M)=J_S(N)\) because \(g\) is supported on \(S\). Only finitely
many profiles and root patterns occur above, so the minimum of the
positive lower constants and maximum of their thresholds give
\[
 R_{\Lambda_{Q_1},\Lambda_{Q_2}}(N)
                 \gg_{\mathcal F}NJ_S(N)/\log N
\]
uniformly on sufficiently large admissible targets. Sieving the
original line only at good inert primes gives the matching upper
bound by the unmarked case of \eqref{eq:norm-marked-sieve}.
Proposition~\ref{prop:norm-class-error} transfers the lower bound
to \(R_{Q_1,Q_2}\).

If \(Q'_i\) is in the genus of \(Q_i\) for each \(i\), their local
value sets coincide exactly, as do the two mixed local shadows.
Subtract their two class-error estimates and divide by the classical
positive lower bound. This proves
part~\ref{part:form-classical} of Theorem~\ref{cor:common-field-form-pairs}. If both doubled groups
vanish the individual global and local value sets are equal, proving
the exact assertion for every target. Nonadmissible targets have zero
count. This argument uses the older quadratic filter estimate, not
Theorem~\ref{thm:general}.
\end{proof}

\subsection{Prefix summation and Haar normalization}

Under the inputs of Section~\ref{sec:inputs}, apply
Theorem~\ref{thm:general} with the coprime fixed \(L\)-supported
coefficients \(A,B\) in each profile. Write \(R_{r,s;1,2}(N)\)
for its contribution in \eqref{eq:form-prefix-decomposition}.
Since \(J_S(N/g)=J_S(N)\), its contribution divided by
\(NJ_S(N)/\log N\) is
\begin{equation}\label{eq:norm-prefix-main}
 c_{\mathcal C}^2 C_{\rm in}\Gamma_{r,s;1,2}(N)\frac g{rs}+o(1).
\end{equation}
The original filters \(A_{r,1},A_{s,2}\) remain in this formula.

To sum the fixed-profile asymptotics, parameterize \(rx+sy=N\).
Its positive interval has length at most \(Ng/(rs)\). At any good
inert prime the two coefficients in this parameterization are units,
and its two forbidden roots merge exactly when the prime divides
\(N\). Discarding both residue filters and applying the unmarked
sieve \eqref{eq:norm-marked-sieve} with \(z=N^{1/10}\) gives
\begin{equation}\label{eq:norm-prefix-tail}
 R_{r,s;1,2}(N)\ll_{\mathcal F}
       \frac{NgJ_S(N)}{rs\log N}+O(N^{9/20}),
\end{equation}
uniformly in \(r,s\). The interval error also covers intervals shorter
than the sieve level. Only
\(O_{\mathcal F}((1+\log N)^{2|S|})\) prefix pairs can occur, while
\[
 \sum_{r,s\in\langle S\rangle}\frac{(r,s)}{rs}
 =\prod_{p\in S}\sum_{i,j\ge0}p^{-\max(i,j)}
 =\prod_{p\in S}\frac{1+p^{-1}}{(1-p^{-1})^2}<\infty.
\]
The total interval error is negligible relative to
\(NJ_S(N)/\log N\); the remaining tail is uniformly small as a fixed
prefix bound tends to infinity. This justifies a finite-prefix limit
followed by a summable-tail limit, without assuming a binary
asymptotic uniform in growing coefficients.

Partition the local Haar line \(a+b=N\) by the exact valuation
prefixes \(r,s\). At a prime with \(v_p(r)=i\), \(v_p(s)=j\), and
\(v_p(L)=k\), each compatible unit-pair root modulo \(p^k\) has
measure \(p^{-\max(i,j)-k}\). For \(i\le j\), prescribing the
second summand modulo \(p^{j+k}\) fixes the first summand's required
unit residue; the case \(j\le i\) is symmetric. Thus each global root
has measure \(g/(Lrs)\), and the full stratum has measure
\[
 \frac g{Lrs}
  \#\{(u,v)\in A_{r,1}\times A_{s,2}:Au+Bv=M\bmod L\}
 =\frac{\alpha^2}{4}\Gamma_{r,s;1,2}(N)\frac g{rs},
\]
because \(|H|=\varphi(L)/2\). Zero values have Haar measure zero.
Summing the disjoint strata gives the single normalization identity
\begin{equation}\label{eq:form-haar-strata}\csname ltx@label\endcsname{eq:norm-haar-strata}
 \sum_{r,s}\Gamma_{r,s;1,2}(N)\frac g{rs}
             =\frac4{\alpha^2}\prod_{p\in S}d_{p;1,2}(N).
\end{equation}
Combining this identity with \eqref{eq:norm-prefix-main} and the
summable tail proves
\[
 R_{\Lambda_{Q_1},\Lambda_{Q_2}}(N)\sim
 \frac{4c_{\mathcal C}^2}{\alpha^2}C_{\rm in}
       \left(\prod_{p\in S}d_{p;1,2}(N)\right)
                              J_S(N)\frac N{\log N}.
\]
The coefficient is explicit. In \eqref{eq:norm-core-constant},
divide the bad-prime factors by
\(\alpha^2=\prod_{p\in S}(1-p^{-1})^2\), and use
\[
 (1-p^{-2})\left(1-\frac1{(p-1)^2}\right)
                         =1-\frac2{p(p-1)}
\]
at every good inert prime. The result is precisely
\(\mathfrak b_{D,S}\) in \eqref{eq:form-universal-factor}.
The finite positive minimum of the local density product and
\eqref{eq:form-class-error} transfer the asymptotic to the actual
values. This proves part~\ref{part:form-leading} of
Theorem~\ref{cor:common-field-form-pairs}.

\subsection{Principal norms}\label{sec:norms}
\label{thm:norm-order}\label{thm:norm-leading}

Take \(Q(x,y)=N_{K/\Q}(x+y\omega)\) in a basis \(1,\omega\) of
a fixed order \(\mathcal O\) of discriminant \(\Delta=f^2D\).
Set \(P=P_Q\), \(\Lambda=\Lambda_Q\), \(R_T=R_{T,T}\),
\(S=\{p:p\mid2\Delta\}\), \(\chi=\chi_D\), and \(J_{\mathcal O}=J_S\).
Theorem~\ref{cor:common-field-form-pairs} with both slots equal to
\(Q\) gives the principal order bounds through part~\ref{part:form-classical}
and, under the analytic inputs, the leading asymptotic through
part~\ref{part:form-leading}. Its coefficient is
\begin{equation}\label{eq:norm-coefficient}
 \mathfrak A_{\mathcal O}(N)=
 \frac{4L(1,\chi)}{\pi}
 \prod_{p\in S}\frac{1-\chi(p)/p}{1-1/p}
 \prod_{\substack{p\notin S\\\chi(p)=-1}}
                  \left(1-\frac2{p(p-1)}\right)
 \prod_{p\in S}d_p(N).
\end{equation}
Here \(d_p=d_{p;1,2}\) for \(Q_1=Q_2=Q\).
The norm form has \(A=\mathcal O\), hence \(c_Q=0\), so
Lemma~\ref{lem:form-ideal-shadow} becomes
\begin{equation}\label{eq:norm-shadow-ideals}
 n\in\Lambda\quad\Longleftrightarrow\quad
 \text{some primitive integral invertible }I
       \text{ has }n(I)=n,\ [I]\in2\operatorname{Pic}(\mathcal O);
\end{equation}
the condition \([I]=0\) characterizes \(P\). In particular,
Proposition~\ref{prop:norm-class-error} gives
\(R_\Lambda-R_P\ll_{\mathcal O}NJ_{\mathcal O}/(\log N)^{6/5}\),
and \(P=\Lambda\) when the doubled class group vanishes.

\begin{remark}[Gaussian normalization]
For $\mathcal O=\Z[i]$ one has $P=C_4\sqcup2C_4$ and
\[
 d_2(N)=
 \begin{cases}
 1/4,&N\bmod8\in\{2,3,6,7\},\\
 1/8,&N\bmod8=4,\\
 0,&N\bmod8\in\{0,1,5\}.
 \end{cases}
\]
Consequently the coefficient of $J_{\mathcal O}(N)N/\log N$ is
$\frac12\prod_{p\equiv3\ (4)}(1-2/(p(p-1)))$ in the first four
classes and half that coefficient in the class $4\bmod8$.
This also calibrates the value-density normalization: the two Gamma
factors are already included in $c_{\mathcal C}^2$, and no additional
beta-integral factor occurs.
\end{remark}

\section{Consequences and boundaries}
\begin{corollary}\label{cor:basis13}
Under the analytic inputs of Section~\ref{sec:inputs}, the set
\(C_{12}\cup\{0\}\) is an asymptotic basis of exact order
thirteen: every sufficiently large integer is a sum of at most
thirteen elements of \(C_{12}\), and twelve do not suffice.
\end{corollary}
\begin{proof}
For a large target choose \(j\in\{0,\ldots,11\}\) with
\(N-j\equiv2\pmod{12}\). Theorem~\ref{thm:c12} writes \(N-j\)
as two core values; add \(j\) copies of one. This uses at most thirteen
positive summands. Conversely, \(25+60r\), \(r\ge0\), is
\(1\pmod{12}\) and divisible by the forbidden prime five.
At most twelve positive core summands would force exactly one
summand on that residue class, which is impossible. Padding with
zeros gives the equivalent fixed-order assertion.
\end{proof}

The field formulation follows from classical abelian class field
theory. A subgroup condition on Frobenius is complete splitting in
its fixed field. Finite conjunctions become intersections of inverse
image subgroups in the compositum. Kronecker--Weber places this
compositum in a cyclotomic field
\cite[Chapter I, Theorem 4.16]{milnecft}, giving precisely the
fixed subgroups in Theorem~\ref{thm:general}, after discarding the
fixed bad-prime set. The reduction is classical; the pointwise
indicator correlation is the analytic assertion used here.

Exact classifications and their finite certificates are treated in
the companion \emph{Exact Sumsets of Primitive Quadratic Norms}
\cite{exactcompanion}.
The separate article
\href{https://hexagonmath.org/2611.00593v1}{\emph{Pointwise representation
counts for Eisenstein restricted sums}} concerns the placement of summands.
Higher-fold sums and independent coupled systems can use classical
linear-form correlations, with a different proof from this reflected
binary problem. Growing-shift uniformity and higher indicator correlations
are not covered by the statements here.

\appendix
\section{An independent classical bound for quadratic cores}\label{sec:older-route}

The quadratic divisor pairing has two factors, so completed
Kloosterman sums and positive restoration give the required classical
order bounds.

Fix an even modulus \(L_0\), a nonprincipal quadratic literal character
\(\chi\bmod L_0\), and put
\[
 \begin{gathered}
 H=\ker\chi,\qquad
 \mathcal C=\langle p:p\nmid L_0,\ \chi(p)=1\rangle,\\
 V(n)=\tau(n)\1_{\mathcal C}(n),\qquad
 J(N)=\prod_{\substack{p\mid N\\p\nmid L_0,\ \chi(p)=-1}}
                          \frac{p-1}{p-2}.
 \end{gathered}
\]
All arithmetic data in the following lemma are fixed.

\begin{lemma}[Classical quadratic core bound]\label{lem:older-quadratic-filter}
Let \(A,B\) be coprime positive \(L_0\)-supported integers and
\(U,V_0\subset H\) fixed residue subsets. For any positive \(A',B'\), define
\[
 \Gamma_{A',B';U,V_0}(N)=\frac{L_0}{|H|^2}
   \#\{(a,b)\in U\times V_0:A'a+B'b=N\bmod L_0\}.
\]
Uniformly over all sufficiently large targets with \(\Gamma_{A,B;U,V_0}(N)>0\),
\[
 \#\{Ax+By=N:x,y\in\mathcal C,\
                  x\bmod L_0\in U,\ y\bmod L_0\in V_0\}
                   \asymp\frac{NJ(N)}{\log N}.
\]
Zero finite density gives zero count identically.  In particular the
lemma applies to unions of unit-square residue states.
Its proof uses only fixed-character means, fixed-progression prime
Mertens, the classical Weil--Estermann bound, and finite sieve inequalities.
\end{lemma}

\paragraph{Proof plan and parameter order.}
Completed two-factor boxes give a shifted weighted family, uniform
under the divisibility predicates needed by a finite lower sieve.
The rough mass has lower-bound scale \(NJ(N)/(\log N)^2\).
Restoration supplies the missing logarithm through the positive
product in \eqref{eq:older-restoredmass}; the valuation patterns
\((1,1),(1,0),(0,1)\) recover its labels \(d,u,v\), and the residual
rough divisor weights are bounded. This makes restoration injective
at the numerical-value level.

Choose the box exponent and core family first, then the saving \(\eta\)
and external exponent \(\theta\) from \eqref{eq:older-uniform-reserve}.
Next choose \(\alpha,\xi\) by \eqref{eq:older-smallbudgets}, the
Bonferroni depth, the roughness exponent, and finally the restoration
range, in that order. All remain fixed as \(N\to\infty\).

\subsection{Two-factor boxes with a fixed period}

For bounded functions \(\omega_1,\omega_2\) of fixed period \(c\),
intervals \(I_1,I_2\), and \((a,m)=(c,m)=1\), let
\[
 S(a)=\sum_{\substack{r\in I_1,\ s\in I_2\\rs\equiv a\ (m)}}
                        \omega_1(r)\omega_2(s).
\]
Then, for every \(\epsilon>0\),
\begin{equation}\label{eq:older-box}
 S(a)-\frac1{\varphi(m)}\sum_{b\bmod m}^{*}S(b)
                         \ll_{c,\epsilon}m^{1/2+\epsilon}.
\end{equation}
The estimate is uniform in the intervals, including intervals longer
than \(cm\), and permits subintervals of arbitrary lengths.

Here is a direct completion proof.  Put \(Q=cm\),
\[
 W_i(\nu)=\sum_{n\in I_i}\ee{-\nu n/Q},\qquad
 G_i(t)=\sum_{r\bmod c}\omega_i(r)\ee{tr/c}.
\]
CRT and finite Fourier inversion express \(S(a)\) as
\[
 Q^{-2}\sum_{\nu,k\bmod Q}
 W_1(\nu)W_2(k)G_1(m^{-1}\nu)G_2(m^{-1}k)
       K_1(c^{-1}\nu,ac^{-1}k;m),
\]
where \(K_1(u,v;m)=\sum_{z\bmod m}^{*}\ee{(uz+v\bar z)/m}\);
inverses in \(G_i\) are taken modulo \(c\).
The terms \(\nu=0\bmod m\) or \(k=0\bmod m\) are independent of
the unit \(a\), since the resulting Ramanujan sum is unchanged by
unit multiplication.  For all other frequencies use
\[
 |K_1(u,v;m)|\le\tau(m)m^{1/2}(u,v,m)^{1/2}.
\]
This is the ordinary Weil--Estermann bound stated and applied in
\cite[printed p.~44]{heathbrown}.  The primary source's
Lemma~6 is its singleton-padded divisor-box consequence.
In centered frequency ranges,
\(|W_i(\nu)|\ll\min(|I_i|,Q/|\nu|)\) and \(|G_i|\le c\).
The remaining double sum is bounded using
\[
 \sum_{\substack{0<|\nu|,|k|\le Q/2\\m\nmid \nu,\ m\nmid k}}
      \frac{(\nu,k,m)^{1/2}}{|\nu k|}
 \le \sum_{d\mid m}d^{1/2}
       \left(\sum_{\substack{0<|\nu|\le Q/2\\d\mid \nu}}\frac1{|\nu|}\right)^2
 \ll_c \log^2(2m)\sum_{d\mid m}d^{-3/2}.
\]
The \(Q^2\) from the two geometric sums cancels the completion
normalization.  This proves a bound
\(O_c(m^{1/2}\tau(m)\log^2(2m))\) around an \(a\)-independent
quantity.  Averaging that quantity over unit \(a\) changes the
error by at most a factor two and proves \eqref{eq:older-box}.
The case \(m=1\) has discrepancy zero.
Tensor partial summation also allows bounded-variation interval weights.

\subsection{Shifted kernels, means, and fixed filters}

For \(t\in\mathcal C\) put \(F_t(n)=V(tn)/V(t)\).
At a split prime with \(b=v_p(t)\),
\[
 F_t(p^e)=1+\frac{e}{b+1};
\]
at a nonsplit prime it is zero for \(e>0\).
With literal characters, its Dirichlet series is
\[
 \sum_{n\ge1}\frac{F_t(n)}{n^s}
   =L(s,\chi_0^{L_0})L(s,\chi)\sum_{j\ge1}\frac{e_t(j)}{j^s}.
\]
The correction factors are \(1-p^{-2s}\) at inert primes and
\(1-\frac{b}{b+1}p^{-s}\) at split primes dividing \(t\).
Deleting all Euler factors at primes dividing another integer \(D\)
replaces its split local correction by \((1-p^{-s})^2\).
Thus all these kernels satisfy, uniformly in valuations,
\begin{equation}\label{eq:older-kernel}
 |e(j)|\le\tau(j)^C,\qquad
 \sum_j|e(j)|j^{-3/5}\ll_{\epsilon}(tD)^\epsilon.
\end{equation}
Every kernel index is an \(L_0\)-unit with \(\chi(j)=1\).

We record explicitly the elementary mean used below.  If
\(\xi\) is a character of a fixed modulus \(c\) with the same prime
support as \(L_0\), then the twisted series has character factors
\(\xi,\xi\chi\).
It has a principal factor only when the restriction of \(\xi\)
is \(1\) or \(\chi\); in those cases \(F_t\xi=F_t\).
Otherwise both character sums have bounded partial sums, and
Dirichlet hyperbola at \(\sqrt X\) gives \(O_c(\sqrt X)\) for
their convolution.
In the principal cases, split the principal factor from the
nonprincipal one by the same hyperbola argument to obtain its residue
times \(X\), with error \(O_c(\sqrt X)\).
Summing kernels by \eqref{eq:older-kernel}, including their main tails,
therefore gives
\begin{equation}\label{eq:older-mean}
 \sum_{n\le X}F_t(n)\1_{(n,D)=1}\xi(n)
 =\delta_{\xi\in\{1,\chi\}}\rho_{t,D}X
                         +O_\epsilon((tDX)^\epsilon X^{3/5}).
\end{equation}
Here the two displayed characters mean their inflations to \(c\).
The error and main remain valid with any fixed periodic unit filter,
after finite character expansion.
For later normalization write
\[
 M_{t,p}=(1-p^{-1})\sum_{e\ge0}F_t(p^e)p^{-e},\qquad
 \rho_t=\frac{\varphi(L_0)}{L_0}\prod_{p\nmid L_0}M_{t,p}.
\]
The product is taken in natural prime order; its convergence and
positivity follow equally from the character factorization and
\(L(1,\chi)>0\).  A deleted split prime divides this mean by its
original local series.  The residue measure on the fixed unit
group is uniform on \(H\).

\begin{proposition}[Uniform shifted weighted family]\label{prop:older-weighted}
There exist \(\eta,\theta>0\), with \(\theta\) chosen sufficiently
small after \(\eta\), such that the following holds.
Let \(t_1,t_2,u,v\in\mathcal C\), each at most \(N^\theta\), with
no restriction on their gcds.  For fixed residue subsets \(U,V_0\subset H\),
\begin{align*}
 &\sum_{\substack{Au x+Bv y=N\\x,y\ge1,\ x\bmod L_0\in U,\ y\bmod L_0\in V_0}}
                      F_{t_1}(x)F_{t_2}(y)\\
 &\quad=\frac{N}{ABuv}\,
       \Gamma_{Au,Bv;U,V_0}(N)\rho_{t_1}\rho_{t_2}
       \prod_{p\nmid L_0}\frac{D_p}{M_{t_1,p}M_{t_2,p}}
                         +O(N^{1-\eta}),
\end{align*}
uniformly in every target valuation.  Here
\[
 D_p=\lim_{m_p\to\infty}p^{-m_p}
 \sum_{\substack{a,b\bmod p^{m_p}\\Au a+Bv b=N\ (p^{m_p})}}
       F_{t_1}(p^{\min(v_p(a),m_p)})
       F_{t_2}(p^{\min(v_p(b),m_p)}).
\]
All factors, including zero factors, are retained.
The same formula holds after imposing any fixed finite combination
of small divisibility and exact-valuation tests whose resulting
shifts and coefficients stay in this family.
\end{proposition}

\begin{proof}
First suppose \((u,v)=(uv,N)=1\).
For \(g=(x,N)\), write \(x=gx'\), \(y=gy'\), \(M=N/g\);
a nonzero term has \(g\in\mathcal C\).
Its exact prefactor is \(F_{t_1}(g)F_{t_2}(g)\), its new shifts
are \(t_1g,t_2g\), and \((x',M)=1\).
Transport the residue states to \(U/g,V_0/g\).
Every such multiplication is a unit bijection, so the corresponding
finite root count, with its transported target, is unchanged.

Choose a small fixed \(\lambda>0\), put \(K=N^\lambda\), and truncate
\(g\) and the two kernels at \(K\).
Large gcds cost \(O(N^{1-\lambda+\epsilon})\).
For either kernel, counting multiples and using
\eqref{eq:older-kernel} gives
\[
 N^{1+\epsilon}\sum_{j>K}\frac{|e(j)|}{j}
   +N^\epsilon\sum_{j\le N}|e(j)|
       \ll N^{1-2\lambda/5+\epsilon}+N^{3/5+\epsilon}.
\]
Discard also \(Au x\le N^{1-\lambda}\), at
\(O(N^{1-\lambda+\epsilon})\).
All implicit valuation and shift factors here are divisor-bounded.

We suppress primes on \(x',y'\) in the remainder of this core argument.
The fixed filter on the original first unit variable is
\[
 w(x)=\1_{x\bmod L_0\in U/g}\,
       \1_{B\mid M-Au x}\,
       \1_{(z,L_0)=1,\ z\bmod L_0\in v(V_0/g)},\qquad
                       z=(M-Au x)/B.
\]
It has fixed period \(c=BL_0\).  Expand it on unit classes by characters
modulo \(c\), retaining every scalar phase of every kernel index.
The filter ensures that
\(t=(M-Au x)/(Bvk)\) is an \(L_0\)-unit with \(\chi(t)=1\).
Consequently the exact symmetric divisor pairing is
\[
 (\chi_0^{L_0}*\chi)(t)
  =2\sum_{\substack{a\mid t\\a<\sqrt t}}\chi(a)
                         +\1_{t=\square}\chi(\sqrt t).
\]
The square endpoint costs \(O(N^{1/2+O(\lambda+\theta)+\epsilon})\).

For a second kernel \(k\le K\), the remaining congruence is
\[
       Au x\equiv M\pmod{vka},\qquad
       Au x+Bvka^2<M,\qquad (x,M)=1 .
\]
The entire bad-supported factor \(B\) has already been imposed in
\(w\); all varying moduli are coprime to \(c\).
Also \((k,uM)=(a,uM)=1\).
Expand the first kernel \(j\) and its two character variables.
A nonzero congruence requires \((uj,vka)=1\); otherwise it would
force a prime divisor of \(uM\) into an excluded place, or make a
unit equal to a nonunit.  The same zero test holds in the unit model.

Partition the two character variables and \(a\) into relative-width
\(\Delta=N^{-\zeta}\) boxes, with \(\zeta=\lambda/100\).
There are \(O(N^{3\zeta+\epsilon})\) boxes.
On a box with its endpoints fixed, apply \eqref{eq:older-box} to
the first two factors, at the full modulus \(m=vka\).
It is enough that \(m\le N^{1/2+O(\lambda+\theta)}\);
summing its square-root error over \(a\le N^{1/2}\),
then the two kernel masses, gives
\begin{equation}\label{eq:older-box-budget}
 O(N^{3/4+100\lambda+100\theta+\epsilon}).
\end{equation}
This deliberately loose bound includes the fine boxes, the fixed
character expansions, and the \(g\)-strata.
In particular, no medium-variable dispersion argument occurs
in this quadratic specialization.

We justify the curved boundary separately.  A crossing box satisfies
\(|Au x+Bvka^2-M|\ll\Delta N\).
On the actual congruence write \(M-Au x=Bvkat\).
Then \(|Bvka(t-a)|\ll\Delta N\), which has
\[
 O\left(\frac{\Delta N\log N}{vk}+\sqrt{\frac N{vk}}\right)
\]
integer pairs \((a,t)\), after discarding only the divisibility
which determines \(x\).
The unique scalar decompositions and divisor-bounded first-kernel
multiplicities cost \(N^\epsilon\).
Summing the second kernel by \eqref{eq:older-kernel} leaves
\begin{equation}\label{eq:older-boundary}
 O_\epsilon\big(N^\epsilon(\Delta N\log N+\sqrt N K^{1/10})\big).
\end{equation}
We next display the finite main sum produced by this replacement,
before completing either the paired-divisor series or the target gcd
sum.  This also fixes the normalization of every local factor.
For a retained \(g\), put
\[
 C_g=F_{t_1}(g)F_{t_2}(g),\qquad
 \mathcal A=F_{t_1g},\quad \mathcal B=F_{t_2g},\quad
 H_{\mathcal A,p}=\sum_{i\ge0}\mathcal A(p^i)p^{-i},
\]
and define
\[
 \rho_{\mathcal A,M}
   =\rho_{\mathcal A}
       \prod_{\substack{p\mid M\\p\nmid L_0}}H_{\mathcal A,p}^{-1},
 \qquad
 E_{\mathcal A}(q)=\prod_{p\mid q}H_{\mathcal A,p}^{-1}.
\]
Let \(\mu_g\) be the average of the displayed filter \(w\) over
\(H_c\), the lift of \(H\) to \(c=BL_0\).
The contribution \(T_g\) of this stratum is \(C_g\mathcal P_g\),
up to the errors estimated above, where the \emph{finite} sum is
\begin{equation}\label{eq:older-finite-main}
 \begin{aligned}
 \mathcal P_g={}&
 \frac{2\mu_g\rho_{\mathcal A,M}}{Au}
 \sum_{\substack{k\le K\\(k,uM)=1}}e_{\mathcal B}(k)\\
 &{}\times
 \sum_{\substack{a<\sqrt{M/(Bvk)}\\(a,L_0uM)=1}}
 \frac{\chi(a)E_{\mathcal A}(vka)}{\varphi(vka)}
                         (M-Bvka^2).
 \end{aligned}
\end{equation}
At a zero-length endpoint the summand is zero, so the strict
endpoint convention has no effect.

To check the step producing \eqref{eq:older-finite-main}, the box
replacement contributes \(1/\varphi(q)\), \(q=vka\), times the
unit population of the original first sequence, with its fixed
filter \(w\).  The fixed-character mean gives, uniformly in \(q\),
\[
 \sum_{\substack{n\le X\\(n,Mq)=1}}\mathcal A(n)w(n)
  =\mu_g\rho_{\mathcal A,M}E_{\mathcal A}(q)X
       +O_\epsilon((t_1gMqX)^\epsilon X^{3/5}).
\]
Indeed the character expansion of \(w\) has principal contributions
only at \(1,\chi\), whose coefficients sum to \(\mu_g\).
The first-kernel indices retain their scalar character phases;
recombining them yields exactly this original-sequence mean.
Their main terms retain \(1/j\), while the errors retain \(j^{-3/5}\).
The interval length is \(X=(M-Bvka^2)/(Au)\), and divisor pairing
supplies the factor two.  This proves the finite formula.
Its mean errors sum to
\(O(N^{3/5+100\lambda+100\theta+\epsilon})\): use
\(1/\varphi(vka)\ll_\epsilon(vka)^{-1+\epsilon}\) and
\eqref{eq:older-kernel}.
The main boundary intervals keep the same reciprocal kernel sums.
Combining these errors with \eqref{eq:older-box-budget},
\eqref{eq:older-boundary}, the truncations, and the square endpoints
leaves the following uniform total reserve:
\begin{align}
 \mathcal R_N\ll_\epsilon{}&
 N^{1-2\lambda/5+\epsilon}+N^{1-\lambda+\epsilon}
       +N^{1-\zeta+\epsilon}\notag\\
 &+N^{3/4+100\lambda+100\theta+\epsilon}
       +N^{3/5+100\lambda+100\theta+\epsilon}
       +N^{1/2+100\lambda+100\theta+\epsilon}.
 \label{eq:older-uniform-reserve}
\end{align}
The sum of the errors over retained \(g\), including their
divisor-bounded \(C_g\), obeys this reserve, after decreasing
\(\epsilon\) in the preceding estimates.

The finite-place coefficient can now be checked independently.
The group \(H_c\) has \(B|H|\) elements.
Each allowed pair modulo \(L_0\) has \(B\) pair lifts modulo \(BL_0\),
since \(Au\) is a unit at primes dividing \(B\).
Each surviving first residue in the definition of \(w\) admits
exactly \(B\) second-residue lifts.
Thus the number of such first residues is the original number of
pairs modulo \(L_0\), and
\begin{equation}\label{eq:older-filter-normalization}
 \Gamma_g:=
 \Gamma_{Au,Bv;U/g,V_0/g}(M)
       =\frac{2BL_0}{\varphi(L_0)}\mu_g
       =\frac{2B}{\alpha_0}\mu_g,\qquad
                  \alpha_0=\varphi(L_0)/L_0 .
\end{equation}
The map \((x',y')\mapsto(gx',gy')\) bijects these root pairs with
the original ones.  Hence \(\Gamma_g\), and therefore \(\mu_g\),
is constant through the \(g\)-sum.  Denote this common mean by \(\mu\).

We complete the two finite sums in
\eqref{eq:older-finite-main} with explicit error bounds.
Set \(M_{\mathcal A,p}=(1-p^{-1})H_{\mathcal A,p}\) and
\(W_p=M_{\mathcal A,p}^{-1}\).
Then
\[
 \frac{E_{\mathcal A}(vka)}{\varphi(vka)}
   =\frac1{vka}\prod_{p\mid vk}W_p
                    \prod_{\substack{p\mid a\\p\nmid vk}}W_p.
\]
For fixed \(k\) put
\[
 b_k(a)=\chi(a)\1_{(a,L_0uM)=1}
                       \prod_{\substack{p\mid a\\p\nmid vk}}W_p,
 \qquad Y_k=\sqrt{M/(Bvk)}.
\]
For every allowed shift,
\(1-p^{-1}\le M_{\mathcal A,p}\le(1-p^{-1})^{-1}\).
Thus \(W_p=1+O(1/p)\), uniformly in all its valuations.
The Dirichlet series of \(b_k\) is \(L(s,\chi)\) times a correction
absolutely convergent for \(\Re s>0\).
At a prime outside \(uMvkL_0\), its correction is explicitly
\(1+(W_p-1)\chi(p)p^{-s}\); at primes of \(uM\) it has the finite
deletion \(1-\chi(p)p^{-s}\), and at primes of \(vk\) no new
correction is required.
Since the partial sums of the fixed nonprincipal \(\chi\) are
bounded, convolution with this correction gives
\[
 \sum_{a\le Y}b_k(a)\ll_\epsilon
                 (t_1g\,uMvkY)^\epsilon .
\]
In particular \(\sum b_k(a)/a\) converges, its tail after \(Y\)
is \(O(N^\epsilon/Y)\), and
\(\sum_{a\le Y}a\,b_k(a)=O(N^\epsilon Y)\).
Consequently
\[
 \begin{split}
 &\sum_{a<Y_k}
 \frac{\chi(a)E_{\mathcal A}(vka)}{\varphi(vka)}
                                      (M-Bvka^2)\\
 &\qquad=
 \frac{M}{vk}\prod_{p\mid vk}W_p
                       \sum_{a\ge1}\frac{b_k(a)}a
        +O_\epsilon\!\left(N^\epsilon\sqrt{\frac{M}{vk}}\right).
 \end{split}
\]
The fixed \(B\) is absorbed in the error constant.
Summing against \(e_{\mathcal B}(k)\) costs at most \(K^{1/10}\)
in the square-root error, by \eqref{eq:older-kernel}.
The completed main has an absolutely bounded correction times
\(e_{\mathcal B}(k)/k\), up to \(N^\epsilon\).
Its \(k>K\) tail is therefore
\(O(N^{1-2\lambda/5+\epsilon})\).
Both errors are included in \eqref{eq:older-uniform-reserve}.

Here is the Euler calculation of the completed sum, without an
implicit normalization.  Write \(z=p^{-1}\), \(q_p=1-z\),
\[
 A_i=\mathcal A(p^i),\qquad B_i=\mathcal B(p^i),\qquad
 P_{\mathcal B,p}(z)=\sum_{i\ge0}e_{\mathcal B}(p^i)z^i.
\]
The kernel identity says
\[
 \frac{P_{\mathcal B,p}(z)}{1-\chi(p)z}
                    =(1-z)\sum_{i\ge0}B_i z^i
                    =M_{\mathcal B,p}.
\]
The double sum over \(k,a\), after extracting its real factor
\(M/(uv)\), has the following local factors:
\begin{align*}
 \mathcal L_p
 &=1+W_p\left(\frac{P_{\mathcal B,p}(z)}{1-\chi(p)z}-1\right)
       &&(p\nmid uvM)\\
 &=\frac{M_{\mathcal A,p}+M_{\mathcal B,p}-1}
                      {M_{\mathcal A,p}}
       &&(p\nmid uvM),\\
 \mathcal L_p&=1 &&(p\mid uM),\\
 \mathcal L_p
 &=\frac{W_pP_{\mathcal B,p}(z)}{1-\chi(p)z}
       =\frac{M_{\mathcal B,p}}{M_{\mathcal A,p}}
       &&(p\mid v).
\end{align*}
For \(p\mid uM\), both local indices \(k,a\) are excluded.
For \(p\mid v\), the factor \(W_p\) is already present even at
\(v_p(k)=v_p(a)=0\); this explains that row of the table.
The first mean is
\[
 \rho_{\mathcal A,M}
 =\alpha_0\prod_{p\nmid L_0}M_{\mathcal A,p}
          \prod_{\substack{p\mid M\\p\nmid L_0}}H_{\mathcal A,p}^{-1}.
\]
At \(p\mid M\) its combined factor is therefore
\(M_{\mathcal A,p}/H_{\mathcal A,p}=q_p\).
Multiplying the other rows by this first mean gives the completed
fixed-\(g\) expression
\begin{equation}\label{eq:older-completed-g}
 \begin{aligned}
 &C_g\,\frac{M}{Auv}\,2\alpha_0\mu
 \prod_{\substack{p\mid M\\p\nmid L_0}}q_p
 \prod_{p\mid u}M_{\mathcal A,p}\\
 &\qquad{}\times
 \prod_{p\mid v}M_{\mathcal B,p}
 \prod_{\substack{p\nmid L_0uvM}}
                (M_{\mathcal A,p}+M_{\mathcal B,p}-1).
 \end{aligned}
\end{equation}
Intermediate products are taken in natural prime order.
Their factorization into the fixed \(L(1,\chi)\) and absolutely
convergent corrections justifies every completion; the final
normalized product in the proposition is absolutely convergent.

We may extend the divisor sum from \(g\le K\) to all \(g\mid N\).
The means and local factors in \eqref{eq:older-completed-g},
with their normalized shifts, are \(N^{O(\epsilon)}\) uniformly,
whereas \(C_gM\ll_\epsilon N^{1+\epsilon}/g\).
The added main is \(O(N^{1-\lambda+O(\epsilon)})\).
At \(p^n\Vert N\), use now the original local coefficients
\(A_i=F_{t_1}(p^i)\), \(B_i=F_{t_2}(p^i)\).
The choices \(v_p(g)=i<n\) in \eqref{eq:older-completed-g}
give \(q_p A_iB_i z^i\).
The choice \(i=n\) gives
\[
 z^nA_nB_n
    \big(M_{A^{[n]},p}+M_{B^{[n]},p}-1\big),\qquad
          A^{[n]}_j=A_{n+j}/A_n,\quad B^{[n]}_j=B_{n+j}/B_n.
\]
A zero \(A_nB_n\) means this choice is absent and no division
by that coefficient is made.
Expanding the two shifted means produces exactly
\[
 \begin{split}
 D_p=I_{p,n}={}&q_p\sum_{i<n}A_iB_i z^i
       +z^n(1-2z)A_nB_n\\
 &+q_p\left(B_n\sum_{i>n}A_i z^i
                         +A_n\sum_{j>n}B_j z^j\right).
 \end{split}
\]
At \(p\mid u\) or \(p\mid v\), the target is a unit and the
respective factors remain \(M_{A,p}\) or \(M_{B,p}\), exactly
the original local line integral.
Finally \eqref{eq:older-filter-normalization} implies
\[
 \frac{2\alpha_0\mu}{Auv}
       =\frac{\Gamma_{Au,Bv;U,V_0}(N)\alpha_0^2}{ABuv},
 \qquad
 \rho_{t_1}\rho_{t_2}=\alpha_0^2
                      \prod_{p\nmid L_0}M_{t_1,p}M_{t_2,p}.
\]
Together with \eqref{eq:older-completed-g}, these identities
give precisely the proposition's main term in the unit-coefficient
target core, including the bad coefficient \(B\).

For completeness, the general coefficient reductions also have an
explicit local check.  At a good prime write
\(a=v_p(u),b=v_p(v),n=v_p(N)\), and \(t=\min(a,b)\).
If \(n<t\), the density is zero.
Otherwise dividing the congruence by \(p^t\) gives
\[
 D_p=p^tD_p^{\rm red}(a-t,b-t,n-t).
\]
The factor \(p^t\) follows from counting the two sets of lifts
and the one-dimensional line normalization \(p^{-m_p}\).
If the remaining valuations are \(a>0,b=0\), solving for the
unit-coefficient variable gives
\[
 D_p^{\rm red}=
 \begin{cases}
    B_nM_{A,p},&n<a,\\
    B_a I_{p,n-a}(A,B^{[a]}),&n\ge a,
 \end{cases}
 \qquad B^{[a]}_j=B_{a+j}/B_a .
\]
There is the symmetric formula when \(b>0,a=0\).
A zero normalizing coefficient makes the relevant density zero.
These formulas retain arbitrarily large coefficient and target
valuations.

Globally remove \(r=(u,v)\), requiring \(r\mid N\), and divide
the coefficients and target by \(r\).
For the coprime remaining coefficients put \(d=(u,N)\), \(e=(v,N)\).
Then \(e\mid x,d\mid y\); the exact substitutions are
\[
 x=ex',\quad y=dy',\quad
 (t_1,t_2)\mapsto(t_1e,t_2d),\quad
 (u,v,N)\mapsto(u/d,v/e,N/(de)),
\]
with prefactor \(F_{t_1}(e)F_{t_2}(d)\) and states
\(U/e,V_0/d\).  The new coefficients are prime to the new target.
All removed factors belong to \(\mathcal C\).
Their unit multipliers biject the finite root sets, with the
transported target and states.  The displayed local formulas show
that the radial integrals transform with exactly the same
prefactors and common-coefficient gains.
Thus reexpressing the transformed main in the original variables
gives the stated \(D_p\)-product, for every gcd profile.

All estimates in \eqref{eq:older-uniform-reserve} are uniform in
these profiles and in the shift valuations.  They depend on
divisor bounds and \eqref{eq:older-kernel}, not on an unbounded
number of local cases.  When the external family exponent is
decreased as below, transformed targets are at least
\(N^{1-3\theta}\), transformed parameters lie in the proved core
range, and the divisor-bounded prefactors preserve a smaller
positive saving.
Choose \(\lambda\) small, then \(\theta\) much smaller.
Equations \eqref{eq:older-box-budget} and
\eqref{eq:older-boundary}, the kernel saving \(2\lambda/5\),
and the gcd and small-summand saving \(\lambda\), leave a fixed
positive core saving.  Choose \(\eta\) below a hundredth of the
boundary saving, and decrease the external family exponent again.
The transformed targets are at least \(N^{1-3\theta}\) and all
transformed parameters stay inside the core family.
Divisor prefactors preserve a smaller saving.
Finally \(b\mid x\) is the exact substitution \(x=bx'\), which
changes the shift to \(t_1b\), coefficient to \(ub\), state to
\(U/b\), and introduces \(F_{t_1}(b)\).
The analogous substitution on the second side, and finite differences
for exact valuations, prove the predicate assertion.
\end{proof}

\subsection{A finite lower sieve and positive restoration}

We now prove Lemma~\ref{lem:older-quadratic-filter}.
Fix an admissible target.  Its finite density \(\Gamma(N)\)
belongs to a fixed finite positive set, and hence has a positive
lower bound depending only on the fixed data.
At good primes write \(x=p^{-1}\), \(q_p=1-x\), and use the
unshifted coefficients \(A_i=B_i=V(p^i)\).
Let \(I_{p,e}\) be the radial integral in the preceding proof,
where \(e=v_p(N)\).
The weighted mass with both variables units is
\[
 c_e(p)=
 \begin{cases}1-2/p,&e=0,\\1-1/p,&e\ge1.\end{cases}
\]
At an inert prime \(I_{p,e}=c_e(p)\).
At a split prime define the bad-event probability
\(g_e(p)=1-c_e(p)/I_{p,e}\).
The explicit geometric series gives, uniformly in \(e\),
\begin{equation}\label{eq:older-roughprob}
             0\le g_e(p)<1,\qquad
             g_e(p)=4/p+O(p^{-2}).
\end{equation}
For example, at \(e=0\),
\(I_{p,0}=2(1-p^{-1})^{-1}-1\).
For \(e\ge1\), the \(i=0,1\) terms give
\(I_{p,e}=1+3/p+O(p^{-2})\); the remaining polynomial-geometric
tail is uniformly \(O(p^{-2})\).
The formula at \(e=1\) has the same expansion directly.
At every finite collection of primes \(I_{p,e}\) is uniformly bounded,
whereas \(c_e(p)\ge1/3\), so the good-event probabilities have a
uniform positive lower bound.

Put \(M_p=(1-p^{-1})\sum_iV(p^i)p^{-i}\) and
\[
 \mathcal V(z)=\prod_{\substack{p\le z\\p\nmid L_0}}
                  \frac{1-2/p}{M_p^2},\qquad
 E_p=\frac{p-1}{p-2},\qquad
 E_{p,e}=\begin{cases}1,&e=0,\\E_p,&e\ge1,\end{cases}
 \qquad a_p=\frac{p-1}{p(p-2)}.
\]
At a split prime \(M_p=(1-p^{-1})^{-1}\), so the factor
in \(\mathcal V\) is \(1-4/p+O(p^{-2})\).
At an inert prime it is \(1+O(p^{-2})\).
Fixed-character prime Mertens, and the split-prime density \(1/2\),
therefore give
\begin{equation}\label{eq:older-roughdimension}
                         \mathcal V(z)\asymp(\log z)^{-2}.
\end{equation}
This is the dimension-two rough weighted count which must recover
one power of \(\log N\).

Choose small fixed \(\alpha,\xi>0\) after the weighted family.
We use squarefree pairwise coprime triples \(d,u,v\) with
\[
 d,u,v\in\mathcal C,\quad d\mid N,\quad (uv,N)=1,\quad
 duv\le N^\alpha,
\]
all primes at most \(w=N^\omega\).
At target \(M=N/d\), prescribe \(v_p(x)=1\) for \(p\mid u\),
\(v_p(y)=1\) for \(p\mid v\), and require that no other prime
at most \(z=N^\gamma\) divide \(xy\).
Count with weight \(V(x)V(y)/(2^{\omega(u)}2^{\omega(v)})\)
and with residue states \(x\in U/d,\ y\in V_0/d\).
Thus the lifted values \(dx,dy\) have the original states.

All these predicates are covered by
Proposition~\ref{prop:older-weighted}.
Indeed exact valuation one is the difference of divisibility by
\(p\) and by \(p^2\), and a roughness predicate expands using
\[
 \1_{p\mid xy}=\1_{p\mid x}+\1_{p\mid y}
                              -\1_{p\mid x,\ p\mid y}.
\]
For a tested squarefree divisor \(\ell\le N^\xi\), each resulting
shift and coefficient is at most \(N^{2\alpha+\xi}\).
The expansion has \(N^{o(1)}\) terms.
It suffices to choose
\begin{equation}\label{eq:older-smallbudgets}
 2\alpha+\xi<(1-\alpha)\theta,\qquad
                         \alpha+\xi<\eta/4 .
\end{equation}
Forbidden shifts give zero terms; common-divisibility terms use the
noncoprime part of the weighted family.

The finite density remains exactly \(\Gamma(N)\) in every lower term.
After writing \(x=ux'\), \(y=vy'\), its states are
\(U/(du),V_0/(dv)\) and its equation is
\(Au x'+Bv y'=N/d\).
The map
\((x',y')\mapsto(du x',dv y')\) bijects its finite roots with the
original roots modulo \(L_0\).
This is why split primes with nontrivial residue states may be
restored.  Restricting restoration to identity-state primes would
lose the exponent needed below.

Here is the exact positive local model after imposing all roughness
conditions.  Its main mass for a triple is
\begin{equation}\label{eq:older-triplemass}
 \mathcal M_{d,u,v}
 =C_N(z)\frac{N}{ABd}\mathcal V(z)
       \prod_{\substack{p\mid N/d\\p\le z,\ p\nmid L_0}}E_p
       \prod_{p\mid uv}a_p .
\end{equation}
The positive factor \(C_N(z)\) includes \(\Gamma(N)\rho_1^2\)
and the Euler tail above \(z\).
It is bounded above and below by fixed positive constants,
uniformly in the target and the triples.
To verify this, outside the target the tail factors are
\(1+O(p^{-2})\); at target primes they are \(1+O(1/p)\) uniformly
in valuation.  At most \(1/\gamma\) target primes exceed \(z\).
All cofactor primes are below \(w<z\), so this tail is unchanged
when dividing by \(d\).
At a unit target prime, a valuation-one directive on either side
has weighted mass \(2(1-p^{-1})/p\).
Its normalization by two and division by \(1-2/p\) give exactly
\(a_p\), proving \eqref{eq:older-triplemass}.

We justify that a finite sieve recovers a fixed fraction of this
fully rough main, rather than asserting an asymptotic sieve.
Treat finitely many small primes exactly so that \(g_e(p)\le1/2\)
at the others.  Split the remaining primes into blocks
\((z^{2^{-j-1}},z^{2^{-j}}]\).
Their sums of local bad probabilities are bounded by a fixed
constant \(M_0\), uniformly in targets and directives.
Let \(U_j\) be the even Bonferroni polynomial through degree
\(2(s+j)\), and \(L_j\) the odd one through degree \(2(s+j)+1\).
If \(I_j\) is the indicator of no bad event in block \(j\), then
\(U_j\ge I_j\ge L_j\), \(U_j\ge0\), and telescoping gives
\begin{equation}\label{eq:older-lowerpoly}
 \prod_jU_j-\sum_j(U_j-L_j)\prod_{i\ne j}U_i
                                      \le\prod_jI_j .
\end{equation}
This avoids multiplying lower polynomials which might be negative.
In the independent prime model,
\[
 \E(U_j-L_j)\le\frac{M_0^{2(s+j)+1}}{(2(s+j)+1)!},
 \qquad \E U_j\ge\prod_{p\text{ in block}}(1-g_e(p))
                                             \ge e^{-2M_0}.
\]
Choose the fixed \(s\) so that the sum of these relative errors
is at most \(1/4\).
Factoring \(\prod_j\E U_j\) shows that
\eqref{eq:older-lowerpoly} has mean at least half the fully rough
probability.  Its squarefree divisor support is at most
\[
 C_0z^{\sum_{j\ge0}(2(s+j)+1)2^{-j}}
                               =C_0z^{4s+6},
\]
with coefficients \(O(\log\log N)\).
Choose \(\gamma<\xi/(4s+7)\).
Then its support is at most \(N^\xi\), so
Proposition~\ref{prop:older-weighted} applies term by term.
There are at most \(N^{\alpha+o(1)}\) cofactor triples.
The total error, including sieve predicates, is
\[
                  O(N^{1-\eta+\alpha+\xi+o(1)}),
\]
negligible by \eqref{eq:older-smallbudgets}.
Thus the summed actual weighted count is at least a fixed fraction
of the sum of \eqref{eq:older-triplemass}, apart from that error.

Temporarily drop the size cap \(duv\le N^\alpha\).
The positive cofactor sum is an Euler product.
At a split prime not dividing \(N\), its options are
\(1,a_p,a_p\).
At a split prime with target valuation \(e\ge1\), the options
omitting or including it in \(d\) give
\[
 E_p+p^{-1}E_{p,e-1}\ge E_p+1/p=1+2a_p.
\]
At inert target primes no restoration option is present; the
factor left over is \(E_p\).
Target factors above \(z\) contribute \(1+o(1)\), uniformly.
Consequently \eqref{eq:older-roughdimension} and prime Mertens give
\begin{equation}\label{eq:older-restoredmass}
 \sum_{d,u,v}\mathcal M_{d,u,v}
 \gg NJ(N)(\log N)^{-2}
          \prod_{\substack{p\le w\\\chi(p)=1}}(1+2a_p)
 \asymp \frac{NJ(N)}{\log N}.
\end{equation}

The cap must be restored in the actual positive model.
After normalizing its Euler product, each prime is included in
\(d,u\), or \(v\) with probability \(O(1/p)\), uniformly in target
valuation.  Therefore
\(\E\log(duv)\le C_1\log w\).
Choose \(\omega<\min(\gamma,\alpha/(4C_1))\).
Markov retains at least three quarters of the mass in
\(duv\le N^\alpha\).  Applied to this model, rather than merely a
comparison product, it proves that the finite sieve retains the
order of magnitude in \eqref{eq:older-restoredmass}.

Finally lift each residual pair to \((dx,dy)\).
At every cofactor prime the numerical valuation patterns are,
respectively,
\((1,1)\) for \(d\), \((1,0)\) for \(u\), and \((0,1)\) for \(v\);
all other small-prime valuations are zero.
Thus the ordered numerical pair recovers the entire triple and
residual pair uniquely.
Writing \(x=ux'\), \(y=vy'\), exact valuation one gives
\[
       \frac{V(x)V(y)}{2^{\omega(u)}2^{\omega(v)}}
                                      =V(x')V(y').
\]
All prime factors of \(x',y'\) exceed \(N^\gamma\), so
\(V(x')V(y')\le4^{1/\gamma}\).
No factor \(V(d)^2\) is inserted: \(d\) lifts values, not their
divisor weights.
Injectivity and this uniform bound turn
\eqref{eq:older-restoredmass} into the required unweighted lower bound.

For the upper bound, sieve the integer parameter of \(Ax+By=N\)
using only good inert primes.  Both coefficients are units there.
There are two forbidden roots, merging exactly when \(p\mid N\).
A dimension-one upper Selberg sieve at fixed power level therefore
gives
\[
 \ll N\prod_{\substack{p\le N^{\gamma_1}\\
                       p\nmid L_0,\ \chi(p)=-1}}
       \left(1-\frac{2-\1_{p\mid N}}p\right)
            +O(N^{1/2}\log^C N)
 \ll \frac{NJ(N)}{\log N}.
\]
For an explicit finite version, use the nonnegative even block
polynomials above at divisor support at most \(N^{1/2}\).
Their model is at most a fixed multiple of the true product,
by the same summable factorial block errors.
CRT gives \(N\nu(\ell)/\ell+O(\nu(\ell))\), \(\nu(p)\le2\);
the total endpoint error is
\(O(N^{1/2}\log^C N)\).
Discarding the residue filters only increases this upper bound.
This completes the proof of Lemma~\ref{lem:older-quadratic-filter}.

\paragraph{AI assistance and review.}
ChatGPT (OpenAI) assisted with mathematical arguments, literature checks,
exposition, proof review and verification code.
\begin{thebibliography}{99}\small

\bibitem{bourgainfuchs}
J. Bourgain and E. Fuchs,
\emph{On representation of integers by binary quadratic forms},
arXiv:1105.3698v2, 23 May 2011.
\url{https://arxiv.org/abs/1105.3698v2}.

\bibitem{richert}
H.-E. Richert,
\emph{Lectures on Sieve Methods}, notes by S. Srinivasan,
Tata Institute of Fundamental Research, Bombay, 1976,
Theorems 11.5--11.8, pp.~138--142.
\href{https://www.math.utoledo.edu/~codenth/Spring_13/3200/NT-books/Lectures_on_Sieve_Methods-Richert.pdf}{Online edition}.

\bibitem{iwaniecpomykala}
H. Iwaniec and J. Pomyka{\l}a,
\emph{Sums and differences of quartic norms},
Mathematika \textbf{40} (1993), 233--245.
\url{https://doi.org/10.1112/S0025579300007014}.

\bibitem{khayutin}
I. Khayutin,
\emph{Joint equidistribution of CM points},
Annals of Mathematics (2) \textbf{189} (2019), 145--276.
\url{https://doi.org/10.4007/annals.2019.189.1.4}.

\bibitem{loughranmatthiesen}
D. Loughran and L. Matthiesen,
\emph{Frobenian multiplicative functions and rational points in fibrations},
Journal of the European Mathematical Society \textbf{26} (2024),
4779--4830.
\url{https://doi.org/10.4171/JEMS/1374}.

\bibitem{milnecft}
J. S. Milne,
\emph{Class Field Theory}, version 4.03, 6 August 2020.
\url{https://www.jmilne.org/math/CourseNotes/CFT.pdf}.

\bibitem{weightedgraphs}
OpenAI,
\emph{Weighted dilation graphs, smooth shifted primes and totient fibers},
24 September 2026, Section 3.
\href{https://github.com/openai/math/blob/4eaee87548250b2682affda7d05d947c05554bae/preprints/Weighted-Dilation-Graphs-Smooth-Shifted-Primes-and-Totient-Fibers-September-24-2026/paper.pdf}{Immutable primary PDF}.

\bibitem{zerofree}
OpenAI,
\emph{The Quasi-Riemann Hypothesis:
A Zero-Free Half-Plane \(\Re(s)>7/8\)},
30 September 2026, Part I, Theorem 3.1.
\href{https://github.com/openai/math/blob/4eaee87548250b2682affda7d05d947c05554bae/preprints/The-Quasi-Riemann-Hypothesis-September-30-2026/paper.pdf}{Immutable primary PDF}.

\bibitem{walsh}
M. N. Walsh,
\emph{Shifted convolutions of divisor functions},
arXiv:2610.40439v2, 12 November 2026.
\url{https://arxiv.org/abs/2610.40439v2}.

\bibitem{smilansky}
Y. Smilansky,
\emph{Sums of two squares---pair correlation and distribution in short intervals},
International Journal of Number Theory \textbf{9} (2013), 1687--1711.
\url{https://doi.org/10.1142/S1793042113500516}.
\url{https://arxiv.org/abs/1208.2419v2}.

\bibitem{heathbrown}
D. R. Heath-Brown,
\emph{The divisor function \(d_3(n)\) in arithmetic progressions},
Acta Arithmetica \textbf{47} (1986), 29--56.
\url{https://doi.org/10.4064/aa-47-1-29-56}.

\bibitem{delabretechetenenbaum}
R. de la Bret\`eche and G. Tenenbaum,
\emph{Remarks on the Selberg--Delange method},
Acta Arithmetica \textbf{200} (2021), 349--369.
\url{https://doi.org/10.4064/aa201024-26-5}.
\href{https://tenenb.perso.math.cnrs.fr/PPP/On-SD.pdf}{Corrected author edition}.

\bibitem{soundararajan}
K. Soundararajan,
\emph{The Liouville function in short intervals
[after Matom\"aki and Radziwi\l\l]},
arXiv:1606.08021v1, 26 June 2016, Lemmas 2.1--2.2.
\url{https://arxiv.org/abs/1606.08021v1}.

\bibitem{taonotes}
T. Tao,
\emph{254A, Notes 2: Complex-analytic multiplicative number theory},
9 December 2014, Proposition 51, Theorem 62(ii),
Exercise 45(i)--(ii).
\url{https://terrytao.wordpress.com/2014/12/09/254a-notes-2-complex-analytic-multiplicative-number-theory/}.

\bibitem{exactcompanion}
R. van Rijn,
\emph{Exact Sumsets of Primitive Quadratic Norms},
version 1, Zenodo, 2026.
\url{https://doi.org/10.5281/zenodo.24371462}.

\end{thebibliography}

\end{document}
